\documentclass[11pt,letterpaper]{article}
\usepackage[T1]{fontenc}
\usepackage[utf8]{inputenc}
\usepackage{lmodern}
\usepackage[margin=1in,headheight=15pt]{geometry}
\usepackage{amsmath,amssymb,amsthm,mathtools}
\usepackage{microtype,booktabs,longtable,array,enumitem}
\usepackage[dvipsnames]{xcolor}
\usepackage{fancyhdr,needspace}
\usepackage{listings}
\usepackage{xurl}
\usepackage[colorlinks=true,linkcolor=MidnightBlue,citecolor=MidnightBlue,urlcolor=MidnightBlue,pdfencoding=auto]{hyperref}
\usepackage{bookmark}
\hypersetup{pdftitle={Stochastic and Thermal Exactness},pdfauthor={AI-assisted research report prepared for Vladimir Reshetnikov},pdfsubject={Diophantine certificates for exact events in uniformly convergent systems: positive stochastic erasure, exact value iteration of stochastic games, and thermal Hamiltonians with undecidable confinement boundaries}}
\urlstyle{same}
\pagestyle{fancy}
\fancyhf{}
\fancyhead[L]{\small Stochastic and Thermal Exactness}
\fancyhead[R]{\small Research report}
\fancyfoot[C]{\thepage}
\setlength{\parskip}{0.32em}
\setlength{\parindent}{1.1em}
\setlength{\emergencystretch}{3em}
\setlist{itemsep=0.25em,topsep=0.4em}
\allowdisplaybreaks[2]
\newtheorem{theorem}{Theorem}[section]
\newtheorem{lemma}[theorem]{Lemma}
\newtheorem{proposition}[theorem]{Proposition}
\newtheorem{corollary}[theorem]{Corollary}
\theoremstyle{definition}
\newtheorem{definition}[theorem]{Definition}
\newtheorem{example}[theorem]{Example}
\theoremstyle{remark}
\newtheorem{remark}[theorem]{Remark}
\newcolumntype{L}[1]{>{\raggedright\arraybackslash}p{#1}}
% Shared by the three sources
\newcommand{\N}{\mathbb N}
\newcommand{\Z}{\mathbb Z}
\newcommand{\Q}{\mathbb Q}
\newcommand{\R}{\mathbb R}
\newcommand{\C}{\mathbb C}
\newcommand{\one}{\mathbf 1}
\DeclareMathOperator{\diag}{diag}
% Source 02 (Part I)
\newcommand{\rk}{\operatorname{rank}}
\newcommand{\im}{\operatorname{im}}
\newcommand{\TV}{\operatorname{TV}}
\newcommand{\Mat}{\operatorname{Mat}}
\newcommand{\EE}{\mathbb E}
\newcommand{\PP}{\mathbb P}
\newcommand{\code}{\operatorname{val}}
\newcommand{\transpose}{^{\mathsf T}}
\newcommand{\repo}{\texttt{928ea97017a25ebe56d240c84f27d2275d818c75}}
% Source 10 (Part II)
\newcommand{\zero}{\mathbf 0}
\newcommand{\opt}{\operatorname{opt}}
\newcommand{\spn}{\operatorname{span}}
\newcommand{\enc}{\operatorname{enc}}
\newcommand{\bits}{\operatorname{bits}}
\newcommand{\lcm}{\operatorname{lcm}}
% Source 03 (Part III)
\newcommand{\Hil}{\mathcal H}
\newcommand{\ket}[1]{\lvert #1\rangle}
\newcommand{\abs}[1]{\left\lvert #1\right\rvert}
\newcommand{\norm}[1]{\left\lVert #1\right\rVert}
\DeclareMathOperator{\Tr}{Tr}
\DeclareMathOperator{\Spec}{spec}
\DeclareMathOperator{\ess}{ess}
\DeclareMathOperator{\rank}{rank}
\DeclareMathOperator{\supp}{supp}
\newcommand{\cB}{\mathcal C_B}
\newcommand{\epsm}{\varepsilon_m}
% Merge
\newcommand{\wtag}{\textup{[write]}}
\lstset{basicstyle=\ttfamily\small,breaklines=true,columns=fullflexible,keepspaces=true,showstringspaces=false,frame=single,rulecolor=\color{black!20}}
\setcounter{tocdepth}{2}

\begin{document}
\hypersetup{pageanchor=false}
\begin{titlepage}
\vspace*{0.45in}
\noindent{\color{MidnightBlue}\large PROVEIT RESEARCH REPORT}\par
\vspace{0.35in}
\noindent{\Huge\bfseries Stochastic and Thermal\\[0.12em] Exactness}\par
\vspace{0.25in}
\noindent{\LARGE Diophantine certificates for exact events\\[0.15em] in uniformly convergent systems}\par
\vspace{0.3in}
\noindent{\large Positive stochastic erasure, exact value iteration of stochastic games, and thermal Hamiltonians with undecidable confinement boundaries}\par
\vspace{0.5in}
\noindent AI-assisted research report prepared for Vladimir Reshetnikov from three manuscripts\par
\noindent(batch~78 of ProveIt's incoming reports, manuscripts~02, 10 and~03), October 2, 2026\par
\vspace{0.3in}
\begin{abstract}
Three independent manuscripts share one pattern: a numerical approximation converges uniformly, at a rate known in advance, while an \emph{exact} event of the same system is undecidable.

Part~I (source~02) compiles integer matrix mortality into strictly positive rational column-stochastic matrices with no common invariant complex line, each entry within any prescribed band around $1/n$ and every product contracting total variation by $\eta^T$. A product has rank one exactly when the source projection is mortal, so exact erasure is undecidable for at most eight positive $7\times7$ matrices and many-one complete for a nine-state family with a variable alphabet. Explicit horizon-$T$ natural quartics have one witness per erasing labelled word.

Part~II (source~10) builds a fixed turn-based stochastic game with discount $1/2$, reward $1/4$, dyadic transition rows of full support and the known value vector $\tfrac12\one$. Value iteration converges at rate $2^{-n}$, yet exact attainment of $\tfrac12\one$ from a dyadic terminal payoff is computably enumerable complete. Each horizon has an orthant-nonnegative quadratic certificate with exactly one complete natural witness.

Part~III (source~03) turns a natural arithmetic circuit into a diagonal Hamiltonian $H=(1+\sum_in_i)Q$ of degree at most five whose zero states biject with the witnesses and whose positive spectrum is effectively confined. Compact resolvent and finite thermal trace are equivalent to finite witness multiplicity and are undecidable on a padded universal family; a uniformly computable critical fugacity $t_*$ satisfies $t_*<1$ exactly on halting inputs.

The report is AI-assisted and unrefereed; nothing in it is formalized in Lean or Rocq, and priority is not certified for any Part.
\end{abstract}
\vfill
\noindent\small\textbf{Status.} Proofs and executable finite checks; not kernel-checked. Section~\ref{ste:sec:front} states the common spine, the notation, the relation to neighbouring reports and to the repository's reviews of these manuscripts, and what is not claimed; Appendix~\ref{ste:app:provenance} gives the provenance.
\end{titlepage}
\hypersetup{pageanchor=true}

\tableofcontents
\newpage

\section{Scope, common spine and conventions}\label{ste:sec:front}

\wtag\ This section was written when the three manuscripts were merged. Parts~\ref{ste:part:ee}, \ref{ste:part:ec} and~\ref{ste:part:ta} print the three sources; text there without a \wtag\ marker is the source's own.

\subsection{The three Parts}\label{ste:sec:parts}

This report is built from three AI-assisted research manuscripts delivered with batch~78 of ProveIt's incoming reports (manuscripts~02, 10 and~03; Appendix~\ref{ste:app:provenance} gives the full provenance). They are called \emph{source~02}, \emph{source~10} and \emph{source~03} after their batch-78 manuscript numbers, which are also the file prefixes of their shipped programs and data. None cites another, they were written from three different snapshots of the repository on the same day, and they share no lemma.

\begin{itemize}
\item \textbf{Part~\ref{ste:part:ee} (source~02), \emph{Exact erasure without a common invariant line}.} A finite alphabet of rational column-stochastic matrices is driven by an external schedule; a word \emph{erases} when its product has rank one, so that the output distribution no longer depends on the input distribution. Two invertible tensor guards lift integer mortality into strictly positive stochastic matrices with no common invariant complex line and the exact rank identity $\rk S_w=1+2\rk A_{\pi(w)}$ (Theorem~\ref{ste:ee:thm:linefree}). Rank-one existence is undecidable for at most eight positive $7\times7$ matrices (Corollary~\ref{ste:ee:cor:seven}) and many-one complete for nine states with a variable alphabet (Corollary~\ref{ste:ee:cor:complete}). Under independent full-support switching, existence of an erasing word, almost-sure finite erasure and finite expected erasure time coincide (Theorem~\ref{ste:ee:thm:iid}). Horizon-$T$ natural quartics have exactly one witness per erasing labelled word (Theorems~\ref{ste:ee:thm:fullcertificate} and~\ref{ste:ee:thm:compressed}).
\item \textbf{Part~\ref{ste:part:ec} (source~10), \emph{Exact convergence as computation}.} A counter program is compiled through homogeneous dyadic min/max circuits into a genuine two-player turn-based stochastic game with discount $1/2$, running reward $1/4$, at most two actions per state and every transition probability at least $1/(2N)$. The Bellman operator is a $1/2$-contraction with the known fixed point $\tfrac12\one$, and exact finite attainment of that fixed point from a dyadic terminal payoff is computably enumerable complete for a suitable fixed game (Theorem~\ref{ste:ec:thm:main}). Point-mass selectors and chance-only games give decidable boundaries, and chance-only scalar equality is Skolem-equivalent (Section~\ref{ste:ec:sec:boundaries}). Each horizon has an orthant-nonnegative quadratic with exactly one complete natural witness (Theorem~\ref{ste:ec:thm:quadratic}).
\item \textbf{Part~\ref{ste:part:ta} (source~03), \emph{Thermal arithmetic}.} A natural arithmetic circuit becomes a diagonal Hamiltonian $H=(1+\sum_in_i)Q$ on occupation modes, of joint degree at most five, whose zero-energy basis states are in bijection with the witnesses while every positive energy is at least $1+\sum_in_i$ (Theorem~\ref{ste:ta:thm:compiler}). Compact resolvent and finite thermal trace hold exactly when the witnesses are finitely many (Theorem~\ref{ste:ta:thm:classification}), and both are undecidable on a padded universal family (Theorem~\ref{ste:ta:thm:padding}). A uniformly computable critical fugacity satisfies $t_*(c)<1$ exactly on halting inputs (Theorem~\ref{ste:ta:thm:critical}).
\end{itemize}

Each Part keeps its own definitions, letters and numbering of statements. Parts~\ref{ste:part:ee} and~\ref{ste:part:ec} are both about stochastic operators with full support (stochastic games extend stochastic matrices by minima and maxima); Part~\ref{ste:part:ta} shares no lemma with either, and it is placed here because of the spine of Section~\ref{ste:sec:spine}, not because of a shared construction.

\subsection{The common spine: uniform approximation, undecidable exactness}\label{ste:sec:spine}

In each Part an approximation statement holds with a computable rate that does not depend on the instance, and an exact statement about the same object is undecidable. Table~\ref{ste:tab:spine} lines them up.

\begin{table}[ht]
\centering\small
\begin{tabular}{@{}L{0.09\textwidth}L{0.27\textwidth}L{0.27\textwidth}L{0.29\textwidth}@{}}
\toprule
Part & Uniform approximation & Exact predicate and its status & Bounded or static certificate\\
\midrule
I (02) & $\TV(S_wp,S_wq)\le\eta^T\TV(p,q)$ for every word of length $T$ (Proposition~\ref{ste:ee:prop:contraction}) & some product has rank one: undecidable on the seven-state class, c.e.-complete on the nine-state class (Corollaries~\ref{ste:ee:cor:seven}, \ref{ste:ee:cor:complete}) & horizon-$T$ quartic, $T((n-1)^2+k)$ natural witnesses, one per erasing labelled word (Theorem~\ref{ste:ee:thm:compressed})\\
\addlinespace
II (10) & $\norm{\Phi^n(v)-\tfrac12\one}_\infty\le2^{-n}\norm{v-\tfrac12\one}_\infty$, and $4^{-n}/2$ on encoded inputs (Section~\ref{ste:ec:sec:mixing}) & $\exists n\;\Phi^n(v)=\tfrac12\one$: c.e.-complete for a fixed game (Theorem~\ref{ste:ec:thm:main}) & horizon-$T$ quadratic, $(N+2b)T$ natural witnesses, one complete witness (Theorem~\ref{ste:ec:thm:quadratic})\\
\addlinespace
III (03) & excited trace, regulated partition function and $t_*$ uniformly computable (Theorems~\ref{ste:ta:thm:remainders}, \ref{ste:ta:thm:partition}, \ref{ste:ta:thm:critical}) & $t_*(c)<1$: c.e.-complete; compact resolvent, finite thermal trace: co-c.e.-complete on a padded family (Theorems~\ref{ste:ta:thm:critical}, \ref{ste:ta:thm:padding}) & static polynomial of degree at most five with a fixed number of modes, one zero state per witness (Theorem~\ref{ste:ta:thm:compiler})\\
\bottomrule
\end{tabular}
\caption{\wtag\ The common spine of the three Parts.}\label{ste:tab:spine}
\end{table}

Three further features recur. First, every Part proves that no total computable function of the input bounds the time or the witness height of a positive instance (Theorem~\ref{ste:ee:thm:nobound}, Corollary~\ref{ste:ec:cor:no-cutoff}, Theorem~\ref{ste:ta:thm:gap}); this is the bounded-search argument of \texttt{cdc:bd:prop:nobound} in~\cite{cdc} and \texttt{pqc:qm:cor:nobound} in~\cite{pqc}, applied to new predicates. Second, a fixed positive tolerance restores decidability or a finite cutoff in Parts~I and~II ($\eta^T\le\varepsilon$; Proposition~\ref{ste:ec:prop:threshold}), while in Part~III no computable positive separation from the boundary exists (Corollary~\ref{ste:ta:cor:separation}). Third, each Part also states an MRDP consequence of fixed arity and degree at most four obtained by circuit quadratization (Proposition~\ref{ste:ee:prop:fixedquartic}; Section~\ref{ste:ec:sec:mrdp}; Corollary~\ref{ste:ta:cor:universal} with Theorem~\ref{ste:ta:thm:compiler}); none of them controls the multiplicity of the original MRDP witnesses.

The slogan itself is not new in the collection: Parts~V and~VII of~\cite{pqc} already prove contraction-versus-exactness results, and sources~02 and~10 say so themselves (Sections~\ref{ste:ee:sec:intro} and~\ref{ste:ec:sec:intro}). What each Part adds is stated in Section~\ref{ste:sec:status}.

\subsection{Notation across the Parts}\label{ste:sec:notation}

No symbol was renamed and no normalization was changed. Each Part keeps its source's letters, and a letter means what its own Part says it means. Table~\ref{ste:tab:notation} lists the letters that the Parts use differently.

\begin{longtable}{@{}L{0.07\textwidth}L{0.28\textwidth}L{0.28\textwidth}L{0.28\textwidth}@{}}
\caption{\wtag\ Letters used differently in the three Parts.}\label{ste:tab:notation}\\
\toprule
Letter & Part I (source 02) & Part II (source 10) & Part III (source 03)\\
\midrule
\endfirsthead
\toprule
Letter & Part I (source 02) & Part II (source 10) & Part III (source 03)\\
\midrule
\endhead
\bottomrule
\endfoot
$n$, $N$ & $n$: state dimension; $N$: nilpotent shift, also $N_T$ (erasing words of length $T$) and the input bit length in $b(N)$ & $N$: number of game states; $n$: iteration count & $N_i$: number operators; $N$: index of the example $H_N$; $n\in\N^m$: occupation tuple\\
$K$ & common lift scale, $K\ge2$ & layer normalizer $K=K_L$, a power of two & $\mathsf K$: halting set; $K$: compact remainder operator\\
$D$ & common column sum $D=nK$ & common probability denominator & $D=(B+1)^m$: tuples in a truncation box\\
$M$ & largest entry of $\abs{\Delta_i}$ & pipeline period $M=L+1$ & number of zero states $\#\mathcal Z$\\
$T$ & word length or horizon; $T_i=p_iS_i$ labelled transitions & horizon & $T_q$: trace-class remainder\\
$H$ & $H_n$: zero-mass hyperplane; $H_i$: integer restriction matrices & first halting time & the Hamiltonian\\
$Q$ & PCP connector $Q=c\ell$; $\mathcal Q_T$: full quartic & $Q_T$: complete quadratic & compiled quartic penalty\\
$S$, $s$ & $S_i$: stochastic matrices; $s_t$: selector residuals & $S$: game operator; $s$: clock register & $s(n)=\sum_in_i$: total occupation\\
$E$ & coordinate map $E$ & $E_T$: endpoint sum & energy level\\
$L$ & guard $L$; length of an erasing word & final circuit layer & minimum zero-state height\\
$R$ & guard $R$ & power of two in the Skolem lift & excited sum $R(q,t)$\\
$P$ & cyclic permutation (Section~\ref{ste:ee:sec:higher}) & integer transition numerators $P_{ia,j}$; chance-only matrix & source polynomial; $P_0$: zero projection\\
$\eta$ & contraction band & reset weight & requested width in the proof of Theorem~\ref{ste:ta:thm:critical}\\
$m$, $k$ & $m$ source matrices, $k$ letters & $m$ control states, $k$ counters & $m$ modes, $k$ source witnesses\\
$q$, $t$ & $q$: probability of a block; $t$: time & $t$: time step & $q$: Boltzmann parameter; $t$: fugacity\\
\end{longtable}

Watch for these readings:
\begin{itemize}
\item \textbf{Erasure.} In Part~I a word \emph{erases} when its stochastic product has rank one: every input distribution is sent to one common output distribution, which need not be uniform (a \emph{uniform} erasing word gives $J_n$). In Part~II the \emph{erasing map} $F_{\rm era}$ and the \emph{erasing halt convention} send a halted counter configuration, and then the whole simulated pipeline vector, to the zero vector; the exact event observed there is attainment of the fixed point $\tfrac12\one$, not rank one of a product. The two notions are unrelated constructions that happen to share a word.
\item \textbf{Uniform.} Part~I: a uniform erasing word (target $J_n$), uniform switching (equal letter probabilities), and the uniform-input quartic (matrix entries as parameters). Part~II: uniformly mixing (every row of full support) and the common uniform reset. Part~III: uniformly computable (one algorithm for every instance).
\item \textbf{Contraction and confinement.} Part~I contracts total variation by $\eta$ per letter; Part~II contracts the supremum norm by $\delta=1/2$ and the span by $\delta(1-\eta)$. Part~III has no dynamics at all: its \emph{confinement} is the spectral inequality $H>0\Rightarrow H\ge1+s$, and its diagonal evolution only changes phases (Section~\ref{ste:ta:sec:gap}).
\item \textbf{Certificate.} The certificates of Parts~I and~II are horizon-indexed: their arity grows with $T$. The Hamiltonian of Part~III has a fixed number of modes for a fixed circuit. None of the three is a fixed-arity single-fold representation of a computably enumerable complete set.
\item \textbf{Fugacity, thermal.} In Part~III these are mathematical parameters of a positive trace; no physical reservoir, laboratory implementation or physical decision procedure is assumed (source~03's own caveat).
\end{itemize}

\subsection{Relation to neighbouring reports and to the repository's reviews}\label{ste:sec:relation}

This report sits in the collection's \texttt{hilbert-tenth-problem} category beside four earlier reports and the signal-machine report written in the same batch. Its closest neighbour is the report on probabilistic, quantum and continuous computation~\cite{pqc}, cited here by label names (it is a separate document):
\begin{itemize}
\item \emph{Part~II of~\cite{pqc} (quantum mortality).} Theorem \texttt{pqc:qm:thm:47} imports the undecidability of mortality for six integer $3\times3$ matrices and packs it into seven rational $4\times4$ Kraus operators, so that a zero-probability outcome word exists exactly when the source is mortal. Corollary~\ref{ste:ee:cor:seven} runs the same pipeline with a different exact predicate (rank one of a positive stochastic product) and Theorem~\ref{ste:ee:thm:nobound} is the argument of \texttt{pqc:qm:cor:nobound}. Proposition \texttt{pqc:qm:prop:nonnegative} decides mortality for nonnegative matrices; Part~I shows that this positivity boundary does not transfer to rank one (note after Corollary~\ref{ste:ee:cor:seven}), which bears on Question \texttt{pqc:q:signed}.
\item \emph{Part~V of~\cite{pqc}.} Source~13's witness-faithful quartic energy (\texttt{pqc:ql:thm:compiler}) is a bounded Boolean energy landscape over $\R$ and expressly treats no Gibbs distribution; Part~III is the unbounded natural counterpart with a thermal trace. Source~17's reversible chains have a uniform spectral gap and an undecidable exact stationary equality, the same slogan as Part~II for a different object.
\item \emph{Part~VII of~\cite{pqc}.} Source~12 proves r.e.-complete exact extinction for a 2-adic $1/2$-contraction (\texttt{pqc:ad:cor:complete}); source~14 a real piecewise-linear $1/2$-contraction whose point-to-point reachability is decidable (\texttt{pqc:sc:thm:main}~(iii)); source~15 the unsquared complementarity device and the degree-two unique-witness trace polynomial (\texttt{pqc:pm:lem:clamp}, \texttt{pqc:pm:thm:trace}). Part~II is complementary to all three: a real, order-preserving $1/2$-contraction whose exact fixed-point attainment is c.e.-complete because its clock is erased at halting (note after Section~\ref{ste:ec:sec:intro}'s comparison), and its quadratic specializes source~15's device to min/max gates. Part~II also answers Part~VII's unlabelled question ``Structural subclasses'' negatively in part (note after the table of Section~\ref{ste:ec:sec:boundaries}).
\item \emph{Parts~VIII and~IX of~\cite{pqc}.} Corollary \texttt{pqc:tu:cor:fixedquartic} and its remark (a fixed MRDP quartic whose witness multiplicity is uncontrolled) are the observation that Part~III's spectral classification makes exact. Theorem \texttt{pqc:iq:thm:geometric} (geometric stopping clocks with undecidable records) is related only to Part~I's zero--one law.
\end{itemize}

The report on canonical Diophantine certificates~\cite{cdc} supplies four results that the Parts reprove, or use without citation, and that are printed here as pointers or marked second routes, never as new: the unique-extension quadratization \texttt{cdc:wf:lem:quadratization} (and \texttt{cdc:cs:lem:quartic}), which is the core of Theorem~\ref{ste:ta:thm:compiler} and of the fixed-arity quartic corollaries of Parts~I and~II; the absence of a computable witness-height bound \texttt{cdc:bd:prop:nobound}; the finite-fold substrate equivalence \texttt{cdc:cs:thm:finitefold}, the template of Theorem~\ref{ste:ta:thm:finitefold}; and the positivity-restricted degree classification \texttt{cdc:of:thm:classification}, which shows that radial isolation below degree four is impossible for a universal family (note in Section~\ref{ste:ta:sec:questions}). The reports \emph{liveness-beyond-halting}, \emph{group-theoretic-substrates} and \emph{signal-machine-collision-certificates} have no overlap with this one.

The Hilbert's-tenth-problem research programme (read-only for this report) had already reviewed all three manuscripts before they were placed here; this report does not present its placement as a first review:
\begin{itemize}
\item \cite{review1977} reviewed sources~02 and~03 (commit \texttt{44b28c395}). It accepts source~02 within its stated domains, replays its suite and all eight certificates and adds independent checks; it reports no defect. It accepts source~03's mathematics for natural arithmetic DAGs and finds that the delivered low-level \texttt{Compiler}/\texttt{Atom} interface accepts invalid descriptors; it supplies the tested patch \path{thermal_exact_dag_guards.patch}. The shipped compiler here is the unpatched original (note in Section~\ref{ste:ta:sec:package}).
\item \cite{review808} reviewed source~10 (commit \texttt{126028588}). It accepts the mathematics and finds two implementation defects: the independent checker does not bind its scaled initial vector to the declared input, and an explicitly supplied integer scale in the optional counter loader is divided in floating point; it supplies the patches \path{exact_convergence_input_binding.patch} and \path{exact_convergence_integer_scale.patch}. The shipped programs here are the unpatched originals (notes in Section~\ref{ste:ec:sec:reproduce} and Appendix~\ref{ste:ec:app:contract}).
\item The note \path{matrix_pcp_trace.md}~\cite{matrixpcp} uses the same classical radix-three encoding of Post correspondence tiles as Lemma~\ref{ste:ee:lem:pcpblock}, with a leading sentinel instead of Part~I's fourth coordinate and idempotent connector.
\end{itemize}

\emph{The formal project.} The report sits in the collection, not in \path{Computability/HilbertTenthProblem}, and placement beside a Lean/Rocq development confers no formal status. The three sources import MRDP only as a classical theorem; source~10 names the project's interface, the theorem \texttt{Diophantine.mrdp} with its variants \texttt{Diophantine.mrdp\_iff} and \texttt{Diophantine.mrdp\_dioph\_iff} in \path{Computability/HilbertTenthProblem/Lean/Diophantine/MRDP.lean} (unchanged from source~10's pin to this write). The project has formalized none of this report's statements: it has no stochastic-matrix, stochastic-game, Bellman-operator or occupation-number Hamiltonian module.

\subsection{Status and what is not claimed}\label{ste:sec:status}

The report claims conventional mathematical proofs, by its sources, of the results summarized in Section~\ref{ste:sec:parts}. What each source presents as its own contribution is narrower: source~02 the two-guard no-common-line compiler, its higher-dimensional subspace exclusion, the exact-forgetting interpretation and the counted certificate interface; source~10 the stochastic-game realization under the simultaneous restrictions of Theorem~\ref{ste:ec:thm:main} with its complete arithmetic interface; source~03 the radial factor $1+\sum_in_i$ with its effective spectral decomposition, the multiplicity classification and the critical-fugacity construction.

Every limitation stated by a source is kept in its Part. Collected here, the report does \textbf{not} claim:
\begin{itemize}
\item historical priority for any Part (all three sources say their searches were targeted), or any refereeing, proof-assistant verification, or rebuild of the repository's Lean or Rocq developments; the finite checks (50,582, 26,077 and 7,936 assertions) illustrate and do not prove;
\item an improvement of the repository's 75-operation complete-certificate or 87-operation universal-polynomial records; no Part instantiates a numerical universal polynomial, alphabet or game;
\item a finite-fold or single-fold Diophantine representation of computably enumerable sets; the bounded certificates of Parts~I and~II have horizon-dependent arity, and the fixed-arity MRDP quartics leave the original witness multiplicity uncontrolled;
\item \emph{Part~I:} irreducibility (a common invariant hyperplane remains); one fixed universal stochastic alphabet with an ordinary input loader (the seven-state result is a family of alphabets parameterized by the input); many-one completeness of the eight-letter bound (only of the nine-state variable-alphabet family); anything about undecidable output entropy of hidden Markov models; optimality of the dimension, alphabet or certificate counts. Its headline answer to a MathOverflow question rests on an external discussion that is not archived in the repository;
\item \emph{Part~II:} a numerical universal game (the shipped 128-state games are one-counter examples); hardness of ordinary value iteration from zero or of solving finite discounted games; a resolution of the Skolem problem or of the one-player Bellman reachability frontier; robustness to perturbed tables or floating point; novelty of the chance-only Skolem specialization (Vahanwala);
\item \emph{Part~III:} a resolution of the finite-fold problem; physical hypercomputation or autonomous step-by-step dynamics; a numerical value of the mode count of the universal family; that $t_*$ is a noncomputable real (it is computable; its exact endpoint is undecidable); geometric locality from four-mode support; robustness to coefficient errors.
\end{itemize}

\part{Exact erasure without a common invariant line (source 02)}\label{ste:part:ee}

\noindent\wtag\ Part~I prints batch-78 manuscript~02, \emph{Exact Erasure without a Common Invariant Line: Positive stochastic computation, tensor guards, and Diophantine certificates}, ``Research report prepared for Vladimir Reshetnikov; Mathematical development and implementation: ChatGPT'', October 2, 2026. Its labels carry the prefix \texttt{ste:ee:}. Its abstract and status statement follow; its two appendices are Appendices~\ref{ste:ee:app:tinyquartic} and~\ref{ste:ee:app:reproduction}.

\begin{quote}\small
\textbf{Abstract (source 02).} We give an explicit reduction from integer matrix mortality to exact finite-time
forgetting by strictly positive rational stochastic matrices. Two tensor-factor
guards remove every common invariant complex line without creating spurious
zero products. From $m$ matrices of dimension $d$, the construction produces
$m+2$ matrices of dimension $2d+1$, with every entry arbitrarily close to the
uniform stochastic matrix and a prescribed uniform total-variation contraction
factor. A product has rank one precisely when its source-letter projection is
mortal. Consequently, rank-one product existence is undecidable already for at
most eight positive rational $7\times7$ column-stochastic matrices having no
common invariant complex line. Clearing one common denominator gives the same
conclusion for strictly positive integer matrices. This answers the specific
no-common-invariant-line restriction raised in a June 2026 MathOverflow discussion;
it does not establish irreducibility.

A separate, fully printed four-dimensional correspondence-problem compiler gives
a many-one complete nine-state stochastic family with a variable-size alphabet.
A higher tensor construction excludes all invariant subspaces of dimensions
$1,\ldots,r-1$. For independent full-support switching, existence of an erasing
word is equivalent both to almost-sure finite exact forgetting and to finite
expected forgetting time, despite deterministic exponential approximate
forgetting on every instance. Explicit mass-compressed horizon-$T$ natural-domain quartics have
$T((n-1)^2+k)$ witnesses and $T((n-1)^2+1)+(n-1)^2$ squared residuals, with
exactly one witness per accepted labelled word. We also prove a sharp $n-1$ erasure horizon
when the induced matrices on the zero-mass hyperplane commute, and exhibit
arbitrarily small non-erasing perturbations preserving all common invariant
subspaces. The package includes exact-arithmetic code, independent certificate
replay, and 50,582 passing finite checks.
\end{quote}
\begin{quote}\small
\textbf{Status (source 02).} Full mathematical proofs are supplied below. The new constructions
are not proof-assistant verified or independently refereed. Established mortality,
PCP, and DPRM results are cited as dependencies. The literature search is targeted,
not a certification of global priority. No improvement to ProveIt's fixed universal
polynomial operation count is claimed.
\end{quote}


\section{Research question, contributions, and provenance}
\label{ste:ee:sec:intro}

A stochastic matrix usually suggests information loss rather than universal
computation. If every entry is positive and close to $1/n$, every step loses a
large fraction of the distinction between two initial distributions. It is
therefore tempting to expect that determining whether all distinction is ever
lost exactly should be straightforward. The distinction between \emph{small}
and \emph{zero}, however, changes the problem completely.

This report studies a finite alphabet $S_1,\ldots,S_k$ of rational stochastic
matrices. A word is an externally selected schedule. The question is whether
some finite schedule makes the output distribution independent of the input
distribution. Equivalently, does a matrix product have rank one? The control
alphabet and rational coefficients are finite data; the word length is unbounded.

\subsection{A concrete question that the construction answers}

In a MathOverflow question posted on June 2, 2026, Victor Miller asked about the
rank-one analogue of matrix mortality. Benjamin Steinberg observed that
$A\mapsto A\oplus[1]$ transfers mortality to rank-one product existence. The
subsequent discussion isolated the restriction that the generators have no common
invariant line, and suggested an affine modification without a complete argument
for that restriction \cite{miller}. We prove the missing restricted statement,
with positivity and strong stochastic contraction added.
(\wtag\ The MathOverflow question and its June~3 discussion are external sources, not archived in the repository; that the construction answers their restriction is this source's claim, not checked at the write.)

The answer has a necessary qualification. Every column-stochastic matrix
preserves the hyperplane $\one\transpose x=0$. Our construction has no common
invariant \emph{line}, even over $\C$, but it is not an irreducible matrix family.
It answers the stated line restriction, not a stronger irreducibility requirement
that may arise in a particular application. It also does not settle any general
question about computing hidden-Markov entropy.

\subsection{Main results and their logical dependencies}

The principal construction, Theorem~\ref{ste:ee:thm:linefree}, is elementary once a matrix
mortality instance is supplied. Its central identity is
\begin{equation}
  \rk S_w=1+2\rk A_{\pi(w)},
  \label{ste:ee:eq:intro-rank}
\end{equation}
where $\pi$ deletes the two guard letters. The guards are invertible on their
private tensor factor, so inserting them anywhere cannot fabricate a mortal
source computation. One guard removes possible shared eigenvectors in a
transverse direction; an affine displacement in the other removes a common
stationary line.

For the small numerical undecidability bounds, we import only the established
mortality bounds of Cassaigne, Halava, Harju, and Nicolas \cite{cassaigne}.
Their conventions allow oracle Turing reductions, so we do not silently promote
every numerical bound in their paper to a many-one completeness statement.
Instead, Section~\ref{ste:ee:sec:pcp} prints a separate PCP reduction. That route gives
many-one completeness in dimension nine, with alphabet size part of the input.
The classical halting-to-PCP dependency is standard and has also been treated by
formal reduction proofs \cite{forster}.

The Diophantine contribution is explicit at a fixed horizon. Its witness
bijection is proved directly, rather than inferred from DPRM. The latter theorem
is used only for the separate fixed-arity, unbounded representation in
Section~\ref{ste:ee:sec:unbounded}. Its use does not preserve a unique witness per
unbounded computation.

\subsection{Relationship to the inspected ProveIt material}

The repository was inspected at the pinned commit
\begin{center}\small\repo.\end{center}
The relevant sources are its Hilbert-tenth-problem README, the October 2 substrate
review, and the README of the combined probabilistic/quantum/continuous report
\cite{repoH10,repoReview,repoProb}. These already contain extensive work on exact
observations, contractions, matrix substrates, and bounded certificates. We do
not present the generic claim ``computably enumerable implies Diophantine'' or
``contraction can coexist with undecidable exactness'' as new.

The proposed contribution is more specific: the two-guard no-invariant-line
compiler, its higher-dimensional invariant-subspace exclusion, its exact
stochastic forgetting interpretation, and the ensuing explicitly counted
certificate interface. The targeted inspection did not amount to rereading every
manuscript in the repository, and the public-source search is not an exhaustive
priority audit.

The October 2 review emphasizes the difference between an external-horizon
polynomial and one fixed polynomial accepting arbitrarily long computations
\cite{repoReview}. We retain that distinction throughout. The review's stated
87-operation fixed universal polynomial is not improved by any count in this
report. State dimension, residual count, arithmetic-circuit size, and number of
unbounded Diophantine witnesses are different resources.

\wtag\ After the pin, the research programme reviewed this manuscript, byte for byte the same package~\cite{review1977}: it accepts the results within their stated domains and quantifiers, replays the suite and all eight certificates, adds independent tensor-rank, natural selector/prefix and signed-selector checks, and reports no defect. The collection's report~\cite{pqc} has a parallel pipeline in its Part~II (Theorem \texttt{pqc:qm:thm:47}: six integer $3\times3$ mortality matrices compiled into seven rational $4\times4$ Kraus operators, zero-probability word $\Leftrightarrow$ mortal); Corollary~\ref{ste:ee:cor:seven} below is its positive stochastic counterpart with a different exact predicate. Neither that report nor the research notes contain the no-common-line compiler, the idempotent PCP connector, the stochastic zero--one law or the compressed quartic. Section~\ref{ste:sec:relation} gives the comparisons.

\section{Exact forgetting and the zero-mass representation}
\label{ste:ee:sec:setup}

All stochastic matrices in this article are \emph{column-stochastic}. Probability
distributions are columns, so $p\mapsto Sp$. We set $\N=\{0,1,2,\ldots\}$.
Matrices called positive have every entry strictly greater than zero. Integer
matrix mortality allows negative entries in the source matrices.

For a word $w=i_1\cdots i_T$ in chronological order, define
\[
 S_w=S_{i_T}\cdots S_{i_1},\qquad S_\epsilon=I.
\]
The same convention applies to every matrix alphabet. A nonempty word is required
for mortality and erasure.

\begin{definition}
An \emph{erasing word} for $S_1,\ldots,S_k$ is a nonempty word $w$ for which
$S_wp$ is the same distribution for every initial distribution $p$. A
\emph{uniform erasing word} satisfies $S_w=J_n$, where
\[
 u_n=\frac1n\one_n,\qquad J_n=u_n\one_n\transpose.
\]
\end{definition}

\begin{lemma}[Equivalent erasure tests]\label{ste:ee:lem:erasure}
For a nonnegative column-stochastic $n\times n$ matrix $S$, the following are
equivalent: $S$ erases all initial distributions; all its columns are equal;
$\rk S=1$; and $S$ vanishes on
\[
 H_n=\{x\in\R^n:\one_n\transpose x=0\}.
\]
If $S$ is also row-stochastic, these conditions are equivalent to $S=J_n$.
\end{lemma}
\begin{proof}
The images of the point-mass distributions are precisely the columns of $S$.
Thus erasure is equivalent to equal columns. Equal columns have rank one, not
zero, since their sums are one. Conversely, rank one makes the columns scalar
multiples of a fixed nonzero vector; their equal nonzero column sums force all
these multiples to be equal. The vectors $e_j-e_1$, $j>1$, span $H_n$, giving the
hyperplane condition. Finally, a matrix with equal columns has the form
$p\one\transpose$; its row sums are $np_a$, which are all one precisely when
$p=u_n$.
\end{proof}

\begin{lemma}[Erasure persists]\label{ste:ee:lem:persistence}
If a word $w$ is erasing, every word containing $w$ as a contiguous block is
erasing. Appending further letters after erasure can change the common output
distribution, but cannot restore dependence on the initial distribution.
\end{lemma}
\begin{proof}
Multiplication on either side cannot increase matrix rank. Every product remains
column-stochastic and so has rank at least one. A product containing the rank-one
factor $S_w$ therefore has rank exactly one.
\end{proof}

To separate total mass from information, let $h=n-1$ and introduce
\begin{equation}
 E=\begin{pmatrix}I_h\\-\one_h\transpose\end{pmatrix},\qquad
 F=[I_h\ 0]-\frac1n\one_h\one_n\transpose,\qquad F_0=nF.
 \label{ste:ee:eq:EF}
\end{equation}
These are explicit rational coordinate maps, not existentially chosen bases.

\begin{lemma}[Coordinates]\label{ste:ee:lem:coordinates}
One has
\[
 FE=I_h,\quad Fu_n=0,\quad \one_n\transpose E=0,\quad
 EF=I_n-J_n.
\]
Every $x\in\Q^n$ has the unique decomposition
\[
 x=Ey+s u_n,\qquad y=Fx,\quad s=\one_n\transpose x.
\]
For any rational column-stochastic $S$, its matrix in these coordinates is
\begin{equation}
 \begin{pmatrix}C&b\\0&1\end{pmatrix},\qquad
 C=FSE,\quad b=FSu_n,
 \label{ste:ee:eq:generalblock}
\end{equation}
and $\rk S=1+\rk C$.
\end{lemma}
\begin{proof}
The displayed identities follow by multiplying the matrices in
\eqref{ste:ee:eq:EF}. They prove existence and uniqueness of the decomposition.
Column-stochasticity preserves $s$, while applying $F$ to $S(Ey+s u_n)$ gives
$Cy+sb$. To compute the rank, subtract from each of the first $h$ output rows its
last-column entry times the last row. This leaves the block diagonal matrix
$\diag(C,1)$ and hence rank $1+\rk C$.
\end{proof}

The same algebra holds over $\C$. We will use the complexified hyperplane and
coordinates when excluding invariant lines. The distinction between a common
invariant subspace and a common eigenvector is important: a one-dimensional
invariant subspace is generated by a common eigenvector, but its eigenvalue may
depend on the generator.

\section{An explicit positive affine lift}
\label{ste:ee:sec:lift}

Let $(C_i,b_i)$, $1\leq i\leq k$, be finitely many integer affine data, with
$C_i\in\Z^{h\times h}$ and $b_i\in\Z^h$. Fix a rational number
$0<\eta<1$, chosen independently of the input coefficients. Define
\begin{align}
 \Delta_i&=EC_iF_0+nEb_i\one_n\transpose,\label{ste:ee:eq:delta}\\
 M&=\max_{i,a,b}|(\Delta_i)_{ab}|,\qquad
 K=\max\{2,1+\lceil M/\eta\rceil\},\qquad D=nK,\label{ste:ee:eq:scale}\\
 B_i&=K\one_n\one_n\transpose+\Delta_i,\qquad S_i=B_i/D.
 \label{ste:ee:eq:positiveB}
\end{align}
Here $B_i$ and $D$ are integers. The same $K$ and $D$ are used for every generator.

\begin{theorem}[Positive affine lift]\label{ste:ee:thm:lift}
The matrices in \eqref{ste:ee:eq:positiveB} are positive rational column-stochastic and
satisfy the strict entrywise bounds
\begin{equation}
 \frac{1-\eta}{n}<(S_i)_{ab}<\frac{1+\eta}{n}.
 \label{ste:ee:eq:entrybounds}
\end{equation}
In the coordinates of Lemma~\ref{ste:ee:lem:coordinates},
\begin{equation}
 S_i\sim\begin{pmatrix}K^{-1}C_i&K^{-1}b_i\\0&1\end{pmatrix}.
 \label{ste:ee:eq:liftblock}
\end{equation}
For every word $w$ of length $T$,
\begin{equation}
 \rk S_w=1+\rk C_w.
 \label{ste:ee:eq:affine-rank}
\end{equation}
All entries of $B_i$ are strictly positive integers and all its column sums equal
$D$.
\end{theorem}
\begin{proof}
Since $K>M/\eta$, each entry of $\Delta_i/(nK)$ has absolute value less than
$\eta/n$. Adding $J_n$ proves \eqref{ste:ee:eq:entrybounds} and positivity.
The column sums of $\Delta_i$ vanish because $\one\transpose E=0$, so the column
sums of $S_i$ are one and those of $B_i$ are $D$.
Using $F_0=nF$, rewrite the lift as
\[
 S_i=J_n+K^{-1}EC_iF+K^{-1}Eb_i\one_n\transpose.
\]
On $Ey+s u_n$ this gives
\[
 S_i(Ey+s u_n)=E\bigl(K^{-1}C_i y+K^{-1}b_i s\bigr)+s u_n,
\]
which proves \eqref{ste:ee:eq:liftblock}. Multiplying these upper block-triangular
matrices gives a top-left block $K^{-T}C_w$. Its translation block need not vanish,
but Lemma~\ref{ste:ee:lem:coordinates} shows that it does not affect the rank. Since
$K\ne0$, \eqref{ste:ee:eq:affine-rank} follows.
\end{proof}

\begin{proposition}[Uniform approximate forgetting]\label{ste:ee:prop:contraction}
For all probability distributions $p,q$ and every word $w$ of length $T$,
\begin{equation}
 \TV(S_wp,S_wq)\leq\eta^T\TV(p,q),\qquad
 \TV(p,q)=\tfrac12\|p-q\|_1.
 \label{ste:ee:eq:contraction}
\end{equation}
\end{proposition}
\begin{proof}
The lower bound in \eqref{ste:ee:eq:entrybounds} allows the decomposition
\[
 S_i=(1-\eta)J_n+\eta Q_i,
\]
where $Q_i$ is nonnegative and column-stochastic. Such a matrix is nonexpansive in
$\ell^1$: for any real vector $z$,
\[
 \|Q_i z\|_1\leq\sum_a\sum_b(Q_i)_{ab}|z_b|=\|z\|_1.
\]
For $z=p-q$ one has $J_nz=0$, giving a factor $\eta$ in one step. Iteration proves
the result.
\end{proof}

In particular, for a fixed accuracy $\varepsilon>0$, a horizon with
$\eta^T\leq\varepsilon$ works for every input instance and every schedule. This
is a bound on the difference between two trajectories driven by the same
schedule. When translations are present, it is \emph{not} a claim that arbitrary
switching converges to one common stationary distribution.

\begin{corollary}[Doubly stochastic special case]\label{ste:ee:cor:double}
For $A_1,\ldots,A_m\in\Z^{d\times d}$, apply the lift with $h=d$, $C_i=A_i$,
and $b_i=0$. Then $S_i$ is doubly stochastic, and
\begin{equation}
 S_w=J_{d+1}+K^{-T}EA_wF.
 \label{ste:ee:eq:doublyproduct}
\end{equation}
Thus $A_w=0$ if and only if $S_w$ has rank one, if and only if $S_w=J_{d+1}$.
\end{corollary}
\begin{proof}
The zero translations imply $S_i u_n=u_n$, which is row-stochasticity. The block
product has zero translation and top-left block $K^{-T}A_w$. Conjugating it back
gives \eqref{ste:ee:eq:doublyproduct}; the equivalences follow from $FE=I_d$.
\end{proof}

The baseline lift alone is a direct transport of mortality, not the main
no-line result: its generators share the stationary line $\C u_n$.

\section{Two guards eliminate every common invariant line}
\label{ste:ee:sec:linefree}

The obstacle in the baseline is structural. Merely adding an affine term need
not remove eigenvectors already shared by the original matrices inside $H_n$.
We therefore reserve a tensor factor for two invertible guards.

Set
\begin{equation}
 R=\begin{pmatrix}1&1\\0&1\end{pmatrix},\qquad
 L=\begin{pmatrix}1&0\\1&1\end{pmatrix}.
 \label{ste:ee:eq:guards}
\end{equation}
Given $A_1,\ldots,A_m\in\Z^{d\times d}$, take $h=2d$, $n=2d+1$, and use the
following $m+2$ affine data:
\begin{align}
 (C_i,b_i)&=(A_i\otimes I_2,0), &&1\leq i\leq m,\label{ste:ee:eq:sourcepairs}\\
 (C_R,b_R)&=(I_d\otimes R,0),\label{ste:ee:eq:Rpair}\\
 (C_L,b_L)&=(I_d\otimes L,e_1).\label{ste:ee:eq:Lpair}
\end{align}
The vector $e_1$ in \eqref{ste:ee:eq:Lpair} is the first standard basis vector of
$\Z^{2d}$. Apply the common-scale positive affine lift to all these pairs.

For a word $w$ in the augmented alphabet, let $\pi(w)$ be its source-letter
subword and $\gamma(w)$ its guard-letter subword. Both retain chronological
order. An empty source subword has product $I_d$.

\begin{theorem}[Line-free mortality compiler]\label{ste:ee:thm:linefree}
The construction \eqref{ste:ee:eq:sourcepairs}--\eqref{ste:ee:eq:Lpair} is effective and has the
following properties.
\begin{enumerate}[label=(\roman*)]
\item All $m+2$ generators are positive rational column-stochastic matrices of
size $2d+1$, satisfying \eqref{ste:ee:eq:entrybounds} and \eqref{ste:ee:eq:contraction} for the
prescribed $\eta$.
\item The generators have no common invariant one-dimensional subspace of
$\C^{2d+1}$.
\item Every word $w$ satisfies
\begin{equation}
  C_w=A_{\pi(w)}\otimes G_{\gamma(w)},\qquad
  \rk S_w=1+2\rk A_{\pi(w)},
  \label{ste:ee:eq:tensorrank}
\end{equation}
where $G_{\gamma(w)}$ is a product of $R,L$ and is invertible.
\item The source family is mortal if and only if the target family has an
erasing word, if and only if it has a uniform erasing word.
\end{enumerate}
The same rank-one existence equivalence and no-common-line property hold for the
positive integer numerators $B_1,\ldots,B_{m+2}$.
\end{theorem}
\begin{proof}
The first assertion is Theorem~\ref{ste:ee:thm:lift} and
Proposition~\ref{ste:ee:prop:contraction}.

For the tensor identity, each source linear part acts only on the first tensor
factor and each guard only on the second. Thus source factors commute with guard
factors. Moving source factors past guards, without changing the order within
either subword, gives the first identity in \eqref{ste:ee:eq:tensorrank}. Both $R$ and
$L$ have determinant one. Therefore $G_{\gamma(w)}$ is invertible, including
$G_\epsilon=I_2$. Multiplying on the second factor by its inverse shows
\[
 \rk\bigl(A_{\pi(w)}\otimes G_{\gamma(w)}\bigr)
 =\rk\bigl(A_{\pi(w)}\otimes I_2\bigr)=2\rk A_{\pi(w)}.
\]
The affine rank formula now proves \eqref{ste:ee:eq:tensorrank}.

We next rule out a common invariant complex line $\ell$. There are two cases.
If $\ell\subseteq H_n^{\C}$, use the $E$ coordinates. A nonzero vector spanning
it would be an eigenvector of both $I_d\otimes R$ and $I_d\otimes L$; scaling
both by $K^{-1}$ does not affect their eigenvectors. The only eigenvalue of each
unipotent guard is one. Their eigenspaces are respectively
\[
 \C^d\otimes\operatorname{span}(e_1)
 \quad\hbox{and}\quad
 \C^d\otimes\operatorname{span}(e_2),
\]
whose intersection is zero. This is impossible.

If $\ell\not\subseteq H_n^{\C}$, choose $v\in\ell$ with
$\one\transpose v=1$. Every column-stochastic generator acts on $v$ by its
line eigenvalue, and preservation of total mass forces that eigenvalue to be
one. Write $v=Ey+u_n$. The unshifted $R$ guard gives
\[
 y=K^{-1}(I_d\otimes R)y.
\]
The matrix on the right has only eigenvalue $K^{-1}\ne1$, since $K\geq2$.
Consequently $y=0$ and $v=u_n$. But the shifted $L$ guard sends
\[
 u_n\longmapsto u_n+K^{-1}Ee_1\ne u_n,
\]
contradicting invariance. This proves (ii) over $\C$, not just over $\Q$.

If $A_v=0$ for a nonempty source word $v$, use the same word without guards.
Its affine translations are zero, so its lifted product is exactly $J_n$.
Conversely, if $S_w$ has rank one, \eqref{ste:ee:eq:tensorrank} gives
$A_{\pi(w)}=0$. The projection cannot be empty, since $A_\epsilon=I_d\ne0$.
This establishes (iv).

Finally, each $B_i=DS_i$ with a common nonzero scalar $D$. Scaling a generator
by a nonzero scalar preserves all its invariant subspaces, and
$B_w=D^{|w|}S_w$ preserves every product rank. The integer conclusion follows.
\end{proof}

\begin{remark}[Why the guard arguments must both be present]
The two unipotent guards without an affine displacement would still fix $u_n$.
An affine displacement without a protected tensor factor could leave an
invariant line inside the information hyperplane. Our two cases address these
independent obstructions. Invertibility of the guards is equally essential:
a singular guard could destroy source information and introduce false erasing
words.
\end{remark}

\subsection{Small numerical consequences}

Cassaigne et al. prove undecidability of mortality for the four parameter pairs
$(d,m)=(3,6),(5,4),(9,3),(15,2)$ \cite{cassaigne}. Applying the two explicit lifts
gives Table~\ref{ste:ee:tab:bounds}. The entries are sufficient bounds transported from
that paper, not a claim to list an optimal or exhaustive current frontier.
(\wtag\ Part~II of~\cite{pqc} imports the same six-matrix $3\times3$ bound from Neary~\cite[Corollary~12]{neary}, through binary tag systems and the Post correspondence problem for five pairs of words. Both attributions are kept; this report does not decide which publication first proved the $(3,6)$ bound.)

\begin{table}[ht]
\centering
\begin{tabular}{@{}ccc@{}}
\toprule
Source mortality $(d,m)$ & Doubly stochastic $(n,k)$ & No common line $(n,k)$\\
\midrule
$(3,6)$ & $(4,6)$ & $(7,8)$\\
$(5,4)$ & $(6,4)$ & $(11,6)$\\
$(9,3)$ & $(10,3)$ & $(19,5)$\\
$(15,2)$ & $(16,2)$ & $(31,4)$\\
\bottomrule
\end{tabular}
\caption{Undecidable exact-erasure families. Generator counts are upper bounds.
All target entries may be confined to any prescribed band
$((1-\eta)/n,(1+\eta)/n)$ with rational $0<\eta<1$.}
\label{ste:ee:tab:bounds}
\end{table}

\begin{corollary}\label{ste:ee:cor:seven}
For any fixed rational $0<\eta<1$, it is undecidable whether a finite product of
at most eight positive rational $7\times7$ column-stochastic matrices has rank
one, even under both promises that the entrywise band
\eqref{ste:ee:eq:entrybounds} holds and that the generators have no common invariant
complex line. The corresponding rank-one problem for positive integer matrices
with one common positive column sum is also undecidable.
\end{corollary}
\begin{proof}
Apply Theorem~\ref{ste:ee:thm:linefree} to the $(3,6)$ mortality family. An algorithm on
the promised target class would decide every source instance, since the compiler
always meets the promises and preserves the answer.
\end{proof}

This is a family of finite alphabets parameterized by the source input. It is
not a claim that the displayed three-state example below, or one fixed numerical
eight-generator alphabet, simulates every program. Rational precision is
unbounded across input instances; fixing all numerator and denominator bounds
would leave only finitely many instances.

\wtag\ \emph{Relation to the nonnegative-mortality boundary.} Part~II of~\cite{pqc} proves that mortality is decidable for entrywise nonnegative matrices (Proposition \texttt{pqc:qm:prop:nonnegative}) and asks, in Question \texttt{pqc:q:signed}, for further decidable classes, among them families with a common invariant cone. The generators of Corollary~\ref{ste:ee:cor:seven} are strictly positive and share the invariant cone of nonnegative vectors, yet their rank-one problem is undecidable: the positivity boundary for the zero predicate does not transfer to rank one, because rank one in dimension $n$ is mortality of the signed induced matrices in dimension $n-1$ (Proposition~\ref{ste:ee:prop:dimension}).

\section{A direct correspondence-problem compiler}
\label{ste:ee:sec:pcp}

To separate completeness from the numerical mortality bounds, we now give a
standalone many-one reduction. The starting problem is PCP: for a finite list of
pairs $(u_i,v_i)$, does some nonempty index word satisfy
\[
 u_{i_1}\cdots u_{i_T}=v_{i_1}\cdots v_{i_T}?
\]
We use its classical computably enumerable completeness. The halting-to-PCP
reduction is an external dependency, not re-proved here; see the formal treatment
of PCP-related reductions by Forster, Heiter, and Smolka \cite{forster}.
Restriction to a two-letter alphabet does not lose generality: encode a finite
alphabet injectively by equal-length binary codewords, which preserves equality
of all concatenations.

\subsection{A four-dimensional integer front end}

Use letters $a,b$ as digits $1,2$ in base three, read most significant digit first.
Write $\code(s)$ for the resulting nonnegative integer, including
$\code(\epsilon)=0$. This encoding is injective, and
\begin{equation}
 \code(st)=3^{|t|}\code(s)+\code(t).
 \label{ste:ee:eq:stringcode}
\end{equation}
Define
\begin{equation}
 M_i=\begin{pmatrix}
 3^{|u_i|}&0&\code(u_i)&0\\
 0&3^{|v_i|}&\code(v_i)&0\\
 0&0&1&0\\
 0&0&0&0
 \end{pmatrix},\qquad
 c=\begin{pmatrix}0\\0\\1\\1\end{pmatrix},\qquad
 \ell=\begin{pmatrix}1&-1&0&1\end{pmatrix},\qquad Q=c\ell.
 \label{ste:ee:eq:pcp-matrices}
\end{equation}
The letter $q$ denotes the extra generator $Q$. It is deliberately idempotent:
$\ell c=1$ and $Q^2=Q\ne0$.

\begin{lemma}[PCP block identity]\label{ste:ee:lem:pcpblock}
For a nonempty source word $v=i_1\cdots i_T$, let
$U_v=u_{i_1}\cdots u_{i_T}$ and $V_v=v_{i_1}\cdots v_{i_T}$. Then
\[
 M_v=\begin{pmatrix}
 3^{|U_v|}&0&\code(U_v)&0\\
 0&3^{|V_v|}&\code(V_v)&0\\
 0&0&1&0\\
 0&0&0&0
 \end{pmatrix},\qquad
 \ell M_vc=\code(U_v)-\code(V_v).
\]
The corresponding scalar for the empty word is $\ell I_4c=1$, not zero.
\end{lemma}
\begin{proof}
The product convention is chronological. Multiplying a later tile on the left
updates the first translation by
$3^{|u_i|}\code(U)+\code(u_i)=\code(Uu_i)$, and similarly for the second.
Induction proves the matrix formula. The scalar follows by direct multiplication.
For the empty word use $M_\epsilon=I_4$; the block formula with a zero final
diagonal entry is asserted only for a nonempty word.
\end{proof}

\begin{theorem}[Full-word PCP reduction]\label{ste:ee:thm:pcp}
The integer alphabet $M_1,\ldots,M_m,Q$ is mortal if and only if the PCP instance
has a solution. More precisely, write any word containing $s\geq1$ connector
letters as
\[
 w=v_0qv_1q\cdots qv_s,
\]
where each $v_j$ is a possibly empty source word. Its product is zero if and only
if at least one internal block $v_j$, $1\leq j<s$, is a nonempty PCP solution.
Words without a connector never have zero product.
\end{theorem}
\begin{proof}
By $Q=c\ell$ and the chronological convention,
\begin{equation}
 M_w=(M_{v_s}c)
 \left(\prod_{j=1}^{s-1}\ell M_{v_j}c\right)(\ell M_{v_0}).
 \label{ste:ee:eq:pcp-factor}
\end{equation}
Both endpoint vectors are nonzero. For a nonempty $v_s$, the third coordinate of
$M_{v_s}c$ is one; for an empty block it is the nonzero vector $c$. For nonempty
$v_0$, the first coordinate of $\ell M_{v_0}$ is $3^{|U_{v_0}|}>0$; for an empty
block it is $\ell\ne0$. Thus the outer product in \eqref{ste:ee:eq:pcp-factor} is
nonzero, and vanishing occurs exactly when an internal scalar is zero.
Lemma~\ref{ste:ee:lem:pcpblock} identifies those zero scalars with nonempty PCP matches;
an empty internal block contributes one. Without $Q$, the third diagonal entry
of every nonempty product is one, excluding mortality.

Conversely, a PCP match $v$ gives the zero word $qvq$. The entire map from tiles
to matrices is an explicit computable map, so this is a many-one reduction.
\end{proof}

The fourth coordinate has a specific purpose. A more naive connector based on a
scalar comparison that vanishes for the empty word would have $Q^2=0$, making
every instance mortal. Here the extra coordinate is killed by the first genuine
tile but makes the empty-block scalar one. This is a boundary condition, not a
cosmetic padding coordinate.

\wtag\ The three-dimensional radix-three tiles are the classical encoding that the research note~\cite{matrixpcp} also uses, there with a leading sentinel digit so that the empty word has a nonzero code. The fourth coordinate and the idempotent connector $Q=c\ell$, $\ell c=1$, are this source's own.

\begin{corollary}[A many-one complete stochastic substrate]\label{ste:ee:cor:complete}
For any fixed rational $0<\eta<1$, exact erasure is computably enumerable complete
under many-one reductions for positive rational $9\times9$ column-stochastic
families with a variable-size alphabet, even with no common invariant complex
line and with the entrywise band \eqref{ste:ee:eq:entrybounds}.
\end{corollary}
\begin{proof}
A PCP instance with $m$ tiles maps by Theorem~\ref{ste:ee:thm:pcp} to $m+1$ integer
$4\times4$ matrices, then by Theorem~\ref{ste:ee:thm:linefree} to $m+3$ stochastic
$9\times9$ matrices. Both maps are many-one and preserve the exact answer.
Erasure is computably enumerable: enumerate all finite words and multiply their
rational matrices exactly. This proves completeness on the promised class.
Equivalently, one may regard the source tile list as the input parameter of the
compiled family, so promise validation is not part of the reduction.
\end{proof}

This gives a precise sense in which the substrate supports universal finite
computation. A halting instance can be transformed into a finite rational gate
alphabet such that a finite exact-forgetting schedule exists exactly when the
computation halts. It does not supply an autonomous universal stochastic machine
with one fixed gate alphabet and a separate ordinary input loader.

\section{Excluding all small common invariant subspaces}
\label{ste:ee:sec:higher}

The tensor guard method is not limited to lines. Let $r\geq2$, and let $P$ be the
cyclic permutation matrix on $\C^r$, while
\[
 Z_r=\diag(1,2,\ldots,r).
\]
For an integer mortality instance of dimension $d$, use $h=rd$, $n=rd+1$, and
\begin{equation}
 (A_i\otimes I_r,0),\qquad (I_d\otimes P,0),\qquad
 (I_d\otimes Z_r,e_1).
 \label{ste:ee:eq:higherpairs}
\end{equation}
Apply the same common-scale lift. The name $Z_r$ here denotes a diagonal matrix,
not an integer witness domain.

\begin{lemma}[The guard algebra]\label{ste:ee:lem:guardalgebra}
The complex unital algebra generated by $P$ and $Z_r$ is
$\Mat_r(\C)$. Every subspace of $\C^d\otimes\C^r$ invariant under
$I_d\otimes P$ and $I_d\otimes Z_r$ has the form $W\otimes\C^r$ for some
$W\leq\C^d$.
\end{lemma}
\begin{proof}
Since the diagonal entries of $Z_r$ are distinct, Lagrange interpolation produces
the coordinate projections:
\[
 E_{jj}=\prod_{a\ne j}\frac{Z_r-aI}{j-a}.
\]
For each $i,j$ choose a power $P^t$ taking $e_j$ to $e_i$. Then
$E_{ii}P^tE_{jj}=E_{ij}$, so the generated algebra contains every matrix unit.

A subspace invariant under the two generators is invariant under their products
and complex linear combinations, hence under all $I_d\otimes E_{ij}$. Define
$W=\{v:v\otimes e_1\text{ belongs to the subspace}\}$. Applying
$I_d\otimes E_{1j}$ to $\sum_jv_j\otimes e_j$ shows that every coefficient
$v_j$ belongs to $W$. Conversely, applying $I_d\otimes E_{j1}$ shows that
$W\otimes\C^r$ is contained in the subspace. This proves equality.
\end{proof}

\begin{theorem}[Invariant-subspace and product-rank gaps]\label{ste:ee:thm:higher}
For the lift of \eqref{ste:ee:eq:higherpairs}, every common invariant complex subspace
has dimension congruent to $0$ or $1$ modulo $r$. There is no common invariant
line and, consequently, no nonzero common invariant subspace of any dimension
$1,\ldots,r-1$. Every product satisfies
\begin{equation}
 \rk S_w=1+r\rk A_{\pi(w)}.
 \label{ste:ee:eq:higherrank}
\end{equation}
In particular, a product is either rank one or has rank at least $r+1$.
Mortality, erasure, and existence of a uniform erasing word remain equivalent.
\end{theorem}
\begin{proof}
Both $P$ and $Z_r$ are invertible. The source/guard tensor factorization therefore
proves \eqref{ste:ee:eq:higherrank} exactly as before. Zero translations on source-only
words give the uniform-target equivalence.

Let $U$ be a common invariant subspace. Its intersection with $H_n^{\C}$ is
invariant under the linear guard restrictions and has dimension divisible by
$r$, by Lemma~\ref{ste:ee:lem:guardalgebra}. If $U\subseteq H_n^{\C}$, this is its
full dimension. Otherwise, restriction of the nonzero mass functional to $U$
has a one-dimensional image, so
$\dim U=\dim(U\cap H_n^{\C})+1$. This proves the congruence statement.

A common invariant line cannot be inside $H_n^{\C}$ because its dimension would
be a positive multiple of $r$. Outside $H_n^{\C}$ normalize a spanning vector to
mass one. The unshifted $P$ guard forces its transverse part $y$ to satisfy
$y=K^{-1}(I_d\otimes P)y$. Every eigenvalue of $P$ has modulus one, whereas
$K>1$, so $y=0$. The translated $Z_r$ guard does not fix $u_n$, again a
contradiction. Excluding dimension one and using the congruence statement
excludes all positive dimensions smaller than $r$.
\end{proof}

The rank gap makes the exactness phenomenon sharper: products cannot drift
through ranks $2,\ldots,r$ on their way to erasure. Their entries can nevertheless
approach a rank-one matrix arbitrarily closely without ever becoming rank one.
The common hyperplane of dimension $rd$ remains, so the theorem must not be
reported as an irreducible-semigroup result.

For example, applying this construction to the $d=3$, $m\leq6$ mortality bound
gives undecidable erasure for at most eight positive stochastic matrices in
dimension $3r+1$, with no invariant subspace of dimensions $1,\ldots,r-1$.
Applying it to the explicit PCP front end gives many-one completeness in
dimension $4r+1$ with variable alphabet size.

\section{Decidable boundaries and a sharp commuting horizon}
\label{ste:ee:sec:decidable}

\subsection{An exact dimension correspondence}

\begin{proposition}\label{ste:ee:prop:dimension}
For $n\geq2$, rank-one product existence for rational column-stochastic
$n\times n$ matrices reduces effectively, without changing the number of
generators, to mortality for integer $(n-1)\times(n-1)$ matrices. Conversely,
integer mortality in dimension $n-1$ reduces to erasure for positive rational
doubly stochastic matrices of dimension $n$, in any prescribed band
\eqref{ste:ee:eq:entrybounds}.
\end{proposition}
\begin{proof}
For the first direction, compute $C_i=FS_iE$ using the explicit rational
coordinates of Lemma~\ref{ste:ee:lem:coordinates}. The restriction of $S_w$ to $H_n$
has coordinate matrix $C_w$, so erasure is equivalent to $C_w=0$. Clear the
rational denominators in each $C_i$, or use one common positive denominator.
Every product is merely multiplied by a nonzero scalar, preserving zero
products. The converse is Corollary~\ref{ste:ee:cor:double}.
\end{proof}

This correspondence is useful independently of particular known dimension
bounds. For example, the three-state positive doubly stochastic erasure problem
has exactly the corresponding two-dimensional mortality difficulty. We do not
infer a current open/closed classification of two-dimensional mortality from a
dated source; the reduction equivalence itself is the assertion here.

\subsection{Commuting information dynamics}

The affine stochastic generators need not commute in the next theorem. It is
enough that their restrictions to the zero-mass hyperplane commute.

\begin{theorem}[Sharp commuting horizon]\label{ste:ee:thm:commuting}
Let $S_1,\ldots,S_k$ be rational column-stochastic $n\times n$ matrices,
$n\geq2$. Suppose that the induced $(n-1)$-dimensional matrices
$C_i=FS_iE$ commute pairwise. If an erasing word exists, one of length at most
$n-1$ exists. A finite exact rank-reduction algorithm finds such a word or
certifies that no erasing word exists. The bound $n-1$ is sharp, even for a
single positive doubly stochastic generator.
\end{theorem}
\begin{proof}
Set $h=n-1$ and start with $P=I_h$, $V=\im P=\Q^h$. Since $P$ is a product of
commuting $C_i$, every $C_i$ leaves $V$ invariant. If some generator satisfies
$\rk(C_iP)<\rk P$, replace $P$ by $C_iP$, append that letter, and continue.
Each such step strictly decreases the nonnegative integer $\dim V$, so at most
$h$ steps are possible.

If $P=0$, the chosen word is erasing. Otherwise suppose that no generator lowers
the rank. Every $C_i:V\to V$ then has full rank on $V$ and is invertible there.
All products remain invertible on this nonzero space, so none can be the zero
matrix on $\Q^h$. Thus failure to find a further rank-lowering generator is a
correct no-erasure certificate. Conversely, if any zero product exists it
annihilates $V$, which is impossible while all $C_i|_V$ are invertible. Hence
the algorithm cannot stop at a nonzero $V$ in a mortal instance.

For sharpness, take the $h\times h$ nilpotent shift $N$ with ones immediately
above the diagonal. It has $N^{h}=0$ and $N^t\ne0$ for $0\leq t<h$. The
single-generator positive doubly stochastic lift of $N$ therefore first erases
at time $h=n-1$.
\end{proof}

The algorithm uses rational matrix multiplication and Gaussian rank tests; it
has no unbounded word search. The proof does not assume that a rank reduction
already makes a matrix zero, and it allows arbitrary choices among available
rank-lowering generators. The common-invariant-image property provided by
commutation is what makes these greedy choices safe.

For one generator, erasure is exactly nilpotency of $C$, and $C^{n-1}=0$ is a
complete test. For $n=2$, every induced $C_i$ is a scalar, so erasure occurs if and
only if at least one generator already has rank one. These are genuine finite
boundaries, not numerical heuristics for the undecidable cases.

\section{Random schedules, exact forgetting, and a hidden-state interpretation}
\label{ste:ee:sec:random}

Now choose each letter independently with probabilities
$p_1,\ldots,p_k>0$, $\sum_i p_i=1$. Define the exact-forgetting time
\[
 \tau=\inf\{T\geq1:\rk(S_{I_T}\cdots S_{I_1})=1\},
\]
with $\inf\varnothing=\infty$. The event concerns exact equality, not a
numerical distance threshold.

\begin{theorem}[A zero--one law for finite exact forgetting]\label{ste:ee:thm:iid}
For a finite column-stochastic alphabet under independent full-support
switching, the following conditions are equivalent:
\begin{enumerate}[label=(\roman*)]
\item an erasing word exists;
\item $\PP(\tau<\infty)>0$;
\item $\PP(\tau<\infty)=1$;
\item $\EE\tau<\infty$.
\end{enumerate}
If $v$ is an erasing word of length $L$ and
$q=\prod_{j=1}^L p_{v_j}>0$, then
\begin{equation}
 \PP(\tau>t)\leq(1-q)^{\lfloor t/L\rfloor},\qquad
 \EE\tau\leq L/q.
 \label{ste:ee:eq:tail}
\end{equation}
For uniform letter probabilities, $q=k^{-L}$.
\end{theorem}
\begin{proof}
If no erasing word exists, no finite prefix can have rank one, so $\tau=\infty$
for every schedule. Suppose instead that an erasing word $v$ of length $L$
exists. Partition the random schedule into successive disjoint blocks of length
$L$. Each block equals $v$ with probability $q$, independently of the others.
The first matching block occurs after a geometric number $G$ of blocks, with
$\EE G=1/q$. By Lemma~\ref{ste:ee:lem:persistence}, regardless of the preceding prefix,
that block makes the entire prefix product rank one. Hence $\tau\leq LG$, which
proves \eqref{ste:ee:eq:tail} and finite expectation. Finite expectation implies
almost-sure finiteness, and almost-sure finiteness implies positive probability.
These implications establish all equivalences.
\end{proof}

\begin{corollary}\label{ste:ee:cor:randomhard}
On the seven-state promised class of Corollary~\ref{ste:ee:cor:seven}, both almost-sure
finite exact forgetting and finiteness of its expected time are undecidable.
They are many-one complete on the variable-alphabet nine-state family of
Corollary~\ref{ste:ee:cor:complete}. Nevertheless, the deterministic approximation
bound \eqref{ste:ee:eq:contraction} holds for every schedule on every instance.
\end{corollary}
\begin{proof}
Theorem~\ref{ste:ee:thm:iid} identifies each of these random-schedule predicates with
existence of a finite erasing word. The reductions and contraction estimates
already proved therefore apply unchanged.
\end{proof}

There is no contradiction with the higher logical complexity of many other
almost-sure termination questions. Here a single finite erasing block is an
absorbing certificate that recurs in independent full-support sampling. That
special structure collapses the almost-sure predicate to an existential
finite-word predicate. Removing independence, full support, or the absorbing
rank-one property changes the argument.

\wtag\ Part~IX of~\cite{pqc} (Theorem \texttt{pqc:iq:thm:geometric}) also combines a geometric stopping clock with an undecidable exact event, for completed records of rational quantum instruments; the statements are different and neither implies the other.

\subsection{A labelled hidden-state model with trivial output statistics}

The stochastic interpretation can be made into a finite labelled transition
model. For label $i$, use the nonnegative transition matrix
\[
 T_i=p_iS_i.
\]
Since $\sum_iT_i$ is column-stochastic, these matrices specify joint probabilities
of the next hidden state and the emitted label. Conditional on any current hidden
state, label $i$ has probability $p_i$, because
$\one\transpose T_i=p_i\one\transpose$. Thus the emitted labels are independent
with the prescribed common distribution.

After observing a word $w=i_1\cdots i_T$, its probability is
$\prod_t p_{i_t}$, independent of the initial hidden distribution. The normalized
hidden-state update is exactly $S_w$. Consequently, exact independence of the
posterior from the initial hidden distribution is precisely rank-one product
existence.

The output entropy rate in this particular construction is simply
$-\sum_i p_i\log p_i$. It is not an undecidable quantity. This observation is
important: even when output statistics are completely elementary, deciding
whether filtering ever forgets its initial distribution \emph{exactly} can be
undecidable. The result addresses exact filtering, not the entropy question that
motivated the MathOverflow discussion.

\subsection{Finite-horizon probabilities}

Under uniform switching, let $N_T$ be the number of labelled words of length $T$
with rank-one product. By persistence,
\begin{equation}
 \PP(\tau\leq T)=N_T/k^T.
 \label{ste:ee:eq:probcount}
\end{equation}
This is an exact rational number computable by finite enumeration. The next
section realizes $N_T$ as the number of natural solutions of an explicit quartic.
Duplicate numerical generators, when present, remain distinct labels; this
matches the probability model and the certificate count.

\section{Explicit natural-domain Diophantine certificates}
\label{ste:ee:sec:quartic}

Fix a numerical input alphabet $B_1,\ldots,B_k\in\Z_{>0}^{n\times n}$, each
with column sum $D>0$, and set $S_i=B_i/D$. In this section $T\geq1$ is an
externally supplied horizon. A polynomial is produced for that horizon; its
arity is not fixed as $T$ varies.

\subsection{A full-prefix certificate with no minors}

Introduce natural selector variables $e_{t,i}$, $1\leq t\leq T$, $1\leq i\leq k$,
and natural prefix matrices $X_t\in\N^{n\times n}$. The matrix $X_0=I_n$ is
constant. Use residuals
\begin{align}
 s_t&=\sum_{i=1}^k e_{t,i}-1,\label{ste:ee:eq:selector}\\
 r_{t,ab}&=(X_t)_{ab}-\sum_{i=1}^k e_{t,i}
                   \sum_{c=1}^n(B_i)_{ac}(X_{t-1})_{cb},\label{ste:ee:eq:prefix}\\
 q_{ab}&=(X_T)_{ab}-(X_T)_{a1},\quad 1\leq a\leq n,\quad 2\leq b\leq n.
 \label{ste:ee:eq:terminal}
\end{align}
Define the integer polynomial
\begin{equation}
 \mathcal Q_T=\sum_{t=1}^T s_t^2+
 \sum_{t=1}^T\sum_{a,b=1}^n r_{t,ab}^2+
 \sum_{a=1}^n\sum_{b=2}^n q_{ab}^2.
 \label{ste:ee:eq:fullquartic}
\end{equation}

\begin{theorem}[Word-bijective full certificate]\label{ste:ee:thm:fullcertificate}
The polynomial $\mathcal Q_T$ has total degree at most four in its natural
witnesses. Its natural zero set is in bijection with the labelled erasing words
of length $T$. It has exactly
\begin{equation}
 T(n^2+k)\quad\text{witness variables},\qquad
 T(n^2+1)+n(n-1)\quad\text{squared residuals}.
 \label{ste:ee:eq:fullcounts}
\end{equation}
\end{theorem}
\begin{proof}
A sum of integer squares vanishes exactly when every residual vanishes. Natural
selectors with sum one have exactly one entry equal to one and every other entry
zero. No separate Boolean equation is needed. Let $i_t$ be that unique selected
letter. The prefix equations then force, successively,
\[
 X_t=B_{i_t}\cdots B_{i_1}.
\]
All these entries are natural, and each column sum equals $D^t$. Thus the
terminal equations say precisely that the stochastic matrix $X_T/D^T$ has equal
columns, which is equivalent to erasure by Lemma~\ref{ste:ee:lem:erasure}.

Conversely, each erasing word supplies those one-hot selectors and exact prefix
products, and all residuals vanish. The selectors and the recurrence determine
every witness uniquely. Different labelled words have different selector tuples,
even when some numerical generators coincide. This proves the bijection.
The prefix residuals have degree at most two because $B_i$ are fixed numerical
coefficients, and all other residuals are linear. The displayed counts simply
count selectors, matrix entries, and equations.
\end{proof}

The terminal test uses equal columns rather than $2\times2$ minors. It is sound
because the common, nonzero column sum has already been established by the
retained recurrence. If the target is specifically $J_n$, replace
\eqref{ste:ee:eq:terminal} by
\[
 (X_T)_{ab}-(X_T)_{11}=0\qquad ((a,b)\ne(1,1)).
\]
This uses $n^2-1$ terminal residuals. It is unnecessary in the doubly stochastic
case, but it distinguishes a general rank-one product from the uniform matrix
in the line-free construction.

\subsection{A smaller certificate in information coordinates}

Mass conservation allows a substantial further reduction. Put $h=n-1$ and
compute the integer restriction matrices
\begin{equation}
 (H_i)_{ac}=(B_i)_{ac}-(B_i)_{an},\qquad 1\leq a,c\leq h.
 \label{ste:ee:eq:integerH}
\end{equation}
These may have negative entries. They are the coordinates of the action of $B_i$
on the zero-mass hyperplane: $B_iE=EH_i$. For a prefix product $X_t$,
\begin{equation}
 (H_{i_t}\cdots H_{i_1})_{ab}=(X_t)_{ab}-(X_t)_{an}.
 \label{ste:ee:eq:restrictiondifference}
\end{equation}
Each entry on the right lies between $-D^t$ and $D^t$, since both nonnegative
entries belong to columns of sum $D^t$.

Introduce natural matrices $Y_t\in\N^{h\times h}$ and use the fixed shifts
\[
 p_t=D^t\quad(0\leq t\leq T),\qquad
 Y_0=I_h+\one_h\one_h\transpose.
\]
The intended value is
\[
 (Y_t)_{ab}=(H_{i_t}\cdots H_{i_1})_{ab}+p_t.
\]
Retain the selector residuals $s_t$ and define
\begin{align}
 \rho_{t,ab}
   &=(Y_t)_{ab}-p_t-
     \sum_{i=1}^k e_{t,i}\sum_{c=1}^h(H_i)_{ac}
                      \bigl((Y_{t-1})_{cb}-p_{t-1}\bigr),\label{ste:ee:eq:compressedrec}\\
 \theta_{ab}&=(Y_T)_{ab}-p_T.\label{ste:ee:eq:compressedterminal}
\end{align}
Let
\begin{equation}
 \mathcal I_T=\sum_t s_t^2+
 \sum_{t,a,b}\rho_{t,ab}^2+\sum_{a,b}\theta_{ab}^2.
 \label{ste:ee:eq:compressedquartic}
\end{equation}

\begin{theorem}[Mass-compressed natural quartic]\label{ste:ee:thm:compressed}
The natural zeros of $\mathcal I_T$ are in bijection with the labelled erasing
words of length $T$. Every residual has degree at most two, and the polynomial
has degree at most four. Its exact counts are
\begin{equation}
 \boxed{\ T\bigl((n-1)^2+k\bigr)\ \text{witnesses},\qquad
 T\bigl((n-1)^2+1\bigr)+(n-1)^2\ \text{residuals}.\ }
 \label{ste:ee:eq:compressedcounts}
\end{equation}
There is a canonical bijection between the natural zero sets of
$\mathcal Q_T$ and $\mathcal I_T$, preserving the selected labelled word.
\end{theorem}
\begin{proof}
The selectors again determine a unique word. Write
$Z_t=Y_t-p_t\one_h\one_h\transpose$. Equation~\eqref{ste:ee:eq:compressedrec} becomes
$Z_t=H_{i_t}Z_{t-1}$, with $Z_0=I_h$. Hence every $Y_t$ is forced uniquely.
The intended values are natural because \eqref{ste:ee:eq:restrictiondifference} bounds
each entry of $Z_t$ below by $-D^t$. Thus no additional range witnesses or
nonnegativity slack equations are needed.

The terminal residuals are zero exactly when $Z_T=0$, which means that the full
product vanishes on the zero-mass hyperplane. Lemma~\ref{ste:ee:lem:erasure} identifies
this with rank-one product existence. This proves both directions and the
word-bijection. The counts and degree bounds follow directly from the formulas.

For the zero-set bijection, keep the selectors and map a full prefix certificate
to the compressed one by
\[
 (Y_t)_{ab}=(X_t)_{ab}-(X_t)_{an}+D^t.
\]
On valid zeros these values are natural. Conversely, the selectors in a
compressed zero uniquely reconstruct the full prefix products. Both directions
preserve the word and are inverse on the zero sets.
\end{proof}

This is a bijection on the valid natural zero sets, not an identity between the
two polynomials away from zero. The map from arbitrary full witness assignments
need not preserve nonnegativity; only its zero-set behavior is asserted.
The distinction matches the source-audit discipline emphasized in ProveIt's
October 2 review \cite{repoReview}.

For the seven-state, eight-letter construction, the full certificate uses
$57T$ witnesses and $50T+42$ residuals. The compressed certificate uses
\begin{equation}
 \boxed{44T\ \text{natural witnesses and }37T+36\ \text{residuals}.}
 \label{ste:ee:eq:sevenledger}
\end{equation}
The improvement is $13T$ witnesses and $13T+6$ residuals. No assertion about an
optimal arithmetic circuit or a smallest possible polynomial follows from
these ledgers.

The constants $p_t=D^t$ are computed for the supplied horizon. Their bit lengths
are $O(t\log D)$, and printing all of them has a horizon-dependent cost. They
are not an uncharged fixed-arity exponentiation primitive. The next subsection
makes the uniform-input version polynomial without treating those powers as
free input constants.

\subsection{When the numerical matrices become polynomial parameters}

A degree statement must specify whether matrix entries are coefficients or
variables. In \eqref{ste:ee:eq:prefix}, treating $(B_i)_{ac}$ as polynomial input
parameters turns a term $e_{t,i}(B_i)_{ac}(X_{t-1})_{cb}$ into a cubic. It would
therefore be incorrect to reuse the full-prefix quartic degree claim unchanged
in that setting.

Here is a fully uniform compressed alternative for fixed $n,k,T$. Let the
positive entries of $B_i$ and the common column sum $D$ be input parameters,
subject to the stated stochastic-numerator promise. Introduce natural variables
$z_t$, $W_t\in\N^{h\times h}$, and $Y_t\in\N^{h\times h}$ at each step,
in addition to the selectors. Set $z_0=1$ and the same constant $Y_0$ as before.
Use the residuals
\begin{align}
 z_t-Dz_{t-1}&=0,\label{ste:ee:eq:clock}\\
 (W_t)_{ac}-\sum_i e_{t,i}\bigl((B_i)_{ac}+D-(B_i)_{an}\bigr)&=0,
 \label{ste:ee:eq:selectH}\\
 (Y_t)_{ab}-z_t-
 \sum_c\bigl((W_t)_{ac}-D\bigr)
       \bigl((Y_{t-1})_{cb}-z_{t-1}\bigr)&=0,
 \label{ste:ee:eq:uniforminfo}
\end{align}
along with the selector equations and terminal equations
$(Y_T)_{ab}-z_T=0$. Square and sum all these residuals.

\begin{proposition}[Uniform-input quartic]\label{ste:ee:prop:uniformquartic}
On the positive, common-column-sum input domain, the resulting polynomial has
total degree at most four, \emph{including its matrix and denominator input
parameters}. It has exactly
\begin{equation}
 T\bigl(k+2(n-1)^2+1\bigr)\quad\text{natural witnesses}
 \label{ste:ee:eq:uniformvar}
\end{equation}
and
\begin{equation}
 T\bigl(2(n-1)^2+2\bigr)+(n-1)^2\quad\text{squared residuals}.
 \label{ste:ee:eq:uniformres}
\end{equation}
Its natural zeros are again in bijection with the labelled erasing words of
length $T$.
\end{proposition}
\begin{proof}
Every displayed residual is quadratic or linear in all inputs and witnesses.
One-hot selectors force $z_t=D^t$ and
\[
 (W_t)_{ac}=(H_{i_t})_{ac}+D.
\]
Since $0\leq(B_i)_{ab}\leq D$, these selected values are natural. The $Y$
recurrence then reduces exactly to \eqref{ste:ee:eq:compressedrec}, so the preceding
proof supplies uniqueness and naturalness of the remaining witnesses.
There are $k+2h^2+1$ variables and $2h^2+2$ residuals per step, and $h^2$ terminal
residuals, yielding the counts.
\end{proof}

The input promises themselves are decidable finite equalities and inequalities.
They are not counted as witnesses in \eqref{ste:ee:eq:uniformvar}. To make a polynomial
whose accepted inputs also enforce these promises, include their finite
arithmetic conditions in the input-domain representation; one should then
charge those additional conditions separately. This proposition supplies the
uniform schema and its proof. The shipped numerical certificate evaluators
implement Theorems~\ref{ste:ee:thm:fullcertificate} and~\ref{ste:ee:thm:compressed}, not an
expanded symbolic version of Proposition~\ref{ste:ee:prop:uniformquartic}.

\subsection{Exact counts of solutions and limits of uniqueness}

The two numerical quartics have at most $k^T$ natural zeros, exactly one per
erasing labelled word. Combining the bijection with \eqref{ste:ee:eq:probcount} gives
\begin{equation}
 \PP(\tau\leq T)=\frac{\#\{x\in\N^{T((n-1)^2+k)}:\mathcal I_T(x)=0\}}{k^T}.
 \label{ste:ee:eq:quarticprobability}
\end{equation}
This identity is about a finite solution set for one externally fixed horizon.
It is not a finite-fold Diophantine representation of arbitrary computably
enumerable sets. A different polynomial is constructed for every $T$, and an
unbounded union over those polynomials is not automatically one fixed-arity
finite-fold polynomial.

\section{Fixed unbounded Diophantine representations and impossibility of bounds}
\label{ste:ee:sec:unbounded}

\subsection{What DPRM supplies, and what it does not}

Fix an effective encoding $e\in\N$ of finite rational stochastic alphabets.
Validity of the finite rational data, positivity, and common-column-sum claims
is decidable. The set of valid codes possessing an erasing word is computably
enumerable by exact word enumeration. The Davis--Putnam--Robinson--Matiyasevich
(DPRM) theorem identifies computably enumerable sets with Diophantine sets
\cite{matiyasevich}. Hence there is an integer polynomial $P$ and a fixed number
of natural witnesses such that
\begin{equation}
 e\text{ admits exact erasure}\quad\Longleftrightarrow\quad
 \exists x_1,\ldots,x_s\in\N:\ P(e,x_1,\ldots,x_s)=0.
 \label{ste:ee:eq:MRDP}
\end{equation}
The same conclusion holds for the compiled mortality or PCP input predicates.

\begin{proposition}[Fixed-arity quartic consequence]\label{ste:ee:prop:fixedquartic}
The unbounded erasure predicate has a fixed-arity integer-polynomial
representation of total degree at most four over natural witnesses.
\end{proposition}
\begin{proof}
Start from \eqref{ste:ee:eq:MRDP}. Split $P=P_+-P_-$, where $P_+$ and $P_-$ have
nonnegative integer coefficients. Evaluate these two polynomials by finite
addition/multiplication circuits over natural numbers. Introduce one natural
wire variable for each gate output and impose the appropriate equation
$z-u-v=0$, $z-uv=0$, or $z-c=0$. Identify the two final outputs by a linear
equation. Squaring and summing all gate and final-output equations gives total
degree at most four, including $e$ as a parameter. For fixed original witnesses,
all new circuit-wire values are unique and natural. Conversely, any zero of the
sum of squares reproduces the two evaluations and satisfies $P=0$.
The circuit and its arity are fixed because $P$ is fixed.
\end{proof}

This is an existence-and-compilation theorem, not a printed small universal
polynomial. It leaves the number and multiplicity of the original DPRM witnesses
uncontrolled. The explicit horizon-dependent polynomials of the previous section
are independent constructions with better traceability, not substitutes for
this fixed-arity argument. None of these observations resolves a general
finite-fold representation problem.

\wtag\ The circuit quadratization in this proof is the unique-extension quadratization \texttt{cdc:wf:lem:quadratization} of~\cite{cdc}, and the quartic lowering \texttt{pqc:lem:quarticlower} of~\cite{pqc}; it is printed as the source's own route, without a novelty claim. Part~\ref{ste:part:ec} proves the same corollary for its own predicate (Section~\ref{ste:ec:sec:mrdp}).

\subsection{No uniform computable exactness horizon}

Let the size of an input be the bit length of its finite binary encoding. The
following statement applies to either the fixed seven-state undecidable family
or the nine-state many-one complete family.
(\wtag\ ``The fixed seven-state undecidable family'' means the seven-state, at-most-eight-letter promised class of Corollary~\ref{ste:ee:cor:seven}: a class of input-dependent alphabets of fixed dimension, not one fixed alphabet, as the paragraph after that corollary explains.)

\begin{theorem}[No computable yes-instance horizon]\label{ste:ee:thm:nobound}
There is no total computable function $b:\N\to\N$ such that every erasing input
of bit length $N$ has an erasing word of length at most $b(N)$.
The conclusion remains true at every fixed rational contraction factor
$0<\eta<1$ used above.
\end{theorem}
\begin{proof}
Such a bound would decide erasure: compute $b(N)$ and enumerate all words up to
that length, multiplying exactly and checking equal columns. If a word erases,
answer yes; otherwise the assumed bound permits a no answer. Applied to the
compiler output, this would decide source mortality or PCP, contradicting the
proved reduction and its external undecidability dependency.
\end{proof}

By contrast, approximate forgetting has the uniform computable horizon
$\eta^T\leq\varepsilon$, independently of the matrix coefficients. The
undecidability concerns exact zero information, not the task of reaching any
fixed positive tolerance.

\wtag\ The proof of Theorem~\ref{ste:ee:thm:nobound} is the bounded-search argument of \texttt{cdc:bd:prop:nobound} in~\cite{cdc} and \texttt{pqc:qm:cor:nobound} in~\cite{pqc}, applied to exact erasure.

\begin{corollary}[No computable uniform finite-expectation bound]\label{ste:ee:cor:expectbound}
Under uniform independent switching, there is no total computable bound in the
input bit length for $\EE\tau$ on all inputs where that expectation is finite.
Nor is there a computable function $b(N)$ guaranteeing
$\PP(\tau\leq b(N))\geq1/2$ on every such input.
\end{corollary}
\begin{proof}
Let $L_{\min}$ be the shortest erasing-word length. Every finite value of $\tau$
is at least $L_{\min}$, so a finite upper bound for $\EE\tau$ bounds
$L_{\min}$. A bound with positive probability of erasure by a given horizon also
implies existence of a word by that horizon. Either would contradict
Theorem~\ref{ste:ee:thm:nobound} after taking an integer ceiling when necessary.
\end{proof}

For any fixed arity representation such as \eqref{ste:ee:eq:MRDP}, an analogous argument
excludes a computable bound on all necessary natural witness coordinates in
terms of the input size. A complete computably enumerable family of
\emph{non-erasure} certificates is also impossible on the undecidable class:
parallel enumeration of yes and no certificates would be a decision algorithm.
These are uniform impossibility statements, not obstacles to finding useful
certificates or tight bounds for particular instances.

\section{Exact erasure is not robust under perturbation}
\label{ste:ee:sec:perturbation}

The preceding hardness holds arbitrarily close to the uniform matrix. That is a
statement about the location of the input family, not robustness of individual
yes instances. Indeed, erasure is destroyed by arbitrarily small perturbations
that preserve the invariant-subspace restrictions.

\begin{theorem}[Subspace-preserving non-erasing perturbations]\label{ste:ee:thm:perturb}
Let $S_1,\ldots,S_k$ be positive rational column-stochastic matrices of dimension
$n\geq2$. For arbitrarily small positive rational $\delta$, all matrices
\begin{equation}
 S_i^{(\delta)}=(1-\delta)S_i+\delta I_n
 \label{ste:ee:eq:perturb}
\end{equation}
are invertible. They remain positive and column-stochastic, have exactly the
same common invariant complex subspaces as the original family, and have no
erasing word. If the original family contracts total variation by at most
$\eta$, the perturbed family contracts it by at most
$\eta_\delta=(1-\delta)\eta+\delta<1$ for $0<\delta<1$.
\end{theorem}
\begin{proof}
For each $i$, $f_i(\delta)=\det((1-\delta)S_i+\delta I_n)$ is a polynomial and
$f_i(1)=1$. It therefore has only finitely many roots. Outside the finite union
of these root sets, all perturbed generators are invertible. Every interval
$(0,\varepsilon)$ contains a positive rational outside that finite set.
A finite product of invertible matrices is invertible and, since $n\geq2$,
cannot have rank one.

Positivity and column-stochasticity follow immediately from the convex
combination for $0<\delta<1$. For every subspace $U$, invariance under $S_i$
implies invariance under $S_i^{(\delta)}$. Conversely,
\[
 S_i=\frac{S_i^{(\delta)}-\delta I_n}{1-\delta}
\]
shows the reverse implication whenever $\delta\ne1$. Applying this to every
generator proves equality of the common invariant-subspace families.
Finally, for $z=p-q$,
\[
 \|S_i^{(\delta)}z\|_1
 \leq(1-\delta)\eta\|z\|_1+\delta\|z\|_1,
\]
which proves the contraction claim.
\end{proof}

For a prescribed open entrywise band, first choose a compiled family strictly
inside that band. Its positive distance from the finitely many band endpoints
allows sufficiently small perturbations to remain inside it. Thus even within
the line-free, strongly contracting promised class, non-erasing instances occur
arbitrarily close to every constructed erasing instance.

This theorem also explains why floating-point rank tests cannot certify exact
erasure. A determinant or singular value that is numerically tiny can be nonzero
for a nearby rational input with a completely different exact answer. The shipped
checks therefore use integer arithmetic and exact rational Gaussian elimination,
with explicit rejection of floating-point and Boolean values in integer witness
fields.

\section{Worked examples and executable verification}
\label{ste:ee:sec:examples}

\subsection{A completely printed three-state example}

Take the scalar source alphabet $A_0=[0]$, tensor degree two, and $\eta=1/2$.
The compiler gives $K=13$, $D=39$, and
\begin{equation}
 B_0=\begin{pmatrix}13&13&13\\13&13&13\\13&13&13\end{pmatrix},\quad
 B_R=\begin{pmatrix}14&14&11\\12&15&12\\13&10&16\end{pmatrix},\quad
 B_L=\begin{pmatrix}18&15&15\\14&14&11\\7&10&13\end{pmatrix}.
 \label{ste:ee:eq:tiny}
\end{equation}
Every column sums to 39, every entry divided by 39 lies between $1/6$ and $1/2$,
and the family has no common invariant complex line by the guard proof.
The source generator is already $S_0=J_3$; this is a readable sanity-check
instance, not a universal numerical machine.

The chronological word $(0,L)$ has numerator product
\[
 B_LB_0=
 \begin{pmatrix}624&624&624\\507&507&507\\390&390&390\end{pmatrix}
\]
and denominator $39^2=1521$. Thus its stochastic product is
\[
 \begin{pmatrix}16/39\\13/39\\10/39\end{pmatrix}\one_3\transpose,
\]
which is rank one but not $J_3$. This explicitly demonstrates why the line-free
construction requires a distinction between erasure and a specified uniform
output, even though existence of the two kinds of word is equivalent.

The full horizon-two certificate has 24 natural witnesses and 26 residuals.
Its compressed version has 14 witnesses and 14 residuals. For the compressed
certificate, the restriction matrices are
\[
 H_0=0,\qquad H_R=3R,\qquad H_L=3L.
\]
After the source letter, the information restriction is zero and remains zero.
Hence the two $Y$ matrices are respectively $39\one_2\one_2\transpose$ and
$1521\one_2\one_2\transpose$. These are the exact values in the delivered JSON
certificate, not floating-point approximations.

Under uniform random switching, erasure occurs precisely when letter $0$ first
appears. Therefore
\[
 N_T=3^T-2^T,\qquad \PP(\tau\leq T)=1-(2/3)^T,\qquad \EE\tau=3.
\]
The enumerated values for $T=1,\ldots,5$ are $1,5,19,65,211$, matching the
certificate counts. This instance also permits direct checks of the zero--one
law and the difference between a uniform erasing product and an arbitrary
rank-one product.

\subsection{A seven-state, eight-letter numerical packet}

The delivered seven-state instance has source dimension three and six labelled
source generators:
\[
 N=\begin{pmatrix}0&1&0\\0&0&1\\0&0&0\end{pmatrix},\quad
 I_3,\quad -I_3,\quad \diag(2,1,1),\quad I_3+E_{12},\quad I_3+E_{23}.
\]
The two tensor guards give eight target letters. With $\eta=1/10$, the compiler
selects $K=131$ and $D=917$. All 392 integer entries of the eight numerator
matrices are printed in \path{examples/seven_state_instance.json}. (\wtag\ Shipped as \path{data/02-exact-erasure-seven_state_instance.json}; every \path{examples/*.json} file of source~02 is shipped as \path{data/02-exact-erasure-*.json}.)

The chronological word $(N,N,N)$ erases because $N^3=0$, while neither $N$ nor
$N^2$ is zero. The numerator product is
\[
 B_N^3=110156459\,\one_7\one_7\transpose,
 \qquad D^3=771095213=7\cdot110156459.
\]
The full certificate has 171 witnesses and 192 residuals. The compressed
certificate has 132 witnesses and 147 residuals. Both independent evaluations
return polynomial value exactly zero. A separate uniform-target full certificate
checks that this source-only word produces precisely $J_7$.

This finite example illustrates the format and counts of the undecidable class.
It does not prove undecidability by enumeration; that conclusion rests on the
all-input reduction and the cited mortality theorem.

\subsection{End-to-end PCP example}

Use the two tiles
\[
 (u_0,v_0)=(a,ab),\qquad (u_1,v_1)=(ba,a).
\]
The index word $(0,1)$ is a match because both concatenations are $aba$.
The four-dimensional front end has three generators, including connector $q=2$.
The chronological word $(q,0,1,q)$ is zero. After the two-guard lift, this becomes
an erasing word in a nine-state, five-letter positive stochastic instance.

Both that numerical instance and its certificates are included. The full
certificate has 344 natural witnesses and 400 residuals; the compressed one has
276 natural witnesses and 324 residuals. Every residual evaluates to zero in
both. The example exercises the complete chain
\[
 \text{PCP match}\ \longrightarrow\ \text{integer zero product}\
 \longrightarrow\ \text{positive stochastic erasure}\
 \longrightarrow\ \text{natural quartic zero}.
\]
The general proof also handles arbitrary connector placement, adjacent
connectors, and empty component strings, rather than checking only this
successful shape.

\subsection{Scope and results of the executable checks}

The shipped suite is deterministic. It uses Python's exact integers and
\texttt{fractions.Fraction}; no external Python package is required. Rational
matrix rank is computed by exact Gaussian elimination. The finite random source
sample uses seed \texttt{20261002}, recorded in the result file.

The completed run performed \textbf{50,582 passing assertions}. Its principal
categories are shown in Table~\ref{ste:ee:tab:tests}; the full category ledger is in
\path{verification/check_results.json}. (\wtag\ Shipped as \path{data/02-exact-erasure-check_results.json}; the suite was rerun at the write on a copy with the delivered layout and reproduced it apart from the recorded Python version.) Counts denote assertions, not distinct
mathematical instances or independently proved theorems.

\begin{table}[ht]
\centering\small
\begin{tabular}{@{}lr@{}}
\toprule
Check category & Assertions\\
\midrule
Mortality versus erasure for enumerated words & 10,188\\
Exact tensor product-rank formula & 10,188\\
Full word coordinate/block identity & 10,188\\
Independent full-prefix quartic replay & 3,448\\
Independent compressed quartic replay & 3,448\\
Full and compressed size ledgers, combined & 6,896\\
Exact total-variation contraction inequalities & 3,601\\
PCP arbitrary connector-word reduction & 1,638\\
All remaining algebraic, domain, mutation, and example checks & 987\\
\midrule
Total & 50,582\\
\bottomrule
\end{tabular}
\caption{Finite exact checks run for this report. The two numerical certificate
formats are checked by a module that imports no generator code.}
\label{ste:ee:tab:tests}
\end{table}

The suite includes scalar source alphabets with coefficients from $-2$ to $2$,
eight seeded two-dimensional source examples, the seven-state packet, a higher
tensor example, one-generator nilpotent examples attaining the commuting bound,
and PCP examples with empty words and adjacent connectors. It also changes
individual witness entries, changes selectors, and rejects values of type
\texttt{float}, \texttt{bool}, and \texttt{str} in natural-number fields.

The independent checker evaluates the defining residuals directly, without
calling the generator. It is independent at the implementation-module level,
not an independent human review. The all-instance reduction, the nonexistence
of common complex lines, the infinite-horizon probability theorem, and the
DPRM dependency are established mathematically in the text, not by these finite
experiments. The uniform-parameter schema has a hand proof but is not one of the
implemented certificate formats.

\section{Integration and formalization plan for ProveIt}
\label{ste:ee:sec:formalization}

The present package is self-contained and does not modify the repository. It is
suited to a separate research-report directory, with a link from the substrate
review rather than a change to the fixed universal-operation record.
(\wtag\ This report is that directory. A link from the substrate review is for the owner of the research programme to add; the review of this manuscript~\cite{review1977} already exists.)

A useful formalization order follows the mathematical dependency graph. First
formalize the explicit maps $E,F$, the affine block decomposition, and the
rank-one/equal-column equivalence over $\Q$. Next establish the tensor product
identity and the two common-line obstructions over $\C$. The mortality
undecidability theorem can initially be an explicitly declared external
interface, rather than an unproved statement hidden inside an implementation.
The PCP front end offers a separate path to a fully concrete reduction chain.

The certificate layer is especially suitable for an early machine-checked
milestone. It needs only finite sums, natural one-hot selectors, a matrix-product
recurrence, and the equivalence between a zero sum of squares and vanishing of
each residual. The compressed form additionally needs the column-sum invariant,
the difference formula \eqref{ste:ee:eq:restrictiondifference}, and the offset bounds.
This proves an exact zero-set bijection, not merely soundness of a successful
example.

For source integration, four interfaces should remain visibly separate:
\begin{enumerate}
\item the computable mortality/PCP-to-stochastic compiler, including its common
scale and all numerical input coefficients;
\item the exact finite-word rank and invariant-subspace theorems;
\item the bounded polynomial family and its literal witness/residual counts;
\item the unbounded DPRM polynomial existence theorem and its unspecified
original witness multiplicities.
\end{enumerate}
No existing Lean or Rocq file in ProveIt was compiled during this task, and no
new theorem here is presented as proof-assistant verified. The accompanying
\path{PROOF_STATUS.md} states these boundaries explicitly. (\wtag\ Shipped as \path{02-exact-erasure-PROOF_STATUS.md}, beside \path{02-exact-erasure-SOURCE_AUDIT.md}.)

\section{Further research questions and concrete next targets}
\label{ste:ee:sec:questions}

The following are research directions exposed by the present proofs. Except for
the precisely cited MathOverflow restriction already answered, this list does
not claim that every formulation is a previously catalogued open problem.

\subsection{Reduce the tensor overhead without losing the line guarantee}\label{ste:ee:q:overhead}

The present no-line compiler uses dimension $2d+1$ and two guard letters. Can one
construct a uniformly valid compiler in dimension $d+2$, or with one additional
letter, while retaining positivity and a word-projection soundness theorem?
The obstacle is not positivity, which the affine lift handles uniformly; it is
removing all source invariant lines without allowing guards to create false
zero products. A useful first goal is to characterize the affine guard systems
whose linear products cannot lower the rank of a protected source word.

\subsection{Irreducibility beyond a missing invariant line}\label{ste:ee:q:irreducible}

Our stochastic models necessarily have a common invariant hyperplane. Thus no
construction confined to normalized column-stochastic matrices can make the
whole linear action irreducible. A meaningful stronger target is either
irreducibility on the information hyperplane, or a positive integer matrix
construction without a common invariant hyperplane after normalization is
abandoned. The first target is not achieved here either: source invariant
subspaces can survive as $W\otimes\C^r$. A reduction must explain how to remove
those spaces while still decoding mortality.

\subsection{Low-dimensional classifications via the exact correspondence}\label{ste:ee:q:lowdim}

Proposition~\ref{ste:ee:prop:dimension} translates state-dimension questions directly
into mortality questions. It can be used to transport new decidability results,
bounded-alphabet results, or structural restrictions in either direction. The
three-state case is especially focused: its transverse dimension is two. Any
claimed classification should consult current primary sources rather than
inherit the status labels of a 2014 survey table.

\subsection{Regularly constrained and non-independent scheduling}\label{ste:ee:q:schedules}

For a finite-state controller, only a regular language of schedules may be
allowed. The unrestricted mortality compiler does not by itself preserve such
constraints when guard letters are erased. One needs a product construction
that tracks both source computation and controller state. On the probabilistic
side, identify conditions under which a reachable erasing block recurs almost
surely. Full-support independent switching is sufficient here, but zero
probabilities, transient controller states, and forbidden blocks can invalidate
the zero--one argument.

\subsection{Sharper canonical certificates and fully charged arithmetic}\label{ste:ee:q:certificates}

The compressed quartic reduces natural witnesses by using the invariant mass
rather than storing it. Can additional structure of the tensor restrictions
reduce the $h^2T$ information entries, while retaining a bijection for every
labelled word? One possibility is to certify the source projection and guard
product separately; another is to exploit the rank gap without recording a full
matrix at every step. Every proposal should count input parsing, selectors,
scale arithmetic, coefficient generation, and reconstruction obligations. An
improvement to $44T$ witnesses is a bounded-compiler result, not automatically
an improvement to the repository's 87-operation universal construction.

\subsection{Fixed-alphabet universality with an ordinary input interface}\label{ste:ee:q:fixedalphabet}

The nine-state many-one complete family encodes the problem in the gate alphabet.
A different objective is one fixed numerical stochastic alphabet with a separate
ordinary integer or finite-word input representation. Such a construction needs
a concrete universal correspondence or matrix system and a fully paid input
loader. It would make the relationship to the fixed universal Diophantine
polynomials in ProveIt substantially tighter.

\subsection{Structural non-erasure certificates}\label{ste:ee:q:nonerasure}

The commuting case supplies a complete finite certificate: a nonzero invariant
image on which all generators are invertible. Can useful larger decidable
classes be characterized by analogous invariant-space certificates? Candidate
restrictions include simultaneously triangularizable transverse actions,
prescribed rank patterns, or semigroups with a bounded descending-image chain.
A proposed algorithm must account for the possibility that noncommuting
invertible-on-current-image maps carry that image elsewhere; the commuting
proof cannot simply be reused without that invariant-image property.

\subsection{Precision-sensitive versus robust observation predicates}\label{ste:ee:q:robust}

The perturbation theorem shows that exact erasure is destroyed by generic small
invertibilizing perturbations, even when the invariant-subspace restrictions
are unchanged. A quantitative alternative is to study coefficient uncertainty
sets together with a fixed tolerance on output dependence. Which robust
predicates admit computable finite certificates, and how do their costs scale
as the tolerance tends to zero? The universal bound
$\eta^T\leq\varepsilon$ solves one natural approximate-forgetting question;
more delicate targets could depend on a specified output distribution, a
restricted schedule language, or nonuniform contraction rates.

\subsection{A formal proof of the guard compiler and the projection bijection}\label{ste:ee:q:formal}

A focused formal milestone is Theorem~\ref{ste:ee:thm:linefree} together with
Theorem~\ref{ste:ee:thm:compressed}, keeping mortality undecidability as a separately
identified dependency. A second milestone is the complete PCP chain and its
connector boundary condition. These targets are concrete enough to test the
repository's arithmetic and matrix APIs, and they would turn the present
hand-checked compiler into a reusable trusted frontend.

\section{Conclusion}

Exact stochastic erasure supports computational hardness even when every gate
is positive, every step uniformly contracts total variation, and no common
invariant complex line exists. The mechanism is explicit: preserve the source
computation on one tensor factor, use invertible guards on the other, and lift
the resulting affine dynamics into the interior of the stochastic simplex.

The strongest precise line-restriction result in this report is the
$(d,m)\mapsto(2d+1,m+2)$ compiler. It supplies the promised seven-state numerical
undecidability bound, while the independent four-dimensional PCP front end
supplies a nine-state many-one complete family. The probabilistic interpretation
then identifies exact finite forgetting, almost-sure finite forgetting, and
finite expected forgetting time without changing the underlying finite-word
predicate.

The Diophantine layer is equally explicit at a bounded horizon. A zero-set
bijection reduces the seven-state certificate from $57T$ to $44T$ natural
witnesses by storing only shifted information coordinates. The fixed-arity
unbounded representation remains a separate DPRM consequence. Together, these
results give a concrete new computational frontend and a carefully delimited
research contribution, rather than a claim that a generic invocation of DPRM
or contraction has solved every arithmetic-compression problem.


\part{Exact convergence as computation (source 10)}\label{ste:part:ec}

\noindent\wtag\ Part~II prints batch-78 manuscript~10, \emph{Exact Convergence as Computation: Diophantine Certificates for Uniformly Mixing Stochastic Games}, subtitled \emph{Fixed-game universality, canonical quadratic witnesses, and the distinction between scalar and vector reachability}, under the kicker ``Diophantine computational substrates'', ``Prepared for Vladimir Reshetnikov; Mathematical development and exact-arithmetic implementation: ChatGPT'', October 2, 2026. Its labels carry the prefix \texttt{ste:ec:}. Its research-status statement and abstract follow; its two appendices are Appendices~\ref{ste:ec:app:audit} and~\ref{ste:ec:app:contract}.

\begin{quote}\small
\textbf{Research status (source 10).} The central results below are supported by complete
mathematical arguments and a delivered, executable compiler. The accompanying
standard-library Python suite passes 26,077 exact-arithmetic checks. These
checks are not a proof-assistant verification, and the manuscript is not
peer reviewed. The literature and repository comparison is targeted rather
than exhaustive; priority for the precise constrained simulation theorem
is not established. No resolution of the general Skolem problem, no
single-fold MRDP theorem, and no improved fixed universal-polynomial operation
bound is claimed.
\end{quote}
\begin{quote}\small
\textbf{Abstract (source 10).} We construct a finite, turn-based, two-player stochastic game whose exact
value iteration is computationally universal despite unusually strong
regularity. The discount is $1/2$, every running reward is $1/4$, every
transition probability is dyadic, every available action assigns probability
at least $1/(2N)$ to each of the $N$ states, and the value vector is the known
constant $\tfrac12\one$. Nevertheless, for a suitable fixed game, deciding
whether an input dyadic terminal-payoff vector ever reaches that value vector
in finitely many Bellman iterations is computably enumerable complete.

The proof gives a compiler from counter programs to homogeneous linear/min/max
circuits and then to genuine stochastic-game transition tables. Paired
positive and negative channels remove signed coefficients; a correctly
initialized cyclic evaluation pipeline removes nested gates; and a common
uniform reset provides full support without destroying the simulated
computation. An erasing halt convention turns computational halting into
finite attainment of the fixed point. A zero-reward variant gives complete
coordinate-equality and finite-consensus problems.

For each externally specified horizon $T$, the complete computation admits
an orthant-nonnegative quadratic Diophantine certificate with exactly one
natural witness tuple when the prescribed endpoint condition holds. With
$N$ game states and $b$ binary-decision states, it uses $(N+2b)T$ witnesses.
Two boundary results clarify the source of hardness: deterministic
min/max-selector games have a finite abstraction, whereas chance-only
fixed-point attainment is a kernel-stabilization problem. Scalar equality
for chance-only games, in contrast, is Skolem-equivalent; an explicit dyadic,
full-support specialization is included and related to the existing ergodic
Markov-chain literature. Complete example polynomials, exact tests, and an
independent serialized-certificate checker accompany the article.
\end{quote}


\section{Research question, contribution, and source boundary}
\label{ste:ec:sec:intro}

A computational substrate need not resemble a tape machine. Here the
substrate is a familiar numerical recursion: take a vector, form several
convex averages of its coordinates, choose a minimum or maximum at each
coordinate, and repeat. All coefficients are rational. A discount makes the
recursion converge rapidly. The question is whether such benign-looking
iteration can still contain an unbounded exact computation.

The substantive question is not whether a computably enumerable relation
has some Diophantine representation. MRDP already answers that question.
Instead we ask for a reduction that respects a restrictive native interface:
nonnegative stochastic coefficients, at most two actions per state, common
positive transition support, a fixed discount, a trivial limiting value,
and a complete, witness-controlled arithmetic description of bounded runs.

\subsection{The central theorem}

Let a state be owned by a minimizing player, a maximizing player, or have
just one action. Write $p_{ia}\in\Q^N$ for the stochastic row of action $a$
at state $i$. Put
\begin{equation}
 S(v)_i=\opt_{a\in A_i}p_{ia}v,
 \qquad
 \Phi(v)=\frac14\one+\frac12S(v).
 \label{ste:ec:eq:main-operator}
\end{equation}
Here $\opt$ is the owner's minimum or maximum; for a one-action state there
is no choice. The vector $v$ is a terminal-payoff table, not a probability
distribution over the states.

\begin{theorem}[Fixed-game universality with a known fixed point]
\label{ste:ec:thm:main}
There are a fixed positive integer $N$ and a fixed finite game of the form
\eqref{ste:ec:eq:main-operator} such that:
\begin{enumerate}[label=(\roman*)]
\item each state has one or two actions; every $p_{ia,j}$ is dyadic and
      $p_{ia,j}\geq 1/(2N)$;
\item $\Phi$ maps $[0,1]^N$ into itself, is a $1/2$-contraction in the
      supremum norm, and has the unique fixed point $v_* = \tfrac12\one$;
\item the set of dyadic $v\in[0,1]^N$ for which
      $\Phi^n(v)=v_*$ for some $n\geq0$ is computably enumerable complete.
\end{enumerate}
The reduction is effective. More generally, the construction compiles every
finite counter program into such a game, with an explicitly computed
terminal-payoff encoding of each initial configuration.
\end{theorem}

The hard input is the initialization. The transition table, rewards,
discount, and target in the fixed-game assertion do not depend on that input.
A numerical transition table for a specifically chosen universal counter
machine is not printed or shipped. Fixed-game existence follows by applying
the fully specified compiler to a fixed universal counter machine. The
shipped numerical tables are worked examples, not universal-machine records.

\subsection{Why the source comparison matters}

The repository contains an MRDP interface stating the existence of one
ordinary integer polynomial and one finite witness dimension for an entire
computably enumerable set \cite{repo-mrdp}. Its computational-substrate
reports already discuss contracting dynamics and quadratic neural
certificates \cite{repoProb}. Consequently, neither the observation
that contraction can coexist with computation nor the generic use of
quadratic complementarity is claimed here as a discovery. The contribution
is the \emph{stochastic-game realization under the simultaneous restrictions
in Theorem~\ref{ste:ec:thm:main}}, together with its complete arithmetic interface.

\wtag\ \emph{Comparison with the collection.} Within this collection a 2-adic $1/2$-contraction already has an r.e.-complete exact extinction (Corollary \texttt{pqc:ad:cor:complete} in Part~VII of~\cite{pqc}, source~12 there), while a real four-dimensional piecewise-linear $1/2$-contraction has decidable point-to-point reachability because its radial clock is divided by $B$ at every step and never vanishes (\texttt{pqc:sc:thm:main}~(iii), source~14 there). The present construction is the real, order-preserving counterpart of the former: its live clock $s'=\sum_ee$ of \eqref{ste:ec:eq:live-clock} is erased at halting (Lemma~\ref{ste:ec:lem:erasure}), so the whole vector reaches the fixed point. It therefore shows that source~14's decidable point reachability comes from that compiler's non-erasable clock, not from real homogeneous contractions as such; nothing printed there is contradicted, and the picture agrees with the comparison section ``Two routes to one phenomenon'' of Part~VII of~\cite{pqc}. Part~V of~\cite{pqc} (source~17 there) has the same slogan for a different object: a uniform spectral gap with an undecidable exact stationary equality.

A particularly close external comparison is the January 26, 2026 revision
of Varonka and Watanabe \cite{varonka}. It proves point-reachability
results for one-player min- or max-Bellman operators, including arbitrary
dimension when source and target are componentwise comparable, and full
decidability in dimension two. Our construction uses both owners, a larger
fixed dimension, and initial errors of both signs. It therefore does not
settle their remaining one-player problem or contradict their positive
results. It does show why monotonicity and rapid fixed-point convergence
alone cannot justify unrestricted extensions to two-player operators.

The chance-only scalar-equality boundary is related to established work,
not a new undecidability assertion. Vahanwala proves Skolem and positivity
completeness for ergodic Markov chains with a dimension-efficient reduction
\cite{vahanwala}. We supply a less dimension-efficient specialization with
dyadic probabilities and a prescribed common reset. The general Skolem
problem remains unresolved in the July 2026 work of Luca, Ouaknine, and
Worrell, whose general decidability implication is conditional
\cite{skolem2026}. We do not resolve that problem.

\subsection{What is and is not computed}

There is no difficulty in obtaining the infinite-horizon game value in
\eqref{ste:ec:eq:main-operator}: every play receives the same running reward $1/4$
at every step, and
\[
 \sum_{t\geq0}2^{-t}\frac14=\frac12.
\]
Every strategy profile therefore has infinite-horizon payoff $1/2$.
The hard property is \emph{finite exact attainment from an input terminal
payoff table}. The usual initialization $v=\zero$ is not the input family
used in the reduction. In particular, this manuscript does not say that
solving ordinary finite discounted games is undecidable.

Value iteration at horizon $n$ is the finite-horizon game value with terminal
payoff $v$ and the indicated running reward. Its computation by backward
induction is immediate: after choosing an action, average the remaining
horizon's values, then optimize according to the current state's owner.
Thus the construction uses actual game semantics, not an arbitrary
piecewise-affine map merely renamed a Bellman operator. Stochastic games
and discounted values originate in Shapley's framework \cite{shapley};
approximation stopping criteria, such as those studied in
\cite{grobelna}, address a different question from finite exact attainment.

\section{Conventions and the three observation problems}
\label{ste:ec:sec:conventions}

We take $\N=\{0,1,2,\ldots\}$. A dyadic rational is an element of
$\Z[1/2]$. All finite rational objects have an effective exact encoding.
``Computably enumerable complete'' means complete under total computable
many-one reductions; no polynomial-time reduction bound is asserted.

We distinguish three predicates for an exact trajectory $v_n$:
\begin{align}
 \text{scalar equality: }&\quad (v_n)_r=(v_n)_s,
 \label{ste:ec:eq:scalar-observation}\\
 \text{consensus: }&\quad v_n\in\R\one,
 \label{ste:ec:eq:consensus-observation}\\
 \text{fixed-point attainment: }&\quad v_n=v_*.
 \label{ste:ec:eq:fixed-observation}
\end{align}
Each may be imposed at a specified $n=T$, or existentially over all $n$.
These are different decision problems. In particular, a scalar zero in a
linear recurrence is not the same as a whole vector becoming zero.

For a vector $x$, set
\[
 \|x\|_\infty=\max_i|x_i|,
 \qquad \spn(x)=\max_i x_i-\min_i x_i.
\]
The latter is a seminorm, since it vanishes on constant vectors. All
stochastic rows in this article act on column vectors of payoffs from the
left. The state of an exact computation is the entire payoff vector.
It is not the state visited by a single random sample path.

A \emph{bounded certificate} has an externally supplied natural horizon
$T$. Its arity is allowed to depend on $T$. A \emph{fixed-arity Diophantine
representation} uses one fixed polynomial and a fixed number of existential
variables for all inputs and all possible running times. Confusing these
notions would invalidate the final MRDP discussion.

\section{A homogeneous counter-machine compiler}
\label{ste:ec:sec:counter}

\subsection{Source instruction set}

A program has $m$ control states and $k\geq1$ natural counters. Its
instructions are
\[
 \operatorname{INC}(i,q),\qquad
 \operatorname{DECJZ}(i,q_+,q_0),\qquad
 \operatorname{HALT}.
\]
An increment increases counter $i$ and jumps to $q$. A decrement instruction
jumps to $q_0$ without changing counters when counter $i$ is zero; otherwise
it decreases that counter and jumps to $q_+$. There is one halt state $h$.
The ordinary operational convention makes it absorbing.

This instruction set supports universal computation with finitely many
counters, indeed with two. The distinction between counter instruction
sets matters: the present zero-test/decrement with independently specified
targets contains the undecidable CM2 instruction set described and analyzed
by Dudenhefner \cite{dudenhefner}. A fixed universal program can be obtained
from a universal interpreter in this standard counter model. Universality
of this source model is the sole external computability ingredient in the
hardness proof; the stochastic and arithmetic translations are proved here.

\subsection{Configuration coordinates and exact zero tests}

Encode a configuration $(q,n_1,\ldots,n_k)$ at scale $s>0$ as
\begin{equation}
 x=(s,p_1,\ldots,p_m,c_1,\ldots,c_k),\qquad
 p_j=s\,\mathbf1_{j=q},\qquad c_i=s2^{-n_i}.
 \label{ste:ec:eq:counter-encoding}
\end{equation}
The dimension is $d=1+m+k$. Every coordinate lies in $[0,s]$. The
homogeneous gate
\begin{equation}
 z_i=\max(0,2c_i-s)
 \label{ste:ec:eq:zero-test}
\end{equation}
returns $s$ when $n_i=0$ and $0$ otherwise. Indeed, a positive counter has
$c_i\leq s/2$. No discontinuous test has been used.

For a nonbranching instruction at state $j$, its enable value is $p_j$.
For a decrement/zero instruction testing counter $i$, its two enable values
are
\begin{equation}
 e_{j,0}=\min(p_j,z_i),\qquad
 e_{j,+}=\min(p_j,s-z_i).
 \label{ste:ec:eq:enable}
\end{equation}
On encoded configurations each enable value is $0$ or $s$, and precisely
one is $s$ when the active control state is included among the branches.

Associate to each branch $e$ its target state $q(e)$ and multipliers
$\alpha_{e,i}\in\{1/2,1,2\}$. An increment of counter $i$ uses $1/2$, a
positive decrement uses $2$, and an unchanged counter uses $1$. Put
\begin{align}
 p'_j&=\sum_{e:q(e)=j}e,
 \label{ste:ec:eq:next-state}\\
 c'_i&=\sum_e\min(2e,\alpha_{e,i}c_i).
 \label{ste:ec:eq:next-counter}
\end{align}
An inactive summand vanishes. An active summand equals its required
counter value: in the decrement case $2c_i\leq s$, and otherwise the value
is also at most $s$. The cap $2e=2s$ is therefore harmless.

\subsection{Two halt conventions}

We use two related maps. The \emph{absorbing map} $F_{\rm abs}$ includes a
halt branch targeting $h$ with all multipliers $1$, and sets $s'=s$.
It exactly implements the ordinary absorbing machine.

The \emph{erasing map} $F_{\rm era}$ omits the halt branch and sets
\begin{equation}
 s'=\sum_e e.
 \label{ste:ec:eq:live-clock}
\end{equation}
On a nonhalting configuration it performs the same machine step at the
same scale. On a halting configuration it outputs the entire zero vector.
Both maps take $\zero$ to $\zero$.

\begin{lemma}[Counter simulation]
\label{ste:ec:lem:counter}
The two maps are continuous, positively homogeneous, dyadic circuits using
only linear, minimum, and maximum gates. The absorbing map sends
\eqref{ste:ec:eq:counter-encoding} to the encoding of its successor. If the source
machine first reaches its halt state after $H$ steps, then
$F_{\rm era}^{H+1}(x)=\zero$, whereas $F_{\rm era}^j(x)$ is a nonzero
encoded configuration for $0\leq j\leq H$. If the machine never halts,
all these iterates remain nonzero and their clock coordinate remains $s$.
\end{lemma}

\begin{proof}
Equation~\eqref{ste:ec:eq:zero-test} gives the stated zero test. Exactly one
branch is enabled at a nonhalting configuration, so
\eqref{ste:ec:eq:next-state} writes a one-hot successor state and
\eqref{ste:ec:eq:next-counter} applies exactly the required multiplier to each
counter coordinate. For the erasing map all branches have enable value
zero at the halt configuration, including every counter summand, and
\eqref{ste:ec:eq:live-clock} vanishes. Homogeneity and continuity hold gate by
gate; the constant zero is a homogeneous linear gate. Induction proves
all trajectory claims.
\end{proof}

Using unrestricted fan-in for linear sums, the displayed circuit can be
organized with depth at most seven: linear zero-test preparation, maximum,
branch preparation, positive-branch selection, enable scaling, gated
counter terms, and final sums. This is a bound for this particular source
construction, not for arbitrary circuits or for binary fan-in arithmetic.
The stochastic implementation permits a finite distribution with more
than two successor states, so such sums do not violate the two-action
restriction.

\section{From signed circuits to genuine stochastic games}
\label{ste:ec:sec:pipeline}

\subsection{Layering and normalization}

Consider any positively homogeneous dyadic circuit
$F:\R^d\to\R^d$ whose gates are finite linear combinations and binary
minima or maxima. Insert copy gates so every edge enters the immediately
following layer. Pad \emph{all outputs} to one common final layer $L$,
chosen at least as deep as every gate retained in the circuit. The latter
qualification handles unused gates and the initially halted edge case.
Inputs are layer zero.

Choose powers of two $C_\ell\geq1$ so that every linear gate
$g=\sum_j a_jg_j$ at layer $\ell$ satisfies
$\sum_j|a_j|\leq C_\ell$. Define
\begin{equation}
 K_0=1,\qquad K_\ell=C_1\cdots C_\ell,
 \qquad K=K_L,\qquad M=L+1.
 \label{ste:ec:eq:normalizers}
\end{equation}
For input $x$, every gate value at layer $\ell$ has absolute value at most
$K_\ell\|x\|_\infty$. This follows by induction for linear gates and by
$|\min(a,b)|,|\max(a,b)|\leq\max(|a|,|b|)$ for the other gates.

\subsection{Paired channels and a baseline state}

Create a baseline state $o$, with its only transition back to itself.
For every circuit input and gate $g$, create two states $g^+,g^-$. They
will hold a value and its negative. On a centered initialization the
baseline value is zero. A retained gate at layer $\ell$ is intended to
hold $g/K_\ell$ in its positive channel and $-g/K_\ell$ in its negative
channel.

For a linear gate $g=\sum_j a_jg_j$ at layer $\ell$, give $g^+$ a single
action with probabilities
\begin{equation}
 \frac{|a_j|}{C_\ell}
 \quad\hbox{to }g_j^+\text{ if }a_j\geq0,
 \quad\hbox{to }g_j^-\text{ if }a_j<0,
 \label{ste:ec:eq:linear-rail}
\end{equation}
and send the remaining probability to $o$. For $g^-$ swap the signs of the
selected channels. Repeated destinations are combined. Nonnegativity of
the remaining probability is exactly the normalization inequality.

For $g=\min(g_1,g_2)$, make $g^+$ a MIN state. Its two actions assign
probability $1/C_\ell$ to $g_1^+$ and $g_2^+$, respectively, and the
remaining probability to $o$. Make $g^-$ a MAX state with the analogous
negative-channel actions. For a maximum gate reverse the two owners.
Thus the negative channel really is the negative output, including at ties.
All rows are stochastic and dyadic; every state has at most two actions.

Finally, each input-register channel has a single action copying its
corresponding output channel at layer $L$. This creates a cyclic pipeline.
Denote its undiscounted, zero-reward operator by $S_0$.

If there are $G$ retained gates in addition to the $d$ inputs, the number
of states before padding is
\begin{equation}
 N_0=1+2(d+G).
 \label{ste:ec:eq:state-count}
\end{equation}
Pad to a power of two $N\geq N_0$, with every extra state copying $o$.
Padding will keep a subsequent uniform reset dyadic. If the circuit has
$g_{\min\max}$ min/max gates, then the game has exactly
$b=2g_{\min\max}$ binary-decision states. No optimality claim is attached
to these counts.

\subsection{The invariant subspace}

Let $V\subseteq\R^N$ consist of vectors whose baseline and padding
coordinates vanish and whose paired channels are negatives of each other.
The construction gives
\begin{equation}
 S_0(V)\subseteq V,\qquad \frac1N\sum_iw_i=0\quad(w\in V).
 \label{ste:ec:eq:dual-subspace}
\end{equation}
The first assertion follows separately for linear gates, min/max pairs,
copy registers, and the baseline. The second is cancellation of the pairs.
In particular, signed arithmetic is represented inside a genuine
order-preserving stochastic operator. The order on the full space does
not prohibit this: distinct generic vectors of $V$ are incomparable in
the componentwise order.

\subsection{Initialization is part of the compiler}

Given $x\in[-1,1]^d$, evaluate the finite circuit once on $x$. Initialize
its input channels to $\pm x_j$, each layer-$\ell$ gate to its paired
values $\pm g(x)/K_\ell$, and the baseline and padding coordinates to
zero. Call the resulting vector $w(x)\in V\cap[-1,1]^N$.

This is a finite, computable loader. It computes only the circuit's
one-step gate values. It does not inspect whether the source program will
halt. Zero-filling the internal gates instead would be unsound: the first
feedback step would copy zero outputs into all registers and could falsely
signal a halted or erased computation.

\begin{lemma}[Every microstep is accounted for]
\label{ste:ec:lem:pipeline}
Let $r_n\in\R^d$ be the positive input-register coordinates of
$S_0^n(w(x))$. Then
\begin{equation}
 r_n=K^{-q_n}F^{q_n}(x),\qquad
 q_n=\left\lceil\frac nM\right\rceil,
 \label{ste:ec:eq:microstep}
\end{equation}
where $q_0=0$. In particular,
$r_0=x$ and $r_1=\cdots=r_M=F(x)/K$.
\end{lemma}

\begin{proof}
Each layer computes its gates from the preceding layer with an additional
factor $1/C_\ell$. Following an input-register vector through all $L$
layers and then the feedback edge therefore gives, for every $n\geq0$,
\begin{equation}
 r_{n+M}=F(r_n)/K.
 \label{ste:ec:eq:feedback-recurrence}
\end{equation}
At time zero each layer was already evaluated consistently on $x$.
For $0\leq t\leq\ell$, the layer-$\ell$ values at time $t$ are still the
normalized gate values evaluated on $x$: prove this by induction on
$\ell$, using the preceding layer at time $t-1$ for $t>0$.
Thus the final layer has its initial output for times $0,\ldots,L$, and
the feedback registers equal $F(x)/K$ for times $1,\ldots,M$.
Positive homogeneity now iterates \eqref{ste:ec:eq:feedback-recurrence}.
For each residue class modulo $M$, its initialized member has exactly
the exponent specified in \eqref{ste:ec:eq:microstep}; every subsequent member
increases that exponent by one.
\end{proof}

\begin{lemma}[Erasure drains the whole pipeline]
\label{ste:ec:lem:erasure}
Use $F=F_{\rm era}$. If the source first reaches its halt state after
$H$ steps, then
\begin{equation}
 S_0^n(w(x))=\zero\quad\hbox{for all }n\geq(H+1)M.
 \label{ste:ec:eq:drain-bound}
\end{equation}
If the source never halts, $S_0^n(w(x))\neq\zero$ for every $n$.
\end{lemma}

\begin{proof}
In the halting case $F^{H+1}(x)=\zero$ and all its later iterates vanish.
Lemma~\ref{ste:ec:lem:pipeline} makes all registers zero from time $HM+1$
onward. After at most $L$ more steps all internal gates have received
only zero data. Every gate is homogeneous and hence maps zero inputs to
zero. The time bound is $HM+1+L=(H+1)M$.
In the nonhalting case the clock register is $K^{-q_n}s>0$ at every finite
time, so the whole vector cannot vanish.
\end{proof}

The extra $L$ steps cannot in general be omitted. Registers can have
become zero while stale, unused, or not-yet-overwritten internal gates
remain nonzero. The full-vector theorem uses the drain argument, not
just the register simulation.

\section{Uniform mixing preserves exact computation}
\label{ste:ec:sec:mixing}

\subsection{A common-reset lemma}

For any zero-reward stochastic min/max operator $S$, any rational
$0<\eta<1$, and $\bar v=N^{-1}\sum_jv_j$, define
\begin{equation}
 \widehat S(v)=(1-\eta)S(v)+\eta\bar v\one.
 \label{ste:ec:eq:reset}
\end{equation}
It is itself a stochastic-game operator: replace every action row by
\begin{equation}
 \widehat p_{ia}=(1-\eta)p_{ia}+\frac\eta N\one^{\mathsf T}.
 \label{ste:ec:eq:reset-rows}
\end{equation}
The common term may be moved outside either a minimum or a maximum.
Every new transition probability is at least $\eta/N$.

\begin{lemma}[Exact preservation of all equality times]
\label{ste:ec:lem:reset}
Let $B=\delta\widehat S$ with $0<\delta<1$. For every initial vector $v$
there are scalars $c_n$ such that
\begin{equation}
 B^n(v)=[\delta(1-\eta)]^nS^n(v)+c_n\one.
 \label{ste:ec:eq:reset-orbit}
\end{equation}
Consequently, every pairwise equality or strict order between coordinates
of $S^n(v)$ is preserved at exactly the same $n$ by $B^n(v)$.
\end{lemma}

\begin{proof}
Stochastic min/max operators are positively homogeneous and additively
homogeneous:
\[
 S(av+c\one)=aS(v)+c\one\qquad(a\geq0).
\]
Put $a_n=[\delta(1-\eta)]^n$. The assertion holds at $n=0$ with $c_0=0$.
Substitution into \eqref{ste:ec:eq:reset} gives the induction step with
\[
 c_{n+1}=\delta c_n+\delta\eta a_n\overline{S^n(v)}.
\]
Subtracting two coordinates removes $c_n$, and $a_n$ is strictly positive
at every finite $n$.
\end{proof}

This lemma does not assume that an input is a valid encoded computation.
It is an identity for every real vector. In particular, the reset is not
an approximation justified only on a finite test set.

\subsection{Contraction estimates}

\begin{lemma}[Supremum and span contractions]
\label{ste:ec:lem:contraction}
For any common reward $r$ and $\Phi=r\one+\delta\widehat S$,
\begin{align}
 \|\Phi(v)-\Phi(u)\|_\infty&\leq\delta\|v-u\|_\infty,
 \label{ste:ec:eq:sup-contract}\\
 \spn(\Phi(v)-\Phi(u))&\leq\delta(1-\eta)\spn(v-u).
 \label{ste:ec:eq:span-contract}
\end{align}
\end{lemma}

\begin{proof}
If $m\one\leq v-u\leq M\one$, monotonicity and additive homogeneity give
$m\one\leq S(v)-S(u)\leq M\one$. This proves its supremum nonexpansiveness
and bounds the span of the difference by $M-m$. The same supremum argument
applies to $\widehat S$, which is a stochastic operator. For the span,
the reset contribution $\eta(\bar v-\bar u)\one$ is constant and therefore
disappears. Multiplication by $\delta$ completes both inequalities.
\end{proof}

On the invariant space $V$ of \eqref{ste:ec:eq:dual-subspace}, the mean vanishes
at every iterate, so a stronger identity is available:
\begin{equation}
 (\delta\widehat S)^n(w)
       =[\delta(1-\eta)]^nS_0^n(w),\qquad w\in V.
 \label{ste:ec:eq:dual-reset}
\end{equation}
Now choose $\delta=\eta=1/2$, define
\begin{equation}
 v(x)=\frac{w(x)+\one}{2},\qquad
 \Phi(v)=\frac14\one+\frac12\widehat S(v),
 \label{ste:ec:eq:affine-shift}
\end{equation}
and use additive homogeneity. One obtains
\begin{equation}
 \boxed{\displaystyle
 \Phi^n(v(x))=\frac12\one+\frac{4^{-n}}2S_0^n(w(x)).}
 \label{ste:ec:eq:main-identity}
\end{equation}
Every initial coordinate belongs to $[0,1]$, all coefficients are dyadic,
and all reset rows have entries at least $1/(2N)$. On these encoded inputs,
the error is bounded by $4^{-n}/2$ in the supremum norm. For arbitrary
$v\in[0,1]^N$, the uniform bound is $2^{-n}\|v-\tfrac12\one\|_\infty$.

\section{Proof of fixed-game universality and its variants}
\label{ste:ec:sec:universality}

\begin{proof}[Proof of Theorem~\ref{ste:ec:thm:main}]
Fix a universal counter program and apply the erasing compiler.
Choose the power-of-two state padding and the reset and reward in
\eqref{ste:ec:eq:affine-shift}. All syntactic assertions about the resulting
fixed game follow from the explicit row construction.
Since a stochastic operator fixes constant vectors,
$\Phi(\tfrac12\one)=\tfrac12\one$. The supremum contraction proves
uniqueness. The cube is preserved because each convex average and its
minimum or maximum lies in $[0,1]$, whence every output of $\Phi$ lies in
$[1/4,3/4]$.

Encode an input of the fixed universal program by the dyadic vector
$v(x)$. Equations~\eqref{ste:ec:eq:main-identity} and
\eqref{ste:ec:eq:drain-bound} show
\[
 \exists n\;\Phi^n(v(x))=\tfrac12\one
 \quad\Longleftrightarrow\quad
 \exists n\;S_0^n(w(x))=\zero
 \quad\Longleftrightarrow\quad
 \text{the source computation halts}.
\]
All factors discarded in these equivalences are strictly positive at
finite times. The map from a source input to the full payoff vector is
computable. Thus the fixed game's attainment set is c.e.-hard. It is
c.e. because rational iteration and equality are decidable at every
finite time, so successive times can be searched. The same reasoning
applies uniformly when the source program is supplied to the compiler.
\end{proof}

\subsection{Uniformly stopping reachability games}

The operator $\Phi$ also has a reachability interpretation. From every
nonterminal state, every action goes to a success terminal with probability
$1/4$, to a failure terminal with probability $1/4$, and otherwise continues
according to its row $\widehat p_{ia}$ with total probability $1/2$.
There are now $N+2$ states. Both terminals are absorbing.

The continuation probability for $n$ steps is exactly $2^{-n}$ under every
strategy profile, and the expected number of transitions before stopping
is exactly two. Success probability is exactly $1/2$ from every
nonterminal state. With a supplied boundary payoff table on nonterminal
states, backward reachability iteration is exactly $\Phi$.

\begin{corollary}[Rapid stopping does not imply decidable exact attainment]\label{ste:ec:cor:stopping}
There is a fixed, uniformly stopping, finite turn-based reachability game
with success-probability vector $\tfrac12\one$ for which finite exact
attainment of that vector from an input dyadic boundary table is
c.e.-complete. The full-support lower bound applies to the conditional
transition kernel on the nonterminal states; it does not apply to the
absorbing terminal rows.
\end{corollary}

Universality here belongs to the arithmetic of the iterates. A typical
random path is short, and we do not assert that it physically executes
an arbitrarily long counter computation.

\subsection{Zero running rewards and scalar observations}

For the absorbing source map, let $B=\tfrac12\widehat S$ have zero
running rewards. The shifted encoding satisfies
\begin{equation}
 B^n(v(x))=\frac{2^{-n}}2\one+\frac{4^{-n}}2S_0^n(w(x)).
 \label{ste:ec:eq:zero-reward-shift}
\end{equation}
Let $r$ be the positive halt-state register, and $s$ the positive clock
register. For initial scale one, put $q_n=\lceil n/M\rceil$. Before
halting, their difference is
\begin{equation}
 (B^n(v(x)))_r-(B^n(v(x)))_s
       =-\frac{4^{-n}}{2K^{q_n}};
 \label{ste:ec:eq:exact-gap}
\end{equation}
after the simulated state is halted, the difference is zero.
If the source's first halting time is $H\geq1$, the first equality time is
\begin{equation}
 n_{\rm eq}=(H-1)M+1.
 \label{ste:ec:eq:equality-time}
\end{equation}
For an initially halted source it is zero. The separate observation
``halt-register value greater than baseline value'' is true exactly at
simulated halted configurations, giving c.e.-complete strict-inequality
reachability as well.

For the erasing source map, consensus in
\eqref{ste:ec:eq:zero-reward-shift} is equivalent to
$S_0^n(w(x))=\zero$: the latter vector belongs to $V$, and the only
constant vector in $V$ is zero. Thus finite consensus is c.e.-complete in
a suitable fixed zero-reward game, although its unique fixed point is zero
and all initial payoffs in the reduction are nonnegative.

\subsection{One-sided errors are easy}

\begin{proposition}\label{ste:ec:prop:onesided}
For the full-support game in Theorem~\ref{ste:ec:thm:main}, if
$v\leq v_*$ or $v\geq v_*$ componentwise, finite attainment occurs if and
only if $v=v_*$ already.
\end{proposition}

\begin{proof}
Suppose $v-v_*\leq0$ and some coordinate is strictly negative.
Every full-support action average of $v-v_*$ is then strictly negative,
so every coordinate of $\Phi(v)-v_*$ is strictly negative. This remains
true at every later iterate. The nonnegative case is identical with
signs reversed.
\end{proof}

The encoded hard initializations straddle the fixed point through their
paired positive and negative errors. This is essential, not a dispensable
technicality.

\section{Two decidable boundaries and one Skolem boundary}
\label{ste:ec:sec:boundaries}

The combination of averaging and mixed min/max control is important. Each
of two natural restrictions has a direct decision procedure for full
fixed-point attainment.

\subsection{No averaging: finite selector dynamics}

Suppose every action row is a point mass. Then $S$ chooses one of the
current coordinates at each state, possibly after a minimum or maximum.
If $v$ has $m_0\leq N$ distinct coordinate values, every vector $S^n(v)$
has all coordinates in this fixed finite set. There are at most $m_0^N$
such vectors.

\begin{proposition}[Deterministic-selector decidability]\label{ste:ec:prop:selectors}
Scalar equality, consensus, and attainment of the uniform fixed point
are decidable for zero-reward or common-reward discounted games with
point-mass action rows. At most $m_0^N$ distinct undiscounted vectors need
be visited before a repetition proves that no new observation will occur.
\end{proposition}

\begin{proof}
For scalar equality and consensus, positive discount factors and common
additive rewards do not affect the observations. Iterate the finite
selector map and test each vector until it repeats. For the uniform fixed
point $v_*=r/(1-\delta)\one$, translate to the error $v-v_*$.
The error at time $n$ is $\delta^n S^n(v-v_*)$, so fixed-point attainment
is zero-vector reachability in the same finite abstraction.
\end{proof}

\subsection{No choices: whole-vector zeros are easy}

Suppose each state has only one action, so $S(v)=Pv$ for one rational
stochastic matrix. For a common reward and its uniform fixed point,
\[
 \Phi^n(v)-v_* = \delta^nP^n(v-v_*).
\]

\begin{proposition}[Kernel-stabilization test]
\label{ste:ec:prop:linear-point}
Finite fixed-point attainment for a chance-only $N$-state game is
 equivalent to $P^N(v-v_*)=\zero$ and is therefore decidable by exact
linear algebra.
\end{proposition}

\begin{proof}
The kernels $\ker P^j$ form an increasing chain. Their dimensions can
strictly increase at most $N$ times. If two consecutive kernels agree,
all subsequent kernels agree, by taking inverse images under $P$.
Hence $\bigcup_{j\geq0}\ker P^j=\ker P^N$.
\end{proof}

This conclusion does \emph{not} solve the scalar zero problem. A single
coordinate may vanish while the whole orbit vector stays nonzero.

\subsection{Scalar equality: a dyadic, full-support Skolem specialization}

Consider the integer-matrix form of the Skolem problem:
\begin{equation}
 \text{given }A\in\Z^{d\times d},\ x\in\Z^d,
 \quad \exists n\geq0\; e_1^{\mathsf T}A^nx=0\;?
 \label{ste:ec:eq:skolem}
\end{equation}
Every integer linear recurrence is put into this form by its companion
matrix. Conversely, Cayley--Hamilton gives a recurrence for every such
matrix sequence.

Choose powers of two
\[
 C\geq\max\left(1,\max_i\sum_j|A_{ij}|\right),\qquad
 R\geq\max(1,\|x\|_\infty).
\]
Use a baseline and $d$ paired channels. For the positive channel $i^+$,
put mass $|A_{ij}|/C$ on $j^+$ or $j^-$ according to the sign of $A_{ij}$;
put the remaining mass on the baseline. Swap channels for $i^-$. Pad
to a power of two $N\geq2d+1$ and add a half-uniform reset.
Initialize centered channels to $\pm x_i/R$, then shift to $[0,1]$.
There are no min/max decisions.

For either zero rewards or a common reward, let $r=1^+$ and let $o$ be
the baseline. With discount $1/2$, the difference is
\begin{equation}
 v_{n,r}-v_{n,o}
 =\frac{e_1^{\mathsf T}A^nx}{2R(4C)^n}.
 \label{ste:ec:eq:skolem-lift}
\end{equation}
Every coefficient is dyadic and every transition has probability at least
$1/(2N)$. The formula follows by induction on the paired channels before
reset and by Lemma~\ref{ste:ec:lem:reset} after reset.

\begin{proposition}[Scalar-equality equivalence]
\label{ste:ec:prop:skolem}
For the class of chance-only dyadic games with discount $1/2$ and a
half-uniform reset, scalar-equality reachability is computably many-one
equivalent to the Skolem problem. The game is part of the input in this
statement; no fixed chance-only game is asserted c.e.-complete.
\end{proposition}

\begin{proof}
Equation~\eqref{ste:ec:eq:skolem-lift} gives one reduction. In the reverse
direction, cancel common rewards and the positive discount. Equality of
two coordinates is the zero of
$(e_r-e_s)^{\mathsf T}P^nv$. Clear the denominators of $P$ and $v$:
with $D P$ integral and $a v$ integral, the integer sequence
$(e_r-e_s)^{\mathsf T}(DP)^n(av)$ has exactly the same zeros.
Cayley--Hamilton supplies its integer recurrence effectively.
\end{proof}

This is an explicit specialization of an established hardness phenomenon.
It is not a claim that Skolem-hardness of ergodic chains was previously
unknown, nor an improvement over the smaller dimension achieved in
\cite{vahanwala}. The specialization is useful because its probability
and reset constraints match the nonlinear compiler exactly.

\begin{center}
\begin{tabular}{>{\raggedright\arraybackslash}p{0.31\textwidth}>{\raggedright\arraybackslash}p{0.27\textwidth}>{\raggedright\arraybackslash}p{0.29\textwidth}}
\toprule
Substrate & Scalar equality & Uniform fixed-point attainment\\
\midrule
Point-mass rows; MIN and MAX allowed & Decidable by finite selector states
 & Decidable by the same finite abstraction\\[3pt]
Chance only; rational stochastic rows & Skolem-equivalent, even with the
 stated dyadic reset & Decidable by kernel stabilization\\[3pt]
MIN, MAX, and stochastic rows & c.e.-complete for a suitable fixed game
 & c.e.-complete for a suitable fixed game\\
\bottomrule
\end{tabular}
\end{center}

The table is not a classification of every one-player subclass. In
particular, allowing stochastic rows and only MAX states, or only MIN
states, is a separate frontier.

\wtag\ \emph{Bears on a question of~\cite{pqc}.} Part~VII of~\cite{pqc} asks (Question ``Structural subclasses'') which restrictions force half-space reachability to be decidable: coordinatewise monotonicity, nonnegative weights, boundedly many activation patterns, or a genuine Bellman-operator structure. Theorem~\ref{ste:ec:thm:main} and Section~\ref{ste:ec:sec:universality} answer this negatively in part: the operators here are coordinatewise monotone, have nonnegative (stochastic) coefficients and a genuine two-player Bellman structure, and still the strict-inequality observation ``halt-register value greater than baseline value'', scalar equality and exact fixed-point attainment are c.e.-complete for suitable fixed games. The one-player case (only MIN or only MAX states) stays open (Section~\ref{ste:ec:q:oneplayer}).

\section{Canonical quadratic Diophantine certificates}
\label{ste:ec:sec:quadratic}

We now compile the native game recursion directly, rather than importing
an MRDP polynomial. Every existential variable below ranges over $\N$.
The nonnegative domain is essential.

\subsection{Integer scaling of zero-reward runs}

Let every transition probability be $P_{ia,j}/D$, with
$P_{ia,j}\in\N$ and $\sum_jP_{ia,j}=D$. Let the discount be $1/2$ and
$v_0=A/d_0$, where $A\in\N^N$ and $d_0\geq1$. Define
\begin{equation}
 X_t=d_0(2D)^t v_t.
 \label{ste:ec:eq:integer-scale-zero}
\end{equation}
For zero rewards the recurrence becomes
\begin{equation}
 X_{t+1,i}=\opt_{a\in A_i}L_{ia}(X_t),\qquad
 L_{ia}(X)=\sum_jP_{ia,j}X_j,
 \qquad X_0=A.
 \label{ste:ec:eq:integer-recurrence}
\end{equation}
In particular, every $X_t$ is integral and nonnegative. The exponentially
large denominator is not an uncharged existential exponentiation gadget:
$T$ is external, and \eqref{ste:ec:eq:integer-recurrence} itself has no
exponentiation in its transition constraints.

\subsection{A min/max gate with unique gaps}

For two nonnegative integers $L_0,L_1$, a claimed output $Z$, and natural
gaps $u,v$, define
\begin{align}
 q_{\min}(L_0,L_1,Z;u,v)
  &=(L_0-Z-u)^2+(L_1-Z-v)^2+uv,
 \label{ste:ec:eq:qmin}\\
 q_{\max}(L_0,L_1,Z;u,v)
  &=(Z-L_0-u)^2+(Z-L_1-v)^2+uv.
 \label{ste:ec:eq:qmax}
\end{align}

\begin{lemma}[Canonical complementarity]
\label{ste:ec:lem:complementarity}
On natural inputs, \eqref{ste:ec:eq:qmin} vanishes exactly when
$Z=\min(L_0,L_1)$ and $(u,v)=(L_0-Z,L_1-Z)$. Formula~\eqref{ste:ec:eq:qmax}
has the analogous statement for the maximum. These polynomials are
nonnegative on the entire real nonnegative orthant.
\end{lemma}

\begin{proof}
Every displayed summand is nonnegative on that orthant. A zero therefore
forces both linear residuals and $uv$ to vanish. In the minimum case,
$L_0=Z+u$ and $L_1=Z+v$, so $Z$ is a lower bound; since $uv=0$, it equals
at least one of the two values. The gaps are then forced. The maximum
case is identical with inequalities reversed. Conversely, the asserted
output and gaps make every term zero.
\end{proof}

When the two actions tie, both gaps are zero. There is no selector bit
whose value can be chosen in two ways. Also, $uv$ is \emph{not squared}.
Squaring it would unnecessarily raise the degree from two to four.
Its sign is already controlled by the domain. This use of complementarity
is consistent with the repository's prior quadratic neural-certificate
work \cite{repoProb}; the contribution here is its full game-level
specialization and resource ledger.
(\wtag\ The repository work meant is source~15 of Part~VII of~\cite{pqc}: Lemma \texttt{pqc:pm:lem:clamp}, an unsquared complementarity certificate on $\N$ of degree two, and Theorem \texttt{pqc:pm:thm:trace}, one degree-two polynomial with canonical scales and a unique natural zero; compare also \texttt{pqc:sc:lem:relucert}. Lemma~\ref{ste:ec:lem:complementarity} specializes that device to minimum and maximum gates; the game-level gate types and counts are this source's.)

For a one-action state use the single square $(Z-L_0)^2$.

\subsection{The complete polynomial}

For each step $t=0,\ldots,T-1$ introduce the next state vector $X_{t+1}$,
and two gaps for each binary-decision state. Sum the state terms
\eqref{ste:ec:eq:qmin}, \eqref{ste:ec:eq:qmax}, or the one-action square. Add an endpoint
sum $E_T$. The resulting \emph{one polynomial} is
\begin{equation}
 Q_T(A;W)=\sum_{t=0}^{T-1}\sum_{i=1}^N q_{t,i}+E_T.
 \label{ste:ec:eq:complete-quadratic}
\end{equation}
For scalar equality use $E_T=(X_{T,r}-X_{T,s})^2$. For consensus use
$E_T=\sum_{i=2}^N(X_{T,i}-X_{T,1})^2$. To test a specified integer endpoint
vector $Y_T$, use $E_T=\sum_i(X_{T,i}-Y_{T,i})^2$.

\begin{theorem}[Unique natural witnesses at every fixed horizon]
\label{ste:ec:thm:quadratic}
For fixed game data and an externally specified $T$, the polynomial
\eqref{ste:ec:eq:complete-quadratic} has integer coefficients and degree at most
two, and is nonnegative on the real nonnegative orthant. For each supplied
initial vector, its natural zero set is empty when the actual endpoint
does not meet the observation, and consists of exactly one complete
witness tuple when the observation holds.

With $b$ binary-decision states and $e$ endpoint squares, it has exactly
\begin{equation}
 \boxed{(N+2b)T\text{ natural witnesses},\quad
        (N+b)T+e\text{ squared linear forms},\quad bT\text{ products}.}
 \label{ste:ec:eq:certificate-counts}
\end{equation}
Here $e=1$ for a scalar pair, $e=N-1$ for consensus, and $e=N$ for a
specified full endpoint. These are counts of the stated construction,
not minimality results.
\end{theorem}

\begin{proof}
There are $N$ next-state coordinates and $2b$ gaps per step. A one-action
state contributes one squared linear form; a binary state contributes
two squares and one complementarity product. This gives the ledger.

All residuals are linear in the variables and all products have degree
two. Nonnegativity prevents cancellation between any two terms. At a
zero, starting from the supplied $X_0$, Lemma~\ref{ste:ec:lem:complementarity}
and the one-action constraints force every $X_1$ coordinate and every
first-step gap. Induction forces the whole trajectory and all gaps.
The endpoint squares then hold exactly when the actual endpoint satisfies
the observation. Conversely that trajectory and those forced gaps make
all terms vanish whenever the observation holds. The case $T=0$ has no
existential variables and reduces to the endpoint test on $X_0$; the
unique tuple in the true case is the empty tuple.
\end{proof}

This proves uniqueness of the \emph{entire} witness tuple, not just of a
projected trajectory. It does not assert that a first-hit witness is
unique when time is also allowed to vary.

\subsection{The fixed-point theorem's nonzero reward}

For $\Phi(v)=\tfrac14\one+\tfrac12S(v)$ use the slightly enlarged scaling
\begin{equation}
 X_t=4d_0(2D)^tv_t,\qquad X_0=4A.
 \label{ste:ec:eq:integer-scale-affine}
\end{equation}
Then
\begin{equation}
 X_{t+1,i}=\opt_a\left(\sum_jP_{ia,j}X_{t,j}
                    +d_0(2D)^{t+1}\right).
 \label{ste:ec:eq:affine-integer-recurrence}
\end{equation}
For fixed $T$, every power appearing here is a known integer coefficient.
The expression is linear in $X_t$ and $d_0$. To test the actual fixed
point $\tfrac12\one$, use
\begin{equation}
 E_T=\sum_{i=1}^N
       \left(X_{T,i}-2d_0(2D)^T\right)^2.
 \label{ste:ec:eq:fixed-endpoint}
\end{equation}
Thus the full fixed-point-attainment certificate is still quadratic,
including in its displayed input parameters $A,d_0$. Its count is
$(N+2b)T$ witnesses, $(N+b)T+N$ squares, and $bT$ products.
The external input conditions are $d_0\geq1$ and $0\leq A_i\leq d_0$
when the payoff vector is required to lie in the unit cube.

For completeness, the implementation permits any rational discount
$\delta=a/b_0\in(0,1)$ and any nonnegative common rational reward $r$.
Let $D$ clear the transition probabilities, let $q$ be a rational constant
point target when one is prescribed, and choose
\[
 m=\lcm(\operatorname{den}(r),\operatorname{den}(q)),\qquad
 \lambda=md_0,
\]
with target denominator one when there is no point target. Then
\[
 X_t=\lambda(b_0D)^tv_t,\quad X_0=mA,
\]
uses the integer recurrence
\[
 X_{t+1,i}=\opt_a\left(a\sum_jP_{ia,j}X_{t,j}
                       +\lambda(b_0D)^{t+1}r\right).
\]
Every forcing term is an emitted integer; no unnamed relation remains in
the exported polynomial. The endpoint $X_{T,i}=\lambda(b_0D)^Tq$ is also
integral. For the main theorem this rule gives exactly
\eqref{ste:ec:eq:integer-scale-affine}--\eqref{ste:ec:eq:fixed-endpoint}.

\subsection{Size and precision}

For a zero-reward run, $\max_iX_{t,i}\leq D^t\max_iA_i$ and each gap is
at most $D^{t+1}\max_iA_i$. The bit length of every witness is therefore
$O(\bits(A)+T\log(D+1))$. For the common-reward construction and cube
inputs, $X_t\leq4d_0(2D)^t\one$, giving the analogous linear-in-$T$
bit-length estimate. The number of witnesses is also linear in $T$.

These bounds concern a horizon supplied in unary or otherwise charged by
its numerical value. They are not polynomial bounds in $\log T$.
The large dense reset rows can make the unexpanded polynomial sizeable:
the count of variables alone does not count the coefficient storage or
the work needed to evaluate all dense linear forms.

\section{Worked examples and complete emitted objects}
\label{ste:ec:sec:examples}

\subsection{A two-state polynomial that can be checked by hand}

Consider the zero-reward game with discount $1/2$ and effective action rows
\begin{align*}
 \text{MIN state: }&\quad(1/2,1/2),\ (1/4,3/4),\\
 \text{MAX state: }&\quad(3/4,1/4),\ (1/2,1/2).
\end{align*}
It is obtained from a half-uniform reset on an averaging/copy game. Start
at $v_0=(1,0)$. The discounted trajectory begins
\[
 v_1=(1/8,3/8),\qquad v_2=(1/8,1/8).
\]
Here $D=4$, and $X_t=8^tv_t$ gives $X_1=(1,3)$ and $X_2=(8,8)$.
Use trajectory variables $(x,y) = X_1$, $(z,w)=X_2$, and gaps
$a,b,c,d,e,f,g,h$. The complete scalar-equality polynomial is
\begin{align}
 Q={}&(2-x-a)^2+(1-x-b)^2+ab\notag\\
 &+(y-3-c)^2+(y-2-d)^2+cd\notag\\
 &+(2x+2y-z-e)^2+(x+3y-z-f)^2+ef\notag\\
 &+(w-3x-y-g)^2+(w-2x-2y-h)^2+gh\notag\\
 &+(z-w)^2.
 \label{ste:ec:eq:small-complete}
\end{align}
Its unique natural zero is
\[
 (x,y,z,w)=(1,3,8,8),\qquad
 (a,b,c,d,e,f,g,h)=(1,0,0,1,0,2,2,0).
\]
The expanded polynomial has 45 nonzero monomials. Both the complete factored
polynomial and its expansion are delivered. The machine-readable variable
ordering interleaves each time slice's coordinates and gaps; its label
array specifies that ordering explicitly.

This small game is not universal. Its purpose is to expose the actual
Diophantine construction, including its degree, natural domain, tie
behavior, and full witness count.

\subsection{A compiled countdown program}

The source has one counter and two control states:
\begin{center}
\begin{tabular}{cl}
\toprule
State & Instruction\\\midrule
0 & If the counter is zero, go to 1; otherwise decrement it and stay at 0.\\
1 & Halt.\\\bottomrule
\end{tabular}
\end{center}
From state zero with counter two, halting is reached after three source
steps. The complete compiler exports two versions.

\begin{center}
\begin{tabular}{lrr}
\toprule
Quantity & Absorbing/scalar version & Erasing/fixed-point version\\
\midrule
Padded game states $N$ & 128 & 128\\
Binary-decision states $b$ & 12 & 10\\
Layered nodes, including inputs & 52 & 41\\
Pipeline period $M$ & 8 & 8\\
Normalization $K$ & 128 & 64\\
Common reward & 0 & $1/4$\\
Initial payoff denominator & 512 & 256\\
First observed event & equality at 17 & full fixed point at 30\\
Natural witnesses at that horizon & 2,584 & 4,440\\
Squared linear forms & 2,381 & 4,268\\
Complementarity products & 204 & 300\\
\bottomrule
\end{tabular}
\end{center}

Both use discount $1/2$, a half-uniform reset, and dyadic full-support rows.
The first equality time is $2M+1=17$, as
\eqref{ste:ec:eq:equality-time} predicts. In the erasing version,
Lemma~\ref{ste:ec:lem:erasure} gives the safe full-drain bound $(H+1)M=32$;
the actual full fixed point is attained earlier, at time 30. The test
checks all coordinates at every preceding time, not only the two observed
registers.

The numerical tables above describe this illustrative one-counter program.
They are not a claim of universality with 128 states. The game exports
include every base action row and the explicit common-reset fraction.
The certificate exports have the reset already expanded into all integer
linear forms; they do not leave transition semantics as an oracle.

\section{Reproducibility, checks, and verification limits}
\label{ste:ec:sec:reproduce}

The archive has one standard-library implementation, a test/export driver,
and a separate checker that does not import the generator. Python 3.13.5
was used; the code uses language features available in Python 3.10 and
later. No floating-point approximations, numerical optimization, or random
search for integer roots are used. The reproducible pseudorandom seed is
20261002; it controls only finite test cases, not theorem selection or
proof validity.

The code exports all linear forms and all complementary products. Thus
\nolinkurl{countdown_certificate.json} and
\nolinkurl{fixedpoint_certificate.json} each specify one complete polynomial
in factored form, a full variable list, a complete natural assignment,
and a resource ledger. Factored representation avoids a gratuitous
expansion of the dense reset terms. The 12-witness example additionally
includes its fully expanded sparse polynomial.

\wtag\ The two large certificates are \emph{not stored} in this report: \path{countdown_certificate.json} (11,921,969 bytes) and \path{fixedpoint_certificate.json} (21,813,796 bytes) are regenerated exactly by the shipped programs, and the README of this report says how and gives their SHA-256 hashes. They also survive in the original archive of the arrival commit. The other delivered artifacts are shipped as \path{data/10-exact-convergence-*.json}; the delivered \path{artifacts/test_run.txt}, a byte copy of \path{verification.json}, is not.

To regenerate the examples and tests:
\begin{lstlisting}
python code/run_checks.py
python code/check_certificate.py artifacts/fixedpoint_certificate.json \
    --game artifacts/fixedpoint_game.json
\end{lstlisting}
\wtag\ These commands use the delivered layout: \path{code/run_checks.py} imports \path{bellman_diophantine} and \path{check_certificate} by their delivered names and rewrites every file of \path{artifacts/}. In this report the programs are \path{code/10-exact-convergence-*.py}; the README gives a recipe that runs on a copy.

The recorded run contains 26,077 assertions. Selected categories are:
\begin{center}
\begin{tabular}{lr}
\toprule
Check category & Assertions\\\midrule
Exhaustive bounded local min/max certificate equivalences & 17,150\\
Raw counter-map successor comparisons & 640\\
Layered counter-map successor comparisons & 640\\
Positive homogeneity comparisons & 640\\
Common-reset coordinate-difference identities & 1,920\\
Original-pipeline checks at every microstep & 248\\
Erasing-pipeline checks at every microstep & 100\\
Chance-only Skolem-lift exact identities & 225\\
Expanded versus factored small-polynomial evaluations & 100\\
Each countdown witness coordinate incremented individually & 2,584\\
Independent complete-export checks & 3\\
\bottomrule
\end{tabular}
\end{center}
The full category list, including ties, input rejection, time-zero cases,
nonzero rewards, rational targets, and initially halted pipelines, is in
\texttt{artifacts/verification.json}. (\wtag\ Shipped as \path{data/10-exact-convergence-verification.json}. The suite was rerun at the write on a copy with the delivered layout: 26,077 assertions, and all eight regenerated exports equal the delivered files apart from line endings.)

For the 2,584 coordinate-mutation checks, the exact change in the polynomial
under a unit increment of one coordinate is computed from its squared-linear
coefficients and the paired gap values; the test does not rerun a dense
polynomial thousands of times. Separate small-instance tests directly
evaluate mutated assignments. The distinction is recorded to avoid
inflating what the test suite independently establishes.

The independent checker validates natural domains, variable indices,
counts, and the integer value of the complete exported polynomial.
When a game export is supplied, it reconstructs the reset rows independently,
checks the input payoff vector, and recomputes every scaled trajectory
coordinate and every gap. This catches serialization or exporter mistakes
that a generator-only self-check could miss.

\wtag\ \emph{A defect in the delivered checker.} The research programme's review of this manuscript~\cite{review808} found that the checker compares the declared input with the game's initialization but replays the trajectory from a separate scaled initial vector that it does not bind to that input. Changing only the declared input numerators of a genuine one-step certificate, and supplying a game with the same changed initial payoffs, still yields polynomial value zero and a successful independent game check although the actual endpoint fails the observation. The integer polynomial itself remains an honest certificate for its own input; the false assertion is the checker's binding. The review's patch \path{exact_convergence_input_binding.patch} binds the input vector, denominator, witness scale, reward, discount, endpoint mode and game, and leaves every valid delivered example accepted; the review also reconstructs every coefficient of the three delivered polynomials from their games and inputs. The checker shipped here, \path{code/10-exact-convergence-check_certificate.py}, is the unpatched original.

None of this proves an infinite theorem by testing. The inductive proofs
in Sections~\ref{ste:ec:sec:counter}--\ref{ste:ec:sec:quadratic} carry that obligation.
No Lean, Rocq, or Isabelle formalization of the new game compiler was
produced or run. The repository's MRDP interface is an inspected source,
not a formal verification of these additional constructions.

\section{From bounded quadratics to a fixed Diophantine relation}
\label{ste:ec:sec:mrdp}

For the fixed game of Theorem~\ref{ste:ec:thm:main}, encode each valid rational
payoff vector by a natural number. Its finite-attainment set $E$ is c.e.
and c.e.-complete. MRDP, in precisely the fixed-arity form exposed by the
repository \cite{repo-mrdp}, yields
\begin{equation}
 \exists k\;\exists P\in\Z[n,y_1,\ldots,y_k]\;
 \forall n\in\N:\quad n\in E\ \Longleftrightarrow\
 \exists y\in\N^k\;P(n,y)=0.
 \label{ste:ec:eq:mrdp}
\end{equation}
The ordinary polynomial $P$ and its witness dimension are fixed before the
input. We do not print its coefficients or claim an operation count for
it. A constructive formal route to MRDP through arithmetic computation is
also described in the published Coq development \cite{h10coq}.

The bounded quadratics do not themselves give a fixed quadratic polynomial
for unbounded time: $Q_T$ has $(N+2b)T$ witnesses, and its time-dependent
coefficients in the reward version also depend on $T$. Writing
``$\exists T,W\;Q_T=0$'' does not turn a varying-arity family into one
first-order polynomial equation.

At each fixed horizon the transition graph is described by finitely many
linear equalities and inequalities, one finite collection for each action
comparison pattern. It is consequently Presburger-definable on its integer
parameters. The quadratic complementarity package compresses this finite
union into one equation over the orthant. The computational difficulty
arises from unbounded iteration, not from a hidden claim that one fixed
quadratic polynomial represents every enumerable set.

\subsection{A quartic existence corollary, with the correct caveat}

One can convert any fixed integer polynomial in \eqref{ste:ec:eq:mrdp} into a
fixed sum of squares of quadratic residuals. Separate its positive and
negative parts, evaluate each by additions and multiplications on natural
registers, and impose equality of their outputs. A multiplication gate
has residual $z-xy$ and an addition gate residual $z-x-y$. Squaring and
summing produces a fixed nonnegative polynomial of degree at most four
with the same projected zero set.

The added circuit registers are uniquely determined by the original
witness $y$. The original MRDP witnesses are not thereby made unique.
Thus neither the single-fold nor the finite-fold Diophantine representation
problem is settled. The horizon-indexed uniqueness in
Theorem~\ref{ste:ec:thm:quadratic} is deliberately a different assertion.

\wtag\ The circuit argument of this subsection is the unique-extension quadratization \texttt{cdc:wf:lem:quadratization} of~\cite{cdc} and \texttt{pqc:lem:quarticlower} of~\cite{pqc}; Part~\ref{ste:part:ee} proves the same corollary for exact erasure (Proposition~\ref{ste:ee:prop:fixedquartic}). No novelty is claimed for it.

\subsection{No computable global search cutoff}

\begin{corollary}
\label{ste:ec:cor:no-cutoff}
For the fixed universal game there is no total computable bound, as a
function of the length of the input payoff encoding, on the first exact
attainment time of all inputs that eventually attain the fixed point.
Likewise, for any fixed polynomial representation of its attainment set,
there is no total computable input-length bound that guarantees a solution
with all witness coordinates below that bound whenever a solution exists.
\end{corollary}

\begin{proof}
Either bound would turn semidecision into decision: compute the bound and
search the finite time interval, or the finite witness box. That would
decide a c.e.-complete set.
\end{proof}

\wtag\ This is the bounded-search argument of \texttt{cdc:bd:prop:nobound} in~\cite{cdc}; for contractions in the collection it appears as \texttt{pqc:sc:cor:gap} and \texttt{pqc:ad:cor:nobound} in Part~VII of~\cite{pqc}, and in Part~\ref{ste:part:ee} as Theorem~\ref{ste:ee:thm:nobound}.

The repository README read during this audit reports 75 arithmetic
operations for its complete certificate and 87 for one universal polynomial
\cite{repoH10}. Nothing here changes either figure. Our counted objects
are native game tables and horizon-dependent quadratic certificates, not
a competing numerically instantiated universal Diophantine polynomial.

\section{Precision, robustness, and why approximation is not a decision}
\label{ste:ec:sec:precision}

The main operator's fixed point is explicit and its approximation error
is uniformly bounded. On compiled initializations the stronger estimate
\[
 \|\Phi^n(v(x))-v_*\|_\infty\leq 4^{-n}/2
\]
holds whether or not the source computation halts. Therefore any prescribed
absolute accuracy is obtained after a computably bounded number of steps,
while exact finite attainment remains undecidable.

For the absorbing scalar construction, a nonzero prehalting register gap
has the exact value in \eqref{ste:ec:eq:exact-gap}. If $K=2^\kappa$, its magnitude
is
\[
 2^{-(1+2n+\kappa\lceil n/M\rceil)}.
\]
For a known finite horizon, its bit requirement is explicit. Two separately
approximated coordinates with absolute error less than one quarter of this
gap can distinguish it from zero, provided a corresponding promise about
the two alternatives is available. There is no contradiction with the
lack of a computable bound on how far one must look for a future exact
zero. Precision requirements and time requirements must not be conflated.

\begin{proposition}[A fixed positive threshold restores a finite cutoff]\label{ste:ec:prop:threshold}
For a zero-reward discounted game on $[0,1]^N$ and a supplied rational
$\varepsilon>0$, the property
$\exists n\;(v_{n,r}-v_{n,s}\geq\varepsilon)$ is decidable.
\end{proposition}

\begin{proof}
Every coordinate of $v_n$ lies in $[0,\delta^n]$, so the difference is at
most $\delta^n$. Once $\delta^n<\varepsilon$, no later time can qualify.
Compute such a cutoff and check the finite prefix exactly.
\end{proof}

\wtag\ The same elementary argument decides threshold reachability for source~14's map in Part~VII of~\cite{pqc} (Propositions \texttt{pqc:sc:prop:away} and \texttt{pqc:sc:prop:threshold}).

A positive absolute tolerance does not solve the zero-threshold predicate:
it changes the predicate. Strict inequalities with threshold zero can
still encode halting, as Section~\ref{ste:ec:sec:universality} shows. Similarly,
full support is not arbitrary robustness against perturbing the transition
table. The theorem proves preservation under one very specific, exactly
known common reset. It does not say that all nearby games have the same
halting language, or that rounded floating-point implementations retain
universality.

Finite-precision implementations have finitely many representable states
under a specified rounding and overflow policy. Their exact reachability
can in principle be decided by finite exploration. That fact is compatible
with the exact-rational result, because the required precision is not
bounded uniformly over all running times.

\section{Research agenda and concrete next targets}
\label{ste:ec:sec:agenda}

\subsection{Instantiate and minimize a fixed universal game}\label{ste:ec:q:instantiate}

The most direct next deliverable is to choose an explicit published
universal counter program, enter its entire instruction table, and export
the resulting game and proof certificate. The output should include the
natural-input loader, not merely a claim about a universal circuit.
Minimize the number of states, binary decisions, stochastic successors,
normalization factors, and the maximum common denominator separately.
These objectives need not agree. The present depth-seven source compiler
makes the microstep period bounded by eight, but it is not a size-optimal
construction.

A particularly concrete optimization is to share delay copies, normalized
linear forms, and zero tests while preserving the two-channel invariant.
Any change must be checked against the actual feedback schedule and full
pipeline drain, not only against one-step circuit equivalence.

\subsection{Remove one player's genuine choices}\label{ste:ec:q:oneplayer}

Does an analogous fixed-game finite-attainment theorem hold when every
nontrivial decision state belongs to MAX, or all belong to MIN? The
present proof does not remove the opposite owner: the negative channel
of a minimum is a maximum. Any attempted one-player translation must
replace this exact dependency, rather than silently switching objectives.
This question is close to the remaining Bellman-operator reachability
frontier in \cite{varonka}, but their exact hypotheses and observation
must be matched before claiming a resolution.

Even a restricted subclass with either a proof of undecidability or a
uniform decision procedure would clarify which portion of the two-player
compiler carries the essential nonlinearity.

\subsection{Prescribed randomness and sparse supports}\label{ste:ec:q:sparse}

Can all chance computations be routed through one shared randomizing
state, or through a constant-size family of fair-coin gadgets? Can one
retain a constant number of stochastic successors per action before the
common reset? A global reset and sparse support cannot both describe the
same expanded row when $N$ grows, so the cost model must distinguish a
native ``reset to uniform'' instruction from its explicitly expanded
probability table.

The current theorem uses an exact common reset. It would be valuable to
characterize all reset kernels that preserve the simulated quotient
modulo $\R\one$, and all nonuniform resets that still admit an invariant
paired-channel subspace. This is a precise algebraic question about the
reset's action on $V$, not a claim of unrestricted noise robustness.

\subsection{One-dimensional observations versus whole-vector attainment}\label{ste:ec:q:observations}

The scalar/whole-vector distinction is already decisive for chance-only
systems: Skolem-equivalence on one side, kernel stabilization on the other.
Investigate the minimum number of observed coordinates required for
hardness in one-player stochastic systems. Related targets include a
specified affine subspace, a consensus line, a nonzero point, and a
fixed point. They should be studied as distinct problems rather than
collapsed under the word ``reachability.''

For the two-player construction, can full fixed-point attainment be reduced
to equality of one designated pair with no additional input promises?
The present reductions give completeness of each problem separately;
they do not furnish a global equivalence on every arbitrary initialization.

\subsection{Witness-economical native certificates}\label{ste:ec:q:economical}

The quadratic compiler spends two natural gaps per binary comparison.
Determine the minimum arity of a single-fold orthant-nonnegative quadratic
representation of the binary min and max graphs under a fixed interface.
One gap can encode distance from one action but loses the sign of the
other difference; naive elimination risks increasing degree or introducing
multiple witnesses. A useful result would be a proved lower bound under
clearly stated interface restrictions, or an explicit smaller compiler.

Time-dependent reward terms are emitted as integer coefficients. A
shared canonical clock could trade coefficient size for additional
variables. Compare the two representations using bit complexity and
literal evaluator costs, not only total polynomial degree.

\subsection{Integrate with the repository's fixed-arity MRDP route}\label{ste:ec:q:mrdp}

A formal interface should expose exact game-step evaluation, its
computability, and the semidecision procedure for the chosen observation.
The repository's MRDP theorem can then supply the abstract fixed polynomial.
Producing a fully expanded native-to-MRDP compiler with audited operation
counts is a separate engineering and mathematical project.

Track three properties independently: preservation of the accepted input
set, preservation of projected trajectories, and preservation of complete
witness tuples. The present bounded compiler proves all three at fixed
$T$. An unbounded history compression must not inherit the strongest
property without a separate proof.

\subsection{Formalize the proof-critical interfaces}\label{ste:ec:q:formalize}

The highest-value formalization targets are the counter encoding,
layer-normalization invariant, every-microstep pipeline lemma, erasing
pipeline drain, common-reset identity, and the zero-set uniqueness theorem.
These are short enough to be individually inspectable and strong enough
to expose most plausible compiler mistakes.

A proof-carrying exporter could attach a machine-checkable derivation of
each stochastic row and each quadratic residual. The independent checker
in this archive is a useful testing boundary, but it does not replace a
kernel-checked theorem connecting exported coefficients to the original
counter semantics.

\section{Conclusion}

The main result isolates a computationally universal exact iteration
inside a tightly constrained stochastic-game interface. Its transition
rows are positive and dyadic; its discount is fixed; its limiting value
is obvious; and its error decays geometrically. None of these facts
supplies a decision procedure for finite exact attainment from arbitrary
rational initial values.

The construction's nontrivial obligations are explicit: signed values are
realized by paired channels, nested computation is realized by a cyclic
pipeline, initialization accounts for every phase, and halting erases the
entire pipeline rather than only its visible registers. Uniform resetting
preserves the computation in a precise invariant subspace.

At fixed time, the same substrate has a particularly economical arithmetic
semantics: one orthant-nonnegative quadratic polynomial with a unique
complete natural witness whenever its endpoint condition holds. At
unbounded time, MRDP supplies a fixed Diophantine relation, but not the
bounded family's uniqueness. The chance-only and selector-only boundaries
explain why both the observation and the computational primitives must be
specified carefully. These distinctions are the foundation for further
native compilers and formal developments, rather than shortcuts around
them.


\part{Thermal arithmetic (source 03)}\label{ste:part:ta}

\noindent\wtag\ Part~III prints batch-78 manuscript~03, \emph{Thermal Arithmetic: Diophantine Witnesses, Effective Excited Spectra, and Undecidable Confinement Boundaries}, under the kicker ``ProveIt research continuation'', ``Research report prepared for Vladimir Reshetnikov'', October 2, 2026 (the manuscript names no assistant). Its PDF metadata gave the shorter subtitle \emph{Diophantine Witnesses and Undecidable Confinement Boundaries}; the printed subtitle is used here. Its labels carry the prefix \texttt{ste:ta:}. Its abstract and status statement follow; its appendix is Appendix~\ref{ste:ta:app:reproduction}.

\begin{quote}\small
\textbf{Abstract (source 03).} An existential Diophantine representation controls whether witnesses exist, but generally not how many there are or how far away they lie. We construct a witness-faithful spectral normal form that separates these difficulties. A natural arithmetic circuit with $k$ witness inputs and $g$ materialized gates gives a diagonal, nonnegative, integer-valued Hamiltonian on $m=k+g+1$ occupation modes, of joint polynomial degree at most five, with an explicit decomposition into positive commuting terms supported on at most four modes. Its zero-energy basis states are in bijection with the original witnesses, while every positive-energy state satisfies $E\ge1+\sum_i n_i$.

This radial isolation makes the positive-energy resolvent effectively compact and the excited thermal trace uniformly computable, even when the zero-energy space is infinite. Compact resolvent and finite unconfined thermal trace are equivalent to finite witness multiplicity. A free-witness padding construction gives unconditional undecidability of both properties on a fixed universal family; requiring compactness on every input instead is exactly a finite-fold representation condition.

A second construction produces a uniformly computable critical fugacity $t_*(c)\in(0,1]$ for each encoded computation. The exact assertion $t_*(c)<1$ is equivalent to halting, although every regulated partition function and every finite-precision approximation to $t_*(c)$ is computable. We prove explicit tail estimates, quantitative witness-height bounds, and the absence of uniform computational cutoffs. An exact Python compiler and 7,936 passing regression assertions accompany the proofs.
\end{quote}
\begin{quote}\small
\textbf{Status (source 03).} The proofs below establish the stated normal-form and boundary results from standard Diophantine representability. Their priority in the literature is not established. This report does not resolve the finite-fold conjecture, claim physical hypercomputation, provide a numerical universal polynomial, or claim proof-assistant verification.
\end{quote}


\section{Research question, scope, and relationship to ProveIt}\label{ste:ta:sec:intro}
\subsection{Why another computational substrate?}
The central question of this report is not simply whether a Turing-complete model has a Diophantine halting predicate. Once a model has an effective finite observation semantics, standard representability already supplies such a predicate. The more informative question is whether an arithmetic representation can separate three independent obstacles:
\[
\text{existence of witnesses},\qquad
\text{multiplicity of witnesses},\qquad
\text{height of witnesses}.
\]
We use a finite number of occupation-number modes, each with infinitely many basis states. The substrate is a static polynomial energy landscape, not a step-by-step simulator. Its zero-energy questions represent computation; its positive-energy states have effective, uniform confinement. In particular, infinitely many witnesses can coexist with an entirely computable excited-state thermal trace.

The repository's October 2 review of computational substrates distinguishes horizon-dependent certificates from one fixed universal polynomial, and existence-preserving compilation from witness-preserving compilation. It also identifies uncontrolled multiplicity in a fixed-quartic MRDP route \cite{repoReview}. (\wtag\ The fixed-quartic route meant is Corollary \texttt{pqc:tu:cor:fixedquartic} of~\cite{pqc}, Part~VIII, and the remark after it, which the cited review discusses.) Here multiplicity is no longer an optional implementation detail: it is exactly the obstruction to compact resolvent and to a finite unconfined thermal trace. The present construction builds on that distinction rather than proposing another uncharged bounded-history encoding.

\subsection{Source snapshot and limits of the audit}
The repository was inspected at commit
\begin{center}
\texttt{6914ccca6685baf53b7a35f25efc89366c76ba74}.
\end{center}
The relevant audit file has blob identifier
\begin{center}
\texttt{8bd07142c8b4298fb636e47b40d2d84e993a40ce}.
\end{center}
The repository's reported 87-operation fixed universal polynomial remains a separate arithmetic-circuit result \cite{repoH10}. This report neither improves that bound nor counts an existential MRDP polynomial as a free primitive in a numerical operation ledger. The source inspection was not a rebuild of the Lean or Rocq developments, and the proofs below do not assume that every research manuscript in the repository has been formally verified.

\subsection{Established foundations and the proposed contribution}
We use the standard fact that every computably enumerable relation over natural numbers has an existential polynomial representation, and that one fixed polynomial can represent a universal such relation \cite{jones}. A natural arithmetic-circuit expansion then replaces a high-degree polynomial equation by quadratic gate equations with uniquely determined intermediate values. We give that expansion explicitly because its witness multiplicity matters here.
(\wtag\ This expansion is the unique-extension quadratization of the collection's report~\cite{cdc}, Lemma \texttt{cdc:wf:lem:quadratization}, with the definition of witness-faithfulness given there; see also \texttt{cdc:cs:lem:quartic} in~\cite{cdc} and, without the uniqueness statement, \texttt{pqc:lem:quarticlower} in~\cite{pqc}. In Theorem~\ref{ste:ta:thm:compiler} the zero-set bijection and the quartic sum of squares are that lemma; new here are the marker $a-1$, the radial factor $1+\sum_in_i$, the bound $H>0\Rightarrow H\ge1+\sum_in_i$, the degree five and the $r(m+1)$ four-mode ledger.)

The basic substitution of number operators into a squared Diophantine polynomial is not new. Kieu already writes a Hamiltonian of the form $D(N_1,\ldots,N_k)^2$ and identifies zero energy with an integer solution \cite{kieu}. We use only that mathematical observation, not any proposed physical decision procedure. Our additional factor $1+\sum_i N_i$ removes all nonzero-energy flat directions while leaving the zero set unchanged. The effective spectral decomposition, multiplicity classification, and critical-fugacity construction are the results developed and proved here.
(\wtag\ Part~V of~\cite{pqc} (source~13 there, Theorem \texttt{pqc:ql:thm:compiler}) builds a witness-faithful quartic energy landscape for Boolean circuits over $\R$ and states that it treats no Gibbs distribution; this Part is the unbounded natural counterpart, with a thermal trace.)

Finite-fold and single-fold representation questions are discussed in the primary literature \cite{matiyasevichFF,cantone}. The equivalences below are reformulations of those representation conditions, not answers to them. In particular, we do not silently replace an arbitrary MRDP representation by one with finitely many witnesses.

\subsection{Main conclusions at a glance}
Let $Q=\sum_j f_j^2$ be the compiled quartic penalty and $s=\sum_i n_i$. The construction is
\begin{equation}\label{ste:ta:eq:main}
H=(1+s)Q.
\end{equation}
Its essential features are:
\begin{enumerate}[label=(\roman*),itemsep=2pt]
\item $\ker H$ has one occupation basis vector for each original witness, with no extra auxiliary multiplicity.
\item The positive-energy sector is effectively confined: $H>0$ implies $H\ge1+s$.
\item Compact resolvent and finite thermal trace hold exactly when there are finitely many witnesses.
\item A regulated partition function has a computable threshold location whose exact contact with the boundary $t=1$ encodes nonhalting.
\end{enumerate}
The novelty claim is deliberately limited: this is a proved package of normal-form, effective-estimation, and boundary results proposed for further research. The literature search and repository audit do not establish that every individual consequence first appears here.

\section{Arithmetic conventions and the witness-faithful compiler}
\subsection{Natural domains and representations}
Throughout, $\N=\{0,1,2,\ldots\}$. A parameter tuple $p\in\N^d$ is fixed when studying an individual operator. Witness coordinates range independently over $\N$. A polynomial $P\in\Z[p,x_1,\ldots,x_k]$ represents a relation by
\begin{equation}\label{ste:ta:eq:representation}
A(p)\quad\Longleftrightarrow\quad
\exists x\in\N^k\; P(p;x)=0.
\end{equation}
A representation is \emph{finite-fold} if each parameter fiber has finitely many witnesses, including zero witnesses outside the relation. It is \emph{single-fold} if each fiber has at most one witness. A representation originally written over positive integers can be moved to our convention by the bijective substitution $x_i=u_i+1$.

All polynomial degrees stated below are total degrees jointly in parameters and occupation variables unless explicitly qualified. Parameter values can be arbitrarily large, but the polynomial expression and the number of coordinates do not change with their values.

\subsection{Natural arithmetic circuits}
Write
\[
P=P^+-P^-,\qquad P^+,P^-\in\N[p,x].
\]
Choose finite directed acyclic circuits for these two nonnegative polynomials using constants in $\N$, parameter and witness inputs, and binary addition and multiplication. Common subexpressions may be shared. Each actual gate introduces one natural coordinate $u_j$, in a topological order. The corresponding residual is one of
\begin{equation}\label{ste:ta:eq:gates}
u_j-v-w,\qquad u_j-vw,
\end{equation}
where $v,w$ are constants, parameters, original witness coordinates, or outputs of earlier gates. The final residual is the difference between the two output atoms. A final, separate coordinate $a$ has residual $a-1$.

The purpose of $a$ is not arithmetical expressiveness. It makes the occupation vacuum a positive-energy state, which will ensure that the thermal threshold starts strictly below its target level. It adds exactly one coordinate and one equation, and is uniquely determined at every zero.

\begin{theorem}[Witness-faithful degree-five compiler]\label{ste:ta:thm:compiler}
Suppose the chosen natural circuit has $k$ witness inputs and $g$ materialized gates. There is an effectively constructed integer polynomial $H(p;n)$ with
\begin{equation}\label{ste:ta:eq:ledger}
m=k+g+1,\qquad r=g+2,
\end{equation}
occupation coordinates and residual equations, respectively, satisfying the following properties.

For every fixed $p$, projection onto the first $k$ coordinates is a bijection
\[
\{n\in\N^m:H(p;n)=0\}
\longrightarrow
\{x\in\N^k:P(p;x)=0\}.
\]
Every occupation value of $H$ is a nonnegative integer, $H(p;0)\ge1$, and
\begin{equation}\label{ste:ta:eq:radial}
H(p;n)>0\quad\Longrightarrow\quad H(p;n)\ge1+\sum_{i=1}^m n_i.
\end{equation}
The joint degree of $H$ is at most five. It has a displayed decomposition into $r(m+1)$ nonnegative commuting terms, each supported on at most four occupation modes.
\end{theorem}
\begin{proof}
Let $f_1,\ldots,f_r$ be all gate residuals, the final output difference, and $a-1$. Put
\[
Q(p;n)=\sum_{j=1}^r f_j(p;n)^2,
\qquad H(p;n)=\left(1+\sum_{i=1}^m n_i\right)Q(p;n).
\]
Since every residual is an integer, $Q=0$ exactly when all residuals vanish. Fix $p$ and the original witness tuple $x$. Topological induction through the circuit forces each $u_j$ to be the natural value of its gate. Thus there is exactly one assignment of intermediate coordinates satisfying the gate equations. The marker is then uniquely $a=1$. At this assignment the remaining output residual is $P(p;x)$. Consequently a source root has exactly one zero-energy lift, and every zero-energy occupation tuple is such a lift.

Each residual has joint degree at most two. Therefore $\deg Q\le4$ and $\deg H\le5$. Every factor $1+\sum_i n_i$ is positive on the natural domain, and every nonzero $Q$ is an integer at least one. This proves \eqref{ste:ta:eq:radial}. At the vacuum, the marker residual equals $-1$, so $Q(p;0)\ge1$.

Finally, expand only the radial factor, not the residual squares:
\begin{equation}\label{ste:ta:eq:local}
H=\sum_{j=1}^r f_j^2+
\sum_{i=1}^m\sum_{j=1}^r n_i f_j^2.
\end{equation}
A gate residual involves at most three occupation modes; the output residual involves at most two, and the marker one. Multiplication by $n_i$ adds at most one mode to the support. Each displayed term is nonnegative on the natural domain. As diagonal multiplication operators all these terms commute. There are $r+mr$ terms in this unsimplified display.
\end{proof}

\begin{remark}[What this ledger does not count]
The number $r(m+1)$ counts displayed positive summands, not additions and multiplications in an optimized evaluation circuit. Formula \eqref{ste:ta:eq:main} is usually cheaper to evaluate than fully expanding \eqref{ste:ta:eq:local}. Constant bit lengths are also not charged by this ledger. Four-mode support does not mean nearest-neighbor geometry, bounded local Hilbert-space dimension, or bounded operator norm. If the degree is measured in creation and annihilation operators rather than number operators, substituting $N_i=a_i^\dagger a_i$ gives a bound of ten instead of five.
\end{remark}

\begin{corollary}[Fixed universal family]\label{ste:ta:cor:universal}
There is a fixed $m$ and one fixed integer polynomial $H(c;n)$ of joint degree at most five such that
\[
c\in\mathsf K\quad\Longleftrightarrow\quad\ker H_c\ne\{0\},
\]
where $\mathsf K$ is a fixed computably enumerable complete halting set, and every fiber has the properties of Theorem~\ref{ste:ta:thm:compiler}.
\end{corollary}
\begin{proof}
Choose once and for all a polynomial $U(c;x)$ representing $\mathsf K$, as supplied by Diophantine representability. Apply the fixed finite compiler to $U$. Only $c$ changes from instance to instance.
\end{proof}

No numerical value of $m$ is claimed here: the universal polynomial is not instantiated in the accompanying code. This distinction prevents a low-degree existence theorem from being mistaken for a competitive numerical universal equation.

\subsection{A small example with a complete resource count}
For $P(p;x,y)=xy-p$, use one gate $u=xy$. The four occupation coordinates are $(x,y,u,a)$ and
\begin{align}
Q_p&=(u-xy)^2+(u-p)^2+(a-1)^2,\\
H_p&=(1+x+y+u+a)Q_p.\label{ste:ta:eq:productexample}
\end{align}
Thus $k=2$, $g=1$, $m=4$, $r=3$, and there are $15$ displayed positive terms. For $p=6$, the four zero-energy tuples are
\[
(1,6,6,1),\ (2,3,6,1),\ (3,2,6,1),\ (6,1,6,1).
\]
For $p=0$, infinitely many tuples occur because $xy=0$ has infinitely many natural solutions. The circuit has preserved that distinction exactly; it has not attempted to regularize the source multiplicity.

\section{The operator and its effectively confined positive spectrum}
\subsection{Domains are part of the construction}
Fix the parameters and suppress them from the notation. Let
\[
\Hil_m=\ell^2(\N^m),\qquad
N_i\ket n=n_i\ket n.
\]
For the polynomial $H(n)$, define the multiplication operator with maximal domain
\begin{equation}\label{ste:ta:eq:domain}
\mathcal D(H)=\left\{\psi=(\psi_n):
\sum_n H(n)^2\abs{\psi_n}^2<\infty\right\},
\qquad (H\psi)_n=H(n)\psi_n.
\end{equation}
Although the formal expression is a polynomial in commuting number operators, \eqref{ste:ta:eq:domain} specifies the actual closed operator used in all statements.

\begin{lemma}[Self-adjoint realization]\label{ste:ta:lem:selfadjoint}
The operator \eqref{ste:ta:eq:domain} is nonnegative and self-adjoint. Its restriction to finitely supported vectors is essentially self-adjoint, with closure equal to \eqref{ste:ta:eq:domain}.
\end{lemma}
\begin{proof}
Symmetry follows from the real diagonal coefficients. If $\psi$ belongs to the adjoint domain of the finite-support restriction, testing against each basis vector forces the adjoint coordinates to be $H(n)\psi_n$. Membership of that vector in $\ell^2$ is precisely \eqref{ste:ta:eq:domain}. Thus the adjoint of the restriction is the maximal multiplication operator. The maximal operator is symmetric, and the same coordinate argument identifies its adjoint with itself.

For the core assertion, truncate any $\psi\in\mathcal D(H)$ to a finite box. Both the $\ell^2$ tail of $\psi$ and that of $H\psi$ tend to zero. Hence finite-support truncations converge in graph norm. Nonnegativity follows from $H(n)\ge0$ coordinatewise.
\end{proof}

\subsection{Energy layers and an exact finite search}
For $n\in\N^m$, put $s(n)=\sum_i n_i$. Let
\[
\mathcal Z=\{n:Q(n)=0\},\qquad
M=\#\mathcal Z\in\N\cup\{\infty\}.
\]
The radial lower bound is assumed in the next results; consequently they apply both to the compiler and to any directly supplied nonnegative integer polynomial $Q$ used in \eqref{ste:ta:eq:main}.

\begin{theorem}[Effective positive-energy layers]\label{ste:ta:thm:layers}
For every integer $E\ge1$, all basis states with $0<H(n)\le E$ lie in the finite simplex $s(n)\le E-1$. Their number, with multiplicity, is at most
\begin{equation}\label{ste:ta:eq:countbound}
\binom{E+m-1}{m}.
\end{equation}
Every positive eigenvalue is an integer of finite multiplicity, and its exact multiplicity is computable uniformly from the polynomial and the eigenvalue. There is no finite accumulation point of distinct eigenvalues.
\end{theorem}
\begin{proof}
The inequality $H(n)\ge1+s(n)$ on nonzero states gives the simplex restriction. The number of natural $m$-tuples of sum at most $E-1$ is \eqref{ste:ta:eq:countbound}: add a slack coordinate to get a tuple of sum $E-1$, and count placements of separators. Polynomial evaluation on this finite set decides each positive-energy membership exactly.

The spectrum of a diagonal multiplication operator is the closure of its diagonal values. To see the only potentially nontrivial direction directly, a complex number at positive distance from all those values has a bounded diagonal inverse for $H-\lambda I$; the other direction follows from basis approximate eigenvectors. Here the values are nonnegative integers, so the set of distinct values has no finite accumulation point and is already closed.
\end{proof}

There is also a useful divisor refinement. If $H(n)=E>0$, then
\begin{equation}\label{ste:ta:eq:divisors}
1+s(n)\mid E,\qquad Q(n)=\frac{E}{1+s(n)}.
\end{equation}
Thus one may enumerate only shells $s=d-1$ for divisors $d\mid E$, rather than the entire simplex. This changes the enumeration strategy, not the decidability result.

\subsection{Separating the zero projection from a computable compact part}
Let $P_0$ be the diagonal projection onto $\ell^2(\mathcal Z)$, and let $C_B$ denote projection onto
\[
\cB=\{0,1,\ldots,B\}^m.
\]
Define $K$ and $T_q$, for $0<q<1$, by their diagonal entries:
\begin{align}
K(n)&=\begin{cases}(1+H(n))^{-1},&H(n)>0,\\0,&H(n)=0,\end{cases}\label{ste:ta:eq:K}\\
T_q(n)&=\begin{cases}q^{H(n)},&H(n)>0,\\0,&H(n)=0.\end{cases}\label{ste:ta:eq:T}
\end{align}

\begin{theorem}[Effective compact and trace-class remainders]\label{ste:ta:thm:remainders}
The following decompositions hold:
\begin{equation}\label{ste:ta:eq:decomposition}
(H+I)^{-1}=P_0+K,\qquad q^H=P_0+T_q.
\end{equation}
The operator $K$ is compact, and $T_q$ is trace class. For the finite-rank truncations $K_B=C_BKC_B$ and $T_{q,B}=C_BT_qC_B$,
\begin{align}
\norm{K-K_B}&\le\frac1{B+3},\label{ste:ta:eq:Ktail}\\
\norm{T_q-T_{q,B}}_1&\le
\frac{m q^{B+2}}{(1-q)^m}.\label{ste:ta:eq:Ttail}
\end{align}
For rational $q$, these are effective approximation procedures using exact rational finite matrices, uniformly in all polynomial parameters.
\end{theorem}
\begin{proof}
The identities hold on each basis vector. Outside $\cB$, at least one coordinate is at least $B+1$, so $s(n)\ge B+1$. On positive-energy states this gives $H(n)\ge B+2$. The supremum of the omitted diagonal entries of $K$ is at most $1/(B+3)$, proving \eqref{ste:ta:eq:Ktail} and compactness.

For the heat remainder, nonnegative diagonal entries give a trace norm equal to their sum. Since $q^{H(n)}\le q^{1+s(n)}$ on positive states, a union bound over the coordinate that leaves the box yields
\begin{align*}
\sum_{n\notin\cB,\,H(n)>0}q^{H(n)}
&\le q\sum_{i=1}^m\left(\sum_{n_i=B+1}^\infty q^{n_i}\right)
\prod_{j\ne i}\left(\sum_{n_j=0}^\infty q^{n_j}\right)\\
&=\frac{m q^{B+2}}{(1-q)^m}.
\end{align*}
The same geometric calculation over the whole space bounds the trace by $q/(1-q)^m$. Exact evaluation of $Q(n)$ tells the finite algorithm which entries are zero and which belong to the remainder.
\end{proof}

\begin{remark}[Effective strong information is not effective norm information]
For each explicitly supplied basis index $n$, membership in $\mathcal Z$ is decidable by evaluating $Q(n)$. The action of $P_0$ is consequently computable on vectors presented with effective $\ell^2$ approximations: project a finite approximation and use $\norm{P_0}\le1$ to control the tail. This does not supply a uniform finite-rank approximation to $P_0$ in operator norm. If any zero state lies outside a proposed cutoff, the omitted projection has norm one. The distinction is exactly where an attempted finite spectral decision procedure can fail.
\end{remark}

\section{Multiplicity, compactness, and thermal finiteness}
\begin{theorem}[Exact multiplicity classification]\label{ste:ta:thm:classification}
For the radially isolated Hamiltonian $H=(1+s)Q$,
\[
\dim\ker H=M.
\]
The following are equivalent:
\begin{enumerate}[label=(\roman*),itemsep=1pt]
\item $M$ is finite;
\item $H$ has compact resolvent;
\item $q^H$ is trace class for some $q\in(0,1)$;
\item $q^H$ is trace class for every $q\in(0,1)$.
\end{enumerate}
When $M<\infty$,
\begin{equation}\label{ste:ta:eq:finiteTrace}
\Tr(q^H)=M+\Tr(T_q).
\end{equation}
When $M=\infty$, the essential spectrum is exactly $\{0\}$. When $M<\infty$, the essential spectrum is empty.
\end{theorem}
\begin{proof}
A vector is killed by $H$ exactly when its support lies in $\mathcal Z$, proving the dimension statement. By Theorem~\ref{ste:ta:thm:remainders}, $(H+I)^{-1}$ differs from $P_0$ by a compact operator. An orthogonal projection is compact exactly when its range is finite dimensional: an infinite orthonormal sequence in its range has no norm-convergent image subsequence, while a finite-rank projection is compact. This proves (i)$\Leftrightarrow$(ii).

Similarly, $q^H=P_0+T_q$ with $T_q$ positive trace class. Its diagonal sum is $M+\Tr(T_q)$, with the usual extended interpretation when $M=\infty$. Thus (i), (iii), and (iv) are equivalent, and \eqref{ste:ta:eq:finiteTrace} follows. The positive eigenvalues are isolated and of finite multiplicity by Theorem~\ref{ste:ta:thm:layers}; only the eigenvalue zero can have infinite multiplicity. This proves the essential-spectrum assertions, using the characterization by failure to be an isolated finite-multiplicity eigenvalue.
\end{proof}

\begin{remark}[Ground states versus zero-energy states]
If $M=0$, the operator still has a ground-state eigenvalue: a nonempty subset of positive integers has a minimum. What fails is specifically the existence of a \emph{zero-energy} state. If $M=1$, there is a unique zero-energy state up to scalar multiples. These are different statements from uniqueness of a positive-energy ground state in a root-free instance.
\end{remark}

\subsection{A spectral equivalent of finite-fold representability}
\begin{definition}
A polynomial diagonal family is a single integer polynomial $F(p;n)$ in the parameter and finitely many occupation coordinates, nonnegative on all natural tuples, interpreted by the maximal multiplication operator. The number of occupation coordinates is fixed, independent of $p$.
\end{definition}

\begin{theorem}[Finite-fold/compact-family equivalence]\label{ste:ta:thm:finitefold}
For a relation $A\subseteq\N^d$, the following conditions are equivalent:
\begin{enumerate}[label=(\roman*),itemsep=1pt]
\item $A$ has a finite-fold natural Diophantine representation.
\item There is a polynomial diagonal family $F_p$ with compact resolvent for every $p$, such that $p\in A$ exactly when $\ker F_p\ne\{0\}$.
\end{enumerate}
In (ii), whenever (i) holds, one may choose joint degree at most five and the four-mode positive decomposition of Theorem~\ref{ste:ta:thm:compiler}. The analogous single-fold condition is equivalent to a polynomial diagonal representation with $\dim\ker F_p\in\{0,1\}$ for every $p$; the forward construction can also have compact resolvent on every fiber.
\end{theorem}
\begin{proof}
For (i)$\Rightarrow$(ii), compile a finite-fold representation. The zero-set bijection preserves finite fiber sizes, and Theorem~\ref{ste:ta:thm:classification} supplies compact resolvent on every fiber.

Conversely, $F(p;n)=0$ is already a natural polynomial representation of $A$. Compact resolvent forces every eigenspace, in particular the zero eigenspace, to be finite dimensional. Distinct zero tuples supply orthogonal basis vectors, so there are finitely many such tuples. This is a finite-fold representation. Replacing ``finitely many'' by ``at most one'' proves the single-fold assertion in both directions.
\end{proof}

The joint polynomial requirement matters. An arbitrary algorithm that selects a polynomial from $p$ does not itself give one polynomial relation in $(p,n)$, and invoking representability for that selection may introduce additional witnesses whose multiplicity is uncontrolled. No such substitution occurs in Theorem~\ref{ste:ta:thm:finitefold}.

This equivalence explains why the apparently stronger assertion ``there is a universal halting family with compact resolvent on every input'' must not be obtained by simply announcing that circuit auxiliaries are unique. Circuit auxiliaries are unique \emph{over each source witness}; the source fiber may still be infinite.

\wtag\ \emph{Credit and a consequence.} Theorem~\ref{ste:ta:thm:finitefold} follows the template of the finite-fold substrate equivalence \texttt{cdc:cs:thm:finitefold} of~\cite{cdc}: a unique-extension quartic in the forward direction, and the polynomial itself as a representation in the backward direction. What it adds is the replacement of ``finite fibre'' by ``compact resolvent'' through Theorem~\ref{ste:ta:thm:classification}. Combined with the equivalence \texttt{cdc:bd:thm:universal} of~\cite{cdc} it gives the following statement, recorded here as an editorial consequence of the two printed theorems, not as a result of source~03.

\begin{remark}[\wtag\ All-fibre compactness for universal halting]\label{ste:ta:rem:universalcompact}
The following are equivalent: (i)~every computably enumerable subset of $\N$ has a finite-fold Diophantine representation over $\N$; (ii)~there is one polynomial diagonal family $F(e,x;n)$, in the sense of the definition above, with compact resolvent for every $(e,x)$ and $\ker F_{(e,x)}\ne\{0\}$ exactly when program $e$ halts on input $x$. When they hold, $F$ can be chosen of joint degree at most five with the four-mode positive decomposition of Theorem~\ref{ste:ta:thm:compiler}. Indeed, (ii) gives by Theorem~\ref{ste:ta:thm:finitefold} a finite-fold representation of universal halting, which is statement~(b) of \texttt{cdc:bd:thm:universal} and hence equivalent to~(i), its statement~(a); conversely (i) gives (b), and Theorem~\ref{ste:ta:thm:finitefold} compiles that representation into a family as in~(ii). The equivalence decides neither statement.
\end{remark}

\subsection{An unconditional undecidable compactness family}
\begin{theorem}[Padding transfers halting to loss of compactness]\label{ste:ta:thm:padding}
Let $Q_c(n)$ be the fixed universal compiled penalty in Corollary~\ref{ste:ta:cor:universal}. Add one free coordinate $y\in\N$ and define
\begin{equation}\label{ste:ta:eq:padding}
\widetilde H_c(n,y)=\left(1+y+\sum_i n_i\right)Q_c(n).
\end{equation}
Then
\begin{align*}
c\notin\mathsf K
&\Longleftrightarrow \widetilde H_c\text{ has compact resolvent}\\
&\Longleftrightarrow \Tr(q^{\widetilde H_c})<\infty\text{ for some/every }0<q<1,
\end{align*}
whereas $c\in\mathsf K$ gives an infinite-dimensional zero eigenspace and infinite thermal trace for every such $q$. Thus compactness and thermal finiteness are co-computably-enumerable complete on this fixed padded family; their negations are computably enumerable complete.
\end{theorem}
\begin{proof}
If $Q_c$ has no root, the padded penalty has no root. If it has one root $n$, every $(n,y)$ is a root, so its root set is infinite. The radial factor includes the new coordinate and hence satisfies the positive-energy lower bound in $m+1$ modes. Theorem~\ref{ste:ta:thm:classification} gives the asserted analytic alternatives. The exact equivalences with membership and nonmembership in $\mathsf K$ give both reductions and the stated enumerability classifications on this family.
\end{proof}

No assumption about finite-fold Diophantine representations occurs in this proof. Also, no claim is made that the unrestricted finite-root problem has exactly this complexity. For a general polynomial, finite root cardinality has the elementary upper-bound description
\[
\exists B\;\forall n\;\bigl(s(n)>B\Longrightarrow Q(n)\ne0\bigr),
\]
but identifying a sharper unrestricted classification would require a separate argument.

\section{Regulated partition functions and their analytic boundary}
\subsection{Two thermal parameters}
Fix $0<q<1$. It is convenient to regard $q=e^{-\beta}$ as an inverse-temperature parameter in exponential coordinates. The second parameter $0\le t<1$ damps large total occupation independently of the energy:
\begin{equation}\label{ste:ta:eq:Z}
Z(q,t)=\sum_{n\in\N^m}q^{H(n)}t^{s(n)}.
\end{equation}
Equivalently, this is the trace of the positive diagonal operator $q^H t^{\sum_i N_i}$. We use $0^0=1$ in evaluating $t=0$. The terminology \emph{fugacity} is a mathematical parametrization; no physical particle reservoir or laboratory implementation is assumed.

Separate the zero-state and excited-state contributions:
\begin{align}
G(t)&=\sum_{n\in\mathcal Z}t^{s(n)},\label{ste:ta:eq:G}\\
R(q,t)&=\sum_{n\notin\mathcal Z}q^{H(n)}t^{s(n)},\label{ste:ta:eq:R}\\
Z(q,t)&=G(t)+R(q,t).\label{ste:ta:eq:ZR}
\end{align}
For finite $M$, $G$ is a polynomial recording witness heights. For infinite $M$, it is an infinite power series with nonnegative integer coefficients. Unique circuit completion makes these coefficients count the source witnesses at their \emph{compiled} total heights; the heights generally differ from those of the original witness coordinates alone.

\begin{theorem}[Effective partition functions]\label{ste:ta:thm:partition}
For rational $q,t$ in the domains below, the following sums are uniformly computable with exact rational enclosures:
\begin{enumerate}[label=(\roman*),itemsep=1pt]
\item $Z(q,t)$ for $0<q<1$ and $0\le t<1$;
\item $R(q,t)$ for $0<q<1$ and $0\le t\le1$.
\end{enumerate}
If $Z_B$ and $R_B$ denote sums over the box $\cB$, then
\begin{align}
0\le Z(q,t)-Z_B(q,t)&\le
\frac{m t^{B+1}}{(1-t)^m},\label{ste:ta:eq:Ztail}\\
0\le R(q,t)-R_B(q,t)&\le
\frac{m q(qt)^{B+1}}{(1-qt)^m}.\label{ste:ta:eq:Rtail}
\end{align}
In particular,
\begin{equation}\label{ste:ta:eq:Rbound}
0\le R(q,t)\le\frac{q}{(1-qt)^m}\le\frac{q}{(1-q)^m}.
\end{equation}
At $t=0$, the sums are determined by the single vacuum tuple.
\end{theorem}
\begin{proof}
Since $q^{H(n)}\le1$, the full sum is dominated by
\[
\sum_{n\in\N^m}t^{s(n)}=(1-t)^{-m}.
\]
A union bound over coordinates exceeding $B$ gives \eqref{ste:ta:eq:Ztail}. For the excited sum, use $q^{H(n)}\le q^{1+s(n)}$ and apply the same argument with $qt$ in place of $t$, proving \eqref{ste:ta:eq:Rtail} and \eqref{ste:ta:eq:Rbound}. For rational parameters every finite term is a rational number computed by integer polynomial evaluation and exponentiation. Increasing $B$ until the displayed rational tail is below a prescribed tolerance gives an effective enclosure. The argument is uniform in the instance parameters, since neither tail uses them.
\end{proof}

This theorem does not assert a fast algorithm close to $t=1$. In the full trace, the sufficient box size from \eqref{ste:ta:eq:Ztail} can be very large. It asserts termination for each rational $t<1$ and each requested precision. The excited trace, unlike the full trace, remains effectively computable at $t=1$.

\subsection{A sharp analytic dichotomy}
\begin{theorem}[All boundary obstructions lie in the witness series]\label{ste:ta:thm:analytic}
For fixed $q\in(0,1)$, the function $R(q,t)$ is holomorphic in the complex disk $\abs t<q^{-1}$. On every smaller closed disk $\abs t\le\rho<q^{-1}$, its box tails satisfy
\begin{equation}\label{ste:ta:eq:complexRtail}
\abs{R(q,t)-R_B(q,t)}\le
\frac{m q(q\rho)^{B+1}}{(1-q\rho)^m}.
\end{equation}
Consequently:
\begin{enumerate}[label=(\roman*),itemsep=1pt]
\item If $M<\infty$, then $Z(q,t)$ is holomorphic on $\abs t<q^{-1}$ and $Z(q,1)=M+R(q,1)$.
\item If $M=\infty$, the Taylor series of $Z(q,t)$ at zero has radius of convergence exactly one and $Z(q,t)\to+\infty$ as $t\uparrow1$.
\end{enumerate}
Thus a finite-valued holomorphic continuation of $Z(q,t)$ through $t=1$ exists exactly when $M<\infty$.
\end{theorem}
\begin{proof}
The majorant $q(q\rho)^{s(n)}$ proves locally uniform absolute convergence of the series for $R$ in the larger disk. Each term is a polynomial in $t$, so the limit is holomorphic; the same majorant gives \eqref{ste:ta:eq:complexRtail}.

If $M$ is finite, $G$ is a polynomial and the first conclusion follows. If $M$ is infinite, monotone convergence along real $t\uparrow1$ gives $G(t)\uparrow\infty$, since every zero-state term tends to one. The full series converges absolutely for $\abs t<1$ by the geometric bound. It cannot have radius greater than one, because that would give a finite value at $t=1$, contradicting the divergence just proved. A finite-valued holomorphic continuation in a neighborhood of $1$ would similarly be locally bounded on its real approach and is impossible.
\end{proof}

\begin{corollary}[A promised analytic-continuation problem is undecidable]\label{ste:ta:cor:radius}
For the padded fixed universal family of Theorem~\ref{ste:ta:thm:padding}, choose any fixed rational $q\in(0,1/2)$. The functions
\[
Z_c(q,t)=\sum_{s=0}^\infty b_s(c)t^s,
\qquad
b_s(c)=\sum_{s(n,y)=s}q^{\widetilde H_c(n,y)},
\]
have uniformly computable positive rational coefficients and are uniformly computable at rational $0\le t<1$. Nevertheless, deciding whether their radius of convergence is one or at least $q^{-1}>2$ is undecidable. Equivalently, deciding whether they extend holomorphically through $t=1$ is co-computably-enumerable complete on this family.
\end{corollary}
\begin{proof}
Each coefficient is a finite exact sum over a simplex shell. A halting input has infinitely many padded zero states and hence radius one; a nonhalting input has no zero states and hence radius at least $q^{-1}$. Theorem~\ref{ste:ta:thm:padding} gives the completeness statement.
\end{proof}

This corollary has a genuine promise gap in the possible radii. It is not caused by uncertainty in the Taylor coefficients or by an inability to evaluate the function at individual interior points. The missing information is a global boundary property of the entire computable series.

\section{A computable critical fugacity with an undecidable endpoint}
\subsection{An instance-independent thermal separation}
The marker coordinate in Theorem~\ref{ste:ta:thm:compiler} guarantees $H(0)\ge1$. For a system of $m\ge1$ modes, fix
\begin{equation}\label{ste:ta:eq:qchoice}
q_m=\frac{1}{8(m+1)},\qquad
\epsm=\frac{q_m}{(1-q_m)^m}.
\end{equation}
Bernoulli's inequality gives the explicit estimate
\begin{equation}\label{ste:ta:eq:epsbound}
\epsm\le\frac{q_m}{1-mq_m}
=\frac1{7m+8}<\frac18.
\end{equation}
Thus the entire excited contribution, at any $0\le t\le1$, is less than $1/8$. Let
\[
F(t)=Z(q_m,t),\qquad c_0=F(0)=q_m^{H(0)}\le q_m\le\frac1{16}.
\]
All Taylor coefficients of $F$ are positive, since each occupation tuple has a positive weight and each total-occupation shell is nonempty. Therefore $F$ is continuous and strictly increasing on $[0,1)$.

\begin{theorem}[The critical fugacity]\label{ste:ta:thm:critical}
For every compiled instance define
\begin{equation}\label{ste:ta:eq:tstar}
t_*=
\begin{cases}
\text{the unique }t\in(0,1)\text{ such that }F(t)=1/2,
&\mathcal Z\ne\varnothing,\\
1,&\mathcal Z=\varnothing.
\end{cases}
\end{equation}
This definition is well posed. The real number $t_*$ is uniformly computable from the finite polynomial description. For the fixed universal family,
\begin{equation}\label{ste:ta:eq:endpoint}
c\in\mathsf K\ \Longleftrightarrow\ t_*(c)<1,
\qquad
c\notin\mathsf K\ \Longleftrightarrow\ t_*(c)=1.
\end{equation}
Hence its strict interior condition is computably enumerable complete, and its exact endpoint condition is co-computably-enumerable complete.
\end{theorem}
\begin{proof}
If there is no zero state, $F=R\le\epsm<1/8$, so there is no crossing. If there is a zero state of total occupation $L$, then $F(t)\ge t^L$ and consequently $\lim_{t\uparrow1}F(t)\ge1$. Since $F(0)\le1/16$, continuity gives a crossing of $1/2$; strict monotonicity makes it unique.

It remains to prove computability without assuming that equality of computable reals is decidable. Write $F(t)=\sum_{s\ge0}b_s t^s$, with $b_s\ge0$. Termwise differentiation inside the unit disk gives, for $0<t<1$,
\begin{equation}\label{ste:ta:eq:slope}
tF'(t)=\sum_{s\ge1}s b_s t^s\ge F(t)-c_0.
\end{equation}
In particular, on any interval where $F(t)\ge3/8$, one has $F'(t)\ge1/4$.

Given a requested rational width $0<\eta\le1$, maintain an interval $[a,b]\subseteq[0,1]$ containing $t_*$, initially $[0,1]$. At its rational midpoint $u<1$, compute a rational enclosure $[L,U]$ of $F(u)$ with width at most $\delta=\eta/16$, using Theorem~\ref{ste:ta:thm:partition}. If $U<1/2$, replace $a$ by $u$. If $L>1/2$, replace $b$ by $u$. These updates preserve the enclosure, including the case $t_*=1$.

In the remaining case, $1/2\in[L,U]$ and hence $\abs{F(u)-1/2}\le\delta$. Such a value is impossible in a no-root instance, whose entire range lies below $1/8$. Thus a crossing exists. Throughout the interval joining $u$ to the crossing, the function is at least $1/2-\delta\ge7/16>3/8$. Integrating the derivative lower bound gives
\[
\abs{u-t_*}\le4\abs{F(u)-1/2}\le4\delta.
\]
Intersecting $[a,b]$ with $[u-4\delta,u+4\delta]$ produces a valid interval of width at most $8\delta=\eta/2$. Otherwise ordinary bisection halves the width at every iteration and terminates when it is at most $\eta$. Every function evaluation is at a rational point strictly below one and terminates. This proves uniform computability.

Finally, the fixed universal polynomial has a zero state exactly on $\mathsf K$. The definition of $t_*$ then gives \eqref{ste:ta:eq:endpoint}, and with it the completeness assertions.
\end{proof}

\begin{remark}[The exact boundary is not a noncomputable real]
It would be incorrect to describe $t_*(c)$ as a noncomputable number. There is one uniform algorithm that encloses it to any requested precision, even for nonhalting inputs. What is undecidable is the exact assertion that its value equals one. A finite-precision interval very close to the boundary does not distinguish the boundary itself from a sufficiently nearby interior crossing.
\end{remark}

\subsection{No computable uniform distance from the boundary}
\begin{corollary}[No instance-computable positive separation]\label{ste:ta:cor:separation}
There is no total computable function $d:\N\to\mathbb Q_{>0}$ such that every halting input in the universal family satisfies
\[
1-t_*(c)\ge d(c).
\]
There is likewise no total computable choice of rational $0\le t(c)<1$ guaranteed to satisfy $F_c(t(c))\ge3/4$ on every halting input.
\end{corollary}
\begin{proof}
For the first assertion, replace $d(c)$ by $\min(d(c),1)$ if necessary and compute $t_*(c)$ to error less than $d(c)/4$. A nonhalting input has value one; a halting input is at most $1-d(c)$. These disjoint alternatives would decide $\mathsf K$.

For the second assertion, evaluate $F_c(t(c))$ to error less than $1/8$. On nonhalting inputs its value is below $1/8$, while the promised halting value is at least $3/4$. Again the alternatives are effectively separated and would decide halting.
\end{proof}

\subsection{Witness heights determine the crossover scale}
\begin{theorem}[Finite-multiplicity height bounds]\label{ste:ta:thm:height}
Suppose $1\le M<\infty$, and let
\[
L=\min\{s(n):n\in\mathcal Z\}\ge1.
\]
Then the critical fugacity satisfies
\begin{equation}\label{ste:ta:eq:heightbounds}
\frac{\log2}{L}\le-\log t_*\le
\frac{\log\bigl(M/(1/2-\epsm)\bigr)}{L}.
\end{equation}
For $M=1$, more precisely,
\begin{equation}\label{ste:ta:eq:uniqueheight}
\frac12-\epsm\le t_*^L\le\frac12.
\end{equation}
If the zero set is infinite but all its heights are at least $L$, one still has the unconditional estimate
\begin{equation}\label{ste:ta:eq:infiniteheight}
G(t)\le\frac{t^{L/2}}{(1-\sqrt t)^m},\qquad 0<t<1.
\end{equation}
\end{theorem}
\begin{proof}
At the crossing, $G(t_*)=1/2-R(q_m,t_*)$, so
\[
1/2-\epsm\le G(t_*)\le1/2.
\]
A minimum-height zero contributes $t_*^L$, while every zero contributes at most $t_*^L$. Hence
\[
t_*^L\le G(t_*)\le M t_*^L.
\]
Combining these inequalities and taking negative logarithms proves \eqref{ste:ta:eq:heightbounds}; with one zero, $G(t)=t^L$ and \eqref{ste:ta:eq:uniqueheight} follows. For the final estimate, use $t^{s(n)}\le t^{L/2}(\sqrt t)^{s(n)}$ whenever $s(n)\ge L$, and sum over all natural tuples.
\end{proof}

The theorem identifies a separate height obstruction even in a single-witness example. A witness can be unique and still require a fugacity arbitrarily close to one before its contribution becomes visible. Moreover, the relevant height is that of the compiled tuple. Large multiplication-gate values can make it much larger than the source witness height.

\section{What remains undecidable despite the positive-energy gap}\label{ste:ta:sec:gap}
\subsection{Bottom-of-spectrum separation}
\begin{theorem}[A gap does not provide a search cutoff]\label{ste:ta:thm:gap}
For the fixed universal family, the bottom spectral value satisfies
\[
\inf\Spec H_c=0\quad\text{or}\quad\inf\Spec H_c\ge1,
\]
and distinguishing the alternatives is undecidable. There is no total computable bound $B(c)$ such that a zero state, whenever one exists, always has a representative in $\{0,\ldots,B(c)\}^m$. In particular, no total finite-truncation procedure can uniformly approximate the bottom spectral value to error less than $1/4$.
\end{theorem}
\begin{proof}
All energies are nonnegative integers, and energy zero occurs exactly on halting inputs. A computable occupation cutoff with the stated property would permit exact finite search for a zero and thereby decide halting. Any uniform approximation to the bottom spectral value with error less than $1/4$, whether obtained by truncation or otherwise, would distinguish zero from values at least one by comparison with $1/2$.
\end{proof}

This does not conflict with Theorem~\ref{ste:ta:thm:layers}. To find all \emph{positive}-energy states below $E$, one needs only a known finite simplex. A zero state can lie outside every proposed computable instance-dependent cutoff, because multiplication by the radial factor leaves every true zero of $Q$ at energy zero. The two searches have different logical content.

\wtag\ The second sentence of Theorem~\ref{ste:ta:thm:gap} is the proposition \texttt{cdc:bd:prop:nobound} of~\cite{cdc} applied to the compiled penalty, whose zero set is that of $H_c$; the spectral form ``$\inf\Spec H_c\in\{0\}\cup[1,\infty)$, undecidable which'' is source~03's restatement.

\subsection{Bounded operator perturbations versus coefficient errors}
\begin{proposition}[A precise perturbative stability statement]\label{ste:ta:prop:perturb}
Let $V$ be any bounded self-adjoint operator with $\norm V\le\delta<1/4$. Then
\[
\abs{\inf\Spec(H+V)-\inf\Spec H}\le\delta.
\]
Thus zero-root and root-free instances retain disjoint bottom-energy ranges, respectively $[-\delta,\delta]$ and $[1-\delta,\infty)$. Moreover, $H+V$ has compact resolvent exactly when $H$ does.
\end{proposition}
\begin{proof}
For every unit vector in the operator domain, the quadratic form of $V$ lies between $-\delta$ and $\delta$. Taking infima proves the spectral-bottom inequality. Bounded self-adjoint perturbation preserves self-adjointness on $\mathcal D(H)$; it can also be seen by a resolvent Neumann series at a sufficiently large imaginary argument. At any common nonreal resolvent point, the resolvent identity expresses the difference of the two resolvents as a product containing either resolvent. If one resolvent is compact, the identity makes the other compact; interchanging the operators proves the converse.
\end{proof}

The proposition is deliberately not a statement of coefficient-noise robustness. For every $\delta>0$, the coefficient perturbation $\delta N_i$ has infinite operator norm on the full occupation space. A very small perturbation of polynomial coefficients can become arbitrarily large at high occupation. Exact coefficient specification, unbounded occupation, and exact boundary predicates are substantive features of this representation, not implementation details that a spectral-gap statement removes.

Nor is diagonal evolution a Turing-complete time evolution in the usual operational sense. For an occupation basis vector,
\[
e^{-itH}\ket n=e^{-itH(n)}\ket n:
\]
it changes a phase and does not execute a sequence of machine configurations. The universality proved here belongs to static spectral predicates and to the partition-function boundary, not to autonomous simulation or to a physical algorithm for an undecidable problem.

\section{Worked examples and exact computational checks}
\subsection{A single witness that moves toward the boundary}
For $N\in\N$, consider the direct two-mode system
\begin{equation}\label{ste:ta:eq:HN}
H_N(x,a)=(1+x+a)\bigl((x-N)^2+(a-1)^2\bigr).
\end{equation}
It has exactly one zero, $(N,1)$, of total height $L=N+1$. Here $m=2$, $q_m=1/24$, and
\[
G_N(t)=t^{N+1}.
\]
The Hamiltonian in this special example has degree three; the degree-five normal form is a general upper bound, not a requirement that every instance have exact degree five.

Table~\ref{ste:ta:tab:critical} reports certified crossings obtained by combining the exact ground polynomial with a rational enclosure of the excited trace. Each enclosure has width at most $2^{-30}$; the displayed decimals have absolute error below $2\times10^{-9}$. The rational endpoints, not just the rounded decimals, are saved in \texttt{verification.json}.
\begin{table}[ht]
\centering
\begin{tabular}{rrr}
\toprule
$N$ & Zero-state height $L$ & Critical fugacity $t_*$\\
\midrule
0 & 1 & 0.458301564\\
1 & 2 & 0.704142627\\
3 & 4 & 0.840894904\\
10 & 11 & 0.938930911\\
100 & 101 & 0.993160652\\
1000 & 1001 & 0.999307785\\
\bottomrule
\end{tabular}
\caption{Single-witness crossover values for \eqref{ste:ta:eq:HN}, at $q=1/24$.}\label{ste:ta:tab:critical}
\end{table}
For example, the enclosure for $N=1000$ is exactly
\[
\frac{1072998563}{1073741824}
\le t_*\le
\frac{268249641}{268435456}.
\]
Every member of this elementary family is decidable directly. Its role is to display the height effect cleanly, not to demonstrate an undecidable individual instance.

\subsection{Infinite zero multiplicity with a computable excited trace}
Take
\begin{equation}\label{ste:ta:eq:infiniteexample}
H(x,a)=(1+x+a)(a-1)^2.
\end{equation}
All $(x,1)$ are zero states, so
\[
G(t)=\frac{t}{1-t}.
\]
For $a\ne1$, summing the geometric series in $x$ gives the independent formula
\begin{equation}\label{ste:ta:eq:explicitR}
R(q,t)=\sum_{\substack{a\ge0\\a\ne1}}
\frac{t^a q^{(1+a)(a-1)^2}}{1-tq^{(a-1)^2}}.
\end{equation}
At $t=1$ this excited series is finite by Theorem~\ref{ste:ta:thm:partition}, even though the full thermal trace is infinite. The critical fugacity at $q=1/24$ is enclosed by
\[
\frac{337179061}{1073741824}\le t_*\le
\frac{168589531}{536870912},
\]
with midpoint approximately $0.314022472$. It is below $1/3$, the crossing of the ground series alone, because the excited contribution is strictly positive.

For an integer $E>0$, the multiplicity of energy $E$ equals the number of natural $a\ne1$ satisfying
\[
(a-1)^2\mid E,\qquad
\frac{E}{(a-1)^2}-a-1\ge0.
\]
The corresponding $x$ is uniquely the expression on the left of the inequality. This divisor formula was checked independently against finite occupation enumeration for every $1\le E\le60$.

\subsection{An arithmetic infinite fiber: Pell witnesses}
A less degenerate infinite zero set comes from
\begin{equation}\label{ste:ta:eq:pell}
H(x,y,a)=(1+x+y+a)
\bigl((x^2-2y^2-1)^2+(a-1)^2\bigr).
\end{equation}
Starting with $(x_0,y_0)=(1,0)$, define
\[
x_{j+1}=3x_j+4y_j,\qquad y_{j+1}=2x_j+3y_j.
\]
Direct expansion gives
\[
(3x+4y)^2-2(2x+3y)^2=x^2-2y^2.
\]
The positive recurrence strictly increases both coordinates after the initial step and therefore produces infinitely many distinct natural solutions. Thus \eqref{ste:ta:eq:pell} has infinitely many zero-energy states and no compact resolvent, while its excited heat operator has the explicit effective tail bound of Theorem~\ref{ste:ta:thm:remainders}. The proof needs only this infinite subfamily, not a classification of every Pell solution. The first 32 recurrence steps were checked with exact integers.

\subsection{Why the radial factor must contain every mode}
The root-free polynomial $Q(x,y)=x^2+1$ has an infinitely degenerate eigenvalue one at $(0,y)$ when used as an unweighted energy. Multiplication by $1+x$ does not help: the omitted $y$ coordinate is still a flat positive-energy direction. Multiplication by $1+x+y$ does help, because every positive state then has energy at least $1+x+y$.

The same issue applies to the padding coordinate in Theorem~\ref{ste:ta:thm:padding} and to intermediate circuit variables. Omitting a mode from the radial factor invalidates the uniform positive-energy bounds unless an additional independent argument controls that mode.

\subsection{Compiler and verification package}\label{ste:ta:sec:package}
The accompanying \texttt{compiler.py} implements natural addition/multiplication circuit construction, commutative gate sharing, exact residual evaluation, canonical witness lifting, rational thermal enclosures, and critical-fugacity enclosure. Sparse polynomial inputs use exponent tuples with parameters first and source witnesses second. Floats and Boolean values are rejected where exact integer inputs are required.

\wtag\ \emph{Editorial note: this holds for \texttt{compile\_polynomial}; the low-level \texttt{Compiler} and \texttt{Atom} interface assumes well-formed natural DAG descriptors and does not check them.} The research programme's review~\cite{review1977} found that the low-level builder accepts a self-referential gate and non-natural constants: \texttt{c = Compiler(1); z = c.add(c.variable(0), Atom('mode', 1)); f = c.finish(z, z)} defines the gate $u=x+u$ and gives \texttt{f.energy((0,7,1)) == f.energy((0,8,1)) == 0}, two completions over one input; \texttt{Atom('const', 0.5)} is accepted and yields energy $0.5$; Boolean atoms are accepted. Both behaviours were reproduced on a copy of the shipped compiler at the write. No theorem is affected: the theorems assume a natural DAG whose gate inputs are constants in $\N$, parameters, witnesses or outputs of \emph{earlier} gates (Section~\ref{ste:ta:sec:intro} and the conventions before Theorem~\ref{ste:ta:thm:compiler}). The review's tested patch \path{thermal_exact_dag_guards.patch} validates atoms and gate descriptors; with it the 7,936-assertion suite still passes and both JSON outputs are unchanged. The shipped \path{code/03-thermal-arithmetic-compiler.py} is the unpatched original. In this report the package files are shipped as \path{code/03-thermal-arithmetic-compiler.py}, \path{code/03-thermal-arithmetic-verify.py}, \path{code/03-thermal-arithmetic-build.sh}, \path{data/03-thermal-arithmetic-examples.json}, \path{data/03-thermal-arithmetic-verification.json}, \path{data/03-thermal-arithmetic-verification_console.txt}, \path{data/03-thermal-arithmetic-requirements.txt} and \path{03-thermal-arithmetic-PROVENANCE.md}; the example below imports the compiler by its delivered name \texttt{compiler}.

For example, the product compiler can be invoked as follows:
\begin{lstlisting}[language=Python]
from compiler import compile_polynomial

# P(p; x, y) = x*y - p; tuple order is (p, x, y).
model = compile_polynomial(
    {(0, 1, 1): 1, (1, 0, 0): -1},
    inputs=2, parameters=1,
)
n = model.canonical((2, 3), (6,))
assert n == (2, 3, 6, 1)
assert model.energy(n, (6,)) == 0
assert (model.modes, model.equations) == (4, 3)
\end{lstlisting}

The verification suite performs 7,936 assertions. These are reproducible finite and symbolic regression checks, not a substitute for the general proofs and not a proof-assistant formalization. Table~\ref{ste:ta:tab:tests} groups the actual assertions without claiming statistical independence.
\begin{table}[ht]
\centering\small
\begin{tabular}{p{0.74\textwidth}r}
\toprule
Check group & Assertions\\
\midrule
Canonical gate/marker values, signed outputs, penalty identities, source-zero equivalence & 1,240\\
Mutation of a single auxiliary coordinate is rejected & 1,120\\
Exhaustive small supplied-tuple zero-fiber equivalence & 2,384\\
Radial positive-energy lower bounds on supplied tuples & 2,384\\
Symbolic degree, exact resource counts, and mode supports & 170\\
Exact input-type rejection & 5\\
Dimension-uniform thermal bound for $1\le m\le200$ & 200\\
Finite energy shells and independent multiplicities & 192\\
Pell recurrence identity and strict growth & 64\\
Certified partition widths, independent overlap, global bounds & 60\\
Critical enclosures, height bounds, endpoint and algorithm checks & 17\\
Omitted-mode confinement counterexample & 100\\
\midrule
Total & 7,936\\
\bottomrule
\end{tabular}
\caption{Executed exact and symbolic checks. Detailed categories and rational outputs appear in \texttt{verification.json}.}\label{ste:ta:tab:tests}
\end{table}

Six source polynomials are exported in \texttt{examples.json}: $x-p$, $x^3-p$, $xy-p$, $x^2+1$, $x^2+xy-py-3$, and the zero polynomial. Their complete compiled polynomial expressions are included. Symbolic expansion independently confirms the degree and support bounds for these examples. Small supplied tuples include noncanonical gate assignments; the suite is not limited to evaluating the intended lifts.

\subsection{Certified truncation without enormous denominators}
The rational partition routine uses two explicit error budgets. First choose a box so that \eqref{ste:ta:eq:Ztail} or \eqref{ste:ta:eq:Rtail} is at most half the requested tolerance. Let $D=(B+1)^m$ be the number of tuples in that box. Next choose an integer energy cutoff $C$ such that
\[
Dq^C\le\frac{\text{tolerance}}{2}.
\]
Terms with $H(n)\ge C$ may be omitted from the finite rational sum, since their total contribution is at most $Dq^C$ when $0\le t\le1$. This avoids constructing denominators such as $24^{H(n)}$ for unnecessarily large energies. The output interval is
\[
[\text{retained exact sum},\,
 \text{retained exact sum}+\text{box tail}+Dq^C].
\]
No omitted term is silently discarded. A configurable resource limit raises an exception rather than claiming a certificate from unfinished work. The optional exact ground-series interface speeds up the examples near $t=1$; supplying such a formula is an extra, independently justified input, not an assumed ability to solve the general zero-set problem.

\section{Further research questions}\label{ste:ta:sec:questions}
The following questions are selected to extend the proved results rather than restate the general halting problem.

\subsection{Can radial isolation be achieved at degree four?}\label{ste:ta:q:degree4}
The direct multiplier converts a quartic sum of gate squares into a quintic. Is there a witness-faithful finite-variable polynomial normal form of degree at most four with both the same zero fibers and an effective bound excluding positive low-energy states at large occupation? The target need not be the exact inequality $H\ge1+s$; any computable proper lower bound on the positive sector would preserve much of the analysis. A construction must account for the possibility that fresh auxiliaries create their own low-energy flat directions. A lower bound for a clearly specified class of constructions would also be useful.
(\wtag\ Degree at most three is impossible for a universal family: $H$ is nonnegative on the real nonnegative orthant, and an orthant-nonnegative polynomial of degree at most three has a semilinear natural zero set (\texttt{cdc:of:thm:classification} and \texttt{cdc:of:cor:no-cubic-universal} in~\cite{cdc}; for globally nonnegative coercive polynomials compare \texttt{pqc:ql:prop:degree} in~\cite{pqc}), whereas the zero sets of Corollary~\ref{ste:ta:cor:universal} project onto a c.e.-complete set. For universal families the open gap is therefore four versus five.)

\subsection{Can four-mode positive terms be reduced to three-mode terms?}\label{ste:ta:q:threemode}
The extra mode enters when $n_i$ multiplies a residual square already involving three modes. A useful refinement would lower the maximum support without losing positivity, introducing uncontrolled witness multiplicity, or increasing degree beyond a stated budget. The appropriate initial task is to classify local identities for one multiplication residual and one external radial coordinate. Geometric locality is a separate, harder requirement.

\subsection{Instantiate a numerical universal arithmetic family}\label{ste:ta:q:instantiate}
Take one audited universal polynomial from ProveIt and apply the compiler literally. Record source witness count, shared gate count, coefficient bit lengths, expanded degree, positive-term supports, and the cost of exact energy evaluation. This would turn Corollary~\ref{ste:ta:cor:universal} from a fixed-family existence result into a reproducible numerical normal form. It should not be reported as an improved universal-operation bound unless that operation count is actually lower under the same convention.

\subsection{Optimize compiled witness height, not only gate count}\label{ste:ta:q:height}
Two zero-set-bijective arithmetic circuits can have very different sums of intermediate values. Theorem~\ref{ste:ta:thm:height} makes this a measurable analytic cost. Define and compare the height overhead of alternative circuits, including product trees, shared powers, redundant checks, and parameter loaders. A circuit with fewer gates might create a much larger critical-fugacity boundary layer. A useful objective is a joint frontier in gate count, maximum intermediate value, and total zero-state occupation.

\subsection{Find unconditional all-fiber compact subclasses}\label{ste:ta:q:compact}
Theorem~\ref{ste:ta:thm:finitefold} identifies the exact representation requirement for all-fiber compactness. Rather than assume it for universal halting, identify interesting restricted computational relations with explicit finite-fold representations and carry out the compact normal form. Preserving known fiber cardinalities under composition is especially important. Any genuinely universal solution would need to confront the finite-fold issue, not merely make the gate completion unique.

\subsection{Characterize effective compactness on restricted families}\label{ste:ta:q:effective}
Even where compact resolvent is mathematically guaranteed, a uniform computable modulus for its finite-rank approximation may fail. Which syntactic or arithmetic hypotheses provide such a modulus? An explicit computable witness-height bound is sufficient. Determine whether weaker information, such as an effective bound on zero-state count combined with a controlled arithmetic geometry, also suffices for useful subclasses. The answer should distinguish compactness, computable compactness, and computability of individual matrix entries.

\subsection{Classify witness-series singularities}\label{ste:ta:q:singularities}
For infinite fibers, $G(t)$ controls all singular behavior at $t=1$, while $R$ extends to a larger disk. For families with structured witnesses, determine asymptotic expansions of $G(t)$ as $t\uparrow1$: polynomial poles for free coordinates, logarithmic growth for exponentially sparse families, and possible logarithmic oscillations. The Pell example is a concrete starting point. Any all-orders expansion would directly become an expansion of the full partition function after adding the analytic excited contribution. This offers a connection with the repository's generating-function and transseries work without importing an unproved asymptotic claim into this article.

\subsection{What survives non-diagonal bounded couplings?}\label{ste:ta:q:couplings}
Proposition~\ref{ste:ta:prop:perturb} preserves compactness and a separated bottom-energy test under bounded self-adjoint perturbations. It does not preserve the exact ground projection or the simple partition decomposition. Seek quantitative heat-trace comparison bounds and effective truncations for a precisely specified class of bounded couplings. An important first distinction is whether the perturbation is effectively approximable in norm or merely has computable matrix entries.

\subsection{Can a physically motivated cutoff retain a useful conditional theorem?}\label{ste:ta:q:cutoff}
A fixed maximum occupation makes the space finite dimensional and every zero test decidable. The meaningful question is therefore not whether this reproduces unrestricted halting, but which bounded computational questions it represents with certified precision and resource costs. Compare known runtime or witness-height bounds with occupation limits, coefficient errors, and energy resolution. A valid theorem should state explicitly which promise supplies the missing cutoff.

\subsection{Formalize the arithmetic and effective-analysis interface}\label{ste:ta:q:formalize}
A modular formalization can begin with the exact zero-fiber bijection, the radial inequality, and finite-box geometric tails. These are independent of a full formalization of unbounded operators. The next layer can treat diagonal operators directly on $\ell^2$, prove the maximal-domain realization, and connect trace and compactness to diagonal sequences. The universal result should depend on a separately audited MRDP interface, and the finite-fold equivalence should remain visibly conditional. The present executable code is a test oracle for this project, not a formal proof artifact.

\section{Dependency and claim ledger}
\begin{longtable}{>{\raggedright\arraybackslash}p{0.30\textwidth}>{\raggedright\arraybackslash}p{0.63\textwidth}}
\toprule
Statement & Status and dependencies\\
\midrule
\endhead
Universal Diophantine representation & Established foundation, used from the MRDP/universal-equation literature; no new proof or numerical instance supplied.\\
Squared number-operator embedding & Established prior construction; explicitly attributed to Kieu. No physical decision claim used.\\
Witness-faithful compiler & Proved here by topological induction; classical circuit method with an explicit marker and quantitative ledger. \wtag\ The circuit core is \texttt{cdc:wf:lem:quadratization} of~\cite{cdc} (note in Section~\ref{ste:ta:sec:intro}).\\
Degree-five radial isolation and four-mode positive display & Proved here by integer positivity and support counting.\\
Positive spectral cutoffs and effective remainders & Proved here from the radial inequality, geometric sums, and diagonal operator arguments.\\
Compactness and thermal multiplicity classification & Proved here; exactly preserves the source root count.\\
All-fiber compact universal representation & Not asserted unconditionally. The exact finite-fold equivalence is proved.\\
Padded compactness and analytic-radius undecidability & Unconditional consequence of fixed universal representability and free-witness padding.\\
Uniformly computable critical fugacity & Proved here using certified interior evaluations and a positive derivative bound.\\
Undecidable exact endpoint and absent cutoffs & Proved by reductions to halting. The critical real itself is computable.\\
Numerical and symbolic checks & 7,936 executed assertions; evidence about the implementation and selected examples, not proof-assistant verification.\\
Priority or breakthrough status & Not established by this targeted search. The article supplies explicit proved results and a research program, not a certified historical first.\\
\bottomrule
\end{longtable}

\section{Conclusion}
A small structural modification of an arithmetic penalty yields a useful separation: difficult zero states are left exactly where the Diophantine witnesses put them, while every positive state is forced into an effective radial energy bound. This separates unknown witness multiplicity from an otherwise computable excited spectrum and gives exact correspondences between finite-fold arithmetic, compact resolvent, and finite thermal trace.

The boundary construction adds a second distinction. A family of partition functions can be uniformly computable throughout its regulated domain, and its threshold location can itself be a uniformly computable real, while the exact endpoint predicate remains complete for nonhalting. No inaccessible approximation oracle is hidden in the proof; the obstruction is the impossibility of a uniform positive separation from the boundary.

The resulting substrate is unconventional but mathematically precise: a fixed finite collection of unbounded occupation modes with polynomial diagonal energies. It is not a physical solution of an undecidable problem. Its value is to expose which parts of a computational representation are existential, which are counting-theoretic, which are controlled by witness height, and which become visible only at an exact analytic boundary.


\appendix
\section{A fully instantiated horizon-two polynomial schema}
\label{ste:ee:app:tinyquartic}

\noindent\wtag\ Appendix~A of source~02.


For the three-state alphabet \eqref{ste:ee:eq:tiny}, let $e_{t,0},e_{t,R},e_{t,L}$ be the
natural selectors at time $t=1,2$. Its compressed information matrix is
\[
 A(e_t)=e_{t,0}H_0+e_{t,R}H_R+e_{t,L}H_L
 =3\begin{pmatrix}
 e_{t,R}+e_{t,L}&e_{t,R}\\
 e_{t,L}&e_{t,R}+e_{t,L}
 \end{pmatrix}.
\]
Let $U=\one_2\one_2\transpose$, $Y_0=I_2+U$, $p_0=1$, $p_1=39$, and
$p_2=1521$. The following is one explicit integer polynomial in fourteen natural
variables (six selectors and eight entries of $Y_1,Y_2$):
\begin{align*}
 \mathcal I_2={}&
 \sum_{t=1}^2(e_{t,0}+e_{t,R}+e_{t,L}-1)^2\\
 &+\sum_{t=1}^2\sum_{a,b=1}^2
 \left((Y_t)_{ab}-p_t-
       \bigl[A(e_t)(Y_{t-1}-p_{t-1}U)\bigr]_{ab}\right)^2\\
 &+\sum_{a,b=1}^2((Y_2)_{ab}-1521)^2.
\end{align*}
All finite sums, coefficients, dimensions, and constant matrices are specified;
expanding the squares is optional and does not change the polynomial. The word
$(0,L)$ gives selectors $(1,0,0)$ and $(0,0,1)$, with $Y_1=39U$ and $Y_2=1521U$.
Every displayed residual is then visibly zero. A word containing only $R,L$
produces an invertible transverse product and cannot satisfy the terminal
conditions.

This example illustrates the natural-domain uniqueness directly. A selector
assignment such as $(2,-1,0)$ is inadmissible, even though its entries sum to one;
using arbitrary integer selectors without additional conditions would invalidate
the one-hot proof. The checker enforces the natural domain exactly.

\section{Reproduction and source ledger}
\label{ste:ee:app:reproduction}

\noindent\wtag\ Appendix~B of source~02.


The archive contains the standalone \path{article.tex}, its compiled
\path{article.pdf}, the Python implementation and independent evaluator,
numerical example JSON files, a test receipt, and separate source/proof-status
notes. From the package directory, run
\begin{verbatim}
python3 code/run_checks.py
python3 code/check_certificate.py \
    examples/seven_state_information_certificate.json
pdflatex -interaction=nonstopmode -halt-on-error article.tex
pdflatex -interaction=nonstopmode -halt-on-error article.tex
pdflatex -interaction=nonstopmode -halt-on-error article.tex
\end{verbatim}
The first command rewrites the delivered examples and verification receipt. The
checker returns exit status zero for an accepted certificate and nonzero for a
rejected or ill-typed certificate. The deliberately incorrect uniform-target
example is expected to fail, not to be treated as a regression.

\wtag\ These commands use the delivered layout and names; the first one overwrites the recorded examples and receipt, and on Windows writes them with CRLF line endings. In this report the programs are \path{code/02-exact-erasure-*.py}, the examples and the three verification records are \path{data/02-exact-erasure-*.json}, and the delivered \path{article.tex} is printed as Part~\ref{ste:part:ee}; the duplicate \path{verification/run_output.txt} (byte-identical to the check results) and the source's PDF are not shipped. The README of this report gives a recipe that runs on a copy and leaves the shipped files untouched.

The numerical implementation uses only the Python standard library. The PDF
build uses pdfLaTeX and standard TeX packages named in the preamble. The version
used for the recorded Python checks was 3.13.5. The final PDF is rendered and
checked separately; the PDF preflight record is included in the verification
folder.

The external dependencies are narrow. The mortality bounds come from
\cite{cassaigne}; the halting/PCP interface is supported by \cite{forster}; the
computably-enumerable/Diophantine correspondence is DPRM \cite{matiyasevich}.
The MathOverflow discussion \cite{miller} supplies the precise research question,
not a theorem used as an unproved premise. The repository sources
\cite{repoH10,repoReview,repoProb} supply project context and claim boundaries,
not automatic mathematical authority for the present results.

\section{Source audit and provenance}\label{ste:ec:app:audit}

\noindent\wtag\ Appendix~A of source~10.


Repository inspection used the connected GitHub interface on October 2,
2026. The relevant reads were:
\begin{enumerate}
\item \nolinkurl{Computability/HilbertTenthProblem/README.md}, first 210 lines,
      returned blob \texttt{880d2a13393ea55110b49e573180b51b011d9bc9}.
      (\wtag\ The link was the unpinned \path{/blob/main/}. Blob \texttt{880d2a13} is the README introduced by commit \texttt{db3b377f0} and replaced by \texttt{b7a9404d4}, both after the pin \texttt{297eb58e4} used for items~2 and~3, whose README blob is \texttt{138f234b}; both versions state the 75- and 87-operation figures quoted in Section~\ref{ste:ec:sec:mrdp}.)
\item \nolinkurl{Computability/HilbertTenthProblem/Lean/Diophantine/MRDP.lean},
      pinned at commit \texttt{297eb58e43157b7723cc28982083fecf8d2f29bf}.
\item The probabilistic/quantum/continuous substrate collection's README,
      pinned at the same commit; it identifies the strict-contraction and
      neural-precision reports as existing research contributions.
\item Targeted repository searches for Diophantine work, contraction,
      self-assembly, and Shapley operators. These are a scope check,
      not an exhaustive audit of every file in the repository.
\end{enumerate}

The probabilistic collection is itself an unrefereed research collection.
This article uses it for provenance and nonduplication, not as a substitute
for proving the new simulation. The code in this archive was written for
this manuscript and does not copy the repository's imported compilers.
No repository files were modified.

External sources were checked through primary publisher or author-preprint
pages. In particular, the Bellman comparison uses arXiv revision 3 dated
January 26, 2026, rather than only the February 2025 abstract. The Skolem
status comparison includes the July 2026 preprint. No claim of exhaustive
priority clearance follows from this targeted audit.

\section{A compact implementation contract}\label{ste:ec:app:contract}

\noindent\wtag\ Appendix~B of source~10.


\texttt{Program} validates a finite source table and natural initial
counters. \texttt{counter\_circuit} emits the absorbing or erasing map.
\texttt{layer\_circuit} inserts shared delay copies and pads all outputs.
\texttt{compile\_program} emits paired-channel actions, normalization
constants, and optional common reset and reward. \texttt{initialize}
emits the complete dyadic payoff vector.
(\wtag\ The review~\cite{review808} found that an explicitly supplied \emph{integer} scale in \texttt{Program.encode} is divided in floating point, so that a counter value $1075$ becomes $0.0$ instead of $2^{-1075}$; the default rational scale, used by every shipped loader and export, is exact. Its patch \path{exact_convergence_integer_scale.patch} normalizes integer or rational scales and rejects floating-point and Boolean ones. The shipped \path{code/10-exact-convergence-bellman_diophantine.py} is the unpatched original.)

\texttt{certificate} accepts the game, initial numerator vector and positive
denominator, a natural horizon, and either a scalar pair, consensus test,
or rational constant full-vector target. It emits every squared linear
form and complementary product, including all reward forcing coefficients.
Its metadata reports the exact witness count, scaling convention, and
endpoint-square count.

A JSON linear form is a list of integer pairs $(j,a_j)$; index $-1$ denotes
the constant monomial, and all other indices refer to the supplied variable
label array. A complementarity term is an index pair $(u,v)$ and means
$W_uW_v$, not a residual that still needs squaring. The complete polynomial
is the sum of the squares of all listed linear forms plus all listed
complementarity terms. This specifies an ordinary integer polynomial
without auxiliary syntax such as ``run,'' ``minimum,'' or ``there exists
a schedule.''

The independent checker evaluates that polynomial and, optionally,
reconstructs the game trajectory from the serialized stochastic rows.
Unsupported domains and malformed instructions are rejected. The implementation
is intended as a transparent research reference, not as a hardened service
for untrusted, resource-exhausting inputs.

\section{Reproduction and file inventory}\label{ste:ta:app:reproduction}

\noindent\wtag\ Appendix~A of source~03.

The archive contains the following principal files:
\begin{center}\small
\begin{tabular}{ll}
\toprule
File & Purpose\\
\midrule
\texttt{article.tex}, \texttt{article.pdf} & Source and compiled research article\\
\texttt{compiler.py} & Exact natural circuit and thermal routines\\
\texttt{verify.py} & Finite and symbolic verification suite\\
\texttt{verification.json} & Exact results and rational threshold enclosures\\
\texttt{examples.json} & Gate-level and expanded compiled examples\\
\texttt{README.md} & Usage, limitations, and reproduction instructions\\
\texttt{PROVENANCE.md} & Repository snapshot and primary-source links\\
\texttt{requirements.txt} & Symbolic verification dependency\\
\texttt{build.sh} & PDF build commands\\
\bottomrule
\end{tabular}
\end{center}
Run the tests with Python 3.10 or later and SymPy installed:
\begin{lstlisting}[language=bash]
python -m pip install -r requirements.txt
python verify.py
\end{lstlisting}
The arithmetic compiler and interval routines themselves use only the Python standard library. The manuscript can be rebuilt with a conventional LaTeX installation containing the packages listed in its preamble:
\begin{lstlisting}[language=bash]
pdflatex -interaction=nonstopmode -halt-on-error article.tex
pdflatex -interaction=nonstopmode -halt-on-error article.tex
\end{lstlisting}
\wtag\ The delivered \path{build.sh} (shipped as \path{code/03-thermal-arithmetic-build.sh}) runs \emph{three} pdflatex passes, not the two shown. File names in the table above are the delivered ones; their shipped names are listed in Section~\ref{ste:ta:sec:package}. \path{verify.py} imports \texttt{compiler} by its delivered name and rewrites \path{examples.json} and \path{verification.json} in its own directory, so it must be run on a copy (README of this report); \path{article.tex} is printed as Part~\ref{ste:part:ta}, and the source's \path{README.md}, \path{article.pdf} and \path{SHA256SUMS.txt} are not shipped.

The verification receipt is generated by the delivered code. In particular, the optional accelerated ground-series path used for the near-boundary examples is identified in the implementation; it is not substituted for the general computability proof.


\section{Provenance of this report}\label{ste:app:provenance}

\wtag\ This appendix was written when the three manuscripts were merged.

\subsection{Sources}

The report was built from three manuscripts of batch~78 of ProveIt's incoming reports, cluster~H2. Two arrived in the archives of commit \texttt{1977e6ea6} and one in those of commit \texttt{808b53ed8}; all three were placed by commit \texttt{798b0c5d4}, which staged their programs, data and audit notes byte-identical to the delivery.

\begin{center}\small
\begin{tabular}{@{}L{0.07\textwidth}L{0.31\textwidth}L{0.12\textwidth}L{0.15\textwidth}L{0.24\textwidth}@{}}
\toprule
Source & Archive and title & Arrival & Pin & Printed as\\
\midrule
02 (base) & \path{ProveIt_Exact_Erasure_Research.zip}; \emph{Exact Erasure without a Common Invariant Line}; 29-page PDF & \texttt{1977e6ea6} & \texttt{928ea9701} & Part~I (Sections~\ref{ste:ee:sec:intro}--\ref{ste:ee:sec:questions} and conclusion), Appendices~A--B\\
\addlinespace
10 & \path{ProveIt_Exact_Convergence_Research.zip}; \emph{Exact Convergence as Computation}; 24-page PDF & \texttt{808b53ed8} & \texttt{297eb58e4} (README read unpinned; Appendix~\ref{ste:ec:app:audit}) & Part~II (Sections~\ref{ste:ec:sec:intro}--\ref{ste:ec:sec:agenda} and conclusion), Appendices~C--D\\
\addlinespace
03 & \path{Thermal_Arithmetic_Diophantine_Hamiltonians.zip}; \emph{Thermal Arithmetic}; 25-page PDF & \texttt{1977e6ea6} & \texttt{6914ccca6}; blob \texttt{8bd07142} of the review it cites & Part~III (Sections~\ref{ste:ta:sec:intro}--\ref{ste:ta:sec:questions}, ledger and conclusion), Appendix~E\\
\bottomrule
\end{tabular}
\end{center}

Contributions: source~02 the positive stochastic erasure compiler, its PCP front end, higher guards, decidable boundaries, zero--one law and natural quartics; source~10 the stochastic-game compiler, its boundaries and its quadratic certificates; source~03 the radially isolated Hamiltonian, its spectral and thermal classification and the critical fugacity. Every result, proof, example, remark, limitation, research question and appendix of the three manuscripts is printed. They share no theorem with one another, so no statement had to be printed once for two sources.

\subsection{Where the merge had to choose}

\begin{itemize}
\item \emph{One report, three Parts.} The three manuscripts' nearest neighbour is the report~\cite{pqc}, which already had nine Parts; they were placed together as a sibling report because they share the spine of Section~\ref{ste:sec:spine}, not a construction. Placing sources~02 and~10 as a tenth Part of~\cite{pqc} and source~03 as an eleventh was the alternative considered at placement.
\item \emph{Base and order.} Source~02 is the base (the longest text, the earliest arrival, the widest set of results); its \path{article.tex} was staged unprefixed and has been replaced in place by this merged text. Source~10 follows because stochastic games extend stochastic matrices by minima and maxima; source~03 shares no lemma with either and comes last.
\item \emph{Re-proofs of collection results.} The core of Theorem~\ref{ste:ta:thm:compiler}, the fixed-arity quartic corollaries of Parts~I and~II, the no-cutoff clauses of all three Parts and the template of Theorem~\ref{ste:ta:thm:finitefold} are results of~\cite{cdc} (and partly of~\cite{pqc}). They are printed in place, with the sources' own proofs, and marked by \wtag\ pointers naming the host labels; no novelty is claimed for them. Remark~\ref{ste:ta:rem:universalcompact} is an editorial consequence of two printed theorems.
\item \emph{Reviews and defects.} The research programme's reviews of all three manuscripts~\cite{review1977,review808} are cited where their findings apply: the thermal compiler's unchecked low-level descriptors (Section~\ref{ste:ta:sec:package}), the Bellman checker's unbound initial vector (Section~\ref{ste:ec:sec:reproduce}) and the Bellman loader's floating-point integer scale (Appendix~\ref{ste:ec:app:contract}). The delivered programs are shipped unpatched; the patches stay in the research programme.
\item \emph{Labels.} Every label carries the prefix \texttt{ste:}. The labels of source~02 are \texttt{ste:ee:} plus their delivered names, those of source~10 \texttt{ste:ec:} plus theirs, and those of source~03 \texttt{ste:ta:} plus theirs; no source label was dropped or renamed apart from the prefix. Added labels: 50 in all; 206 delivered labels give 256. Front section, Parts and provenance (12): \path{ste:sec:front}, \path{ste:sec:parts}, \path{ste:sec:spine}, \path{ste:tab:spine}, \path{ste:sec:notation}, \path{ste:tab:notation}, \path{ste:sec:relation}, \path{ste:sec:status}, \path{ste:part:ee}, \path{ste:part:ec}, \path{ste:part:ta}, \path{ste:app:provenance}. Part~I (9), the nine research questions: \path{ste:ee:q:overhead}, \path{ste:ee:q:irreducible}, \path{ste:ee:q:lowdim}, \path{ste:ee:q:schedules}, \path{ste:ee:q:certificates}, \path{ste:ee:q:fixedalphabet}, \path{ste:ee:q:nonerasure}, \path{ste:ee:q:robust}, \path{ste:ee:q:formal}. Part~II (13), its four unlabelled statements, two appendices and seven research directions: \path{ste:ec:cor:stopping}, \path{ste:ec:prop:onesided}, \path{ste:ec:prop:selectors}, \path{ste:ec:prop:threshold}, \path{ste:ec:q:instantiate}, \path{ste:ec:q:oneplayer}, \path{ste:ec:q:sparse}, \path{ste:ec:q:observations}, \path{ste:ec:q:economical}, \path{ste:ec:q:mrdp}, \path{ste:ec:q:formalize}, \path{ste:ec:app:audit}, \path{ste:ec:app:contract}. Part~III (16), four sections, its appendix, the ten research questions and the editorial remark: \path{ste:ta:sec:intro}, \path{ste:ta:rem:universalcompact}, \path{ste:ta:sec:gap}, \path{ste:ta:sec:package}, \path{ste:ta:sec:questions}, \path{ste:ta:q:degree4}, \path{ste:ta:q:threemode}, \path{ste:ta:q:instantiate}, \path{ste:ta:q:height}, \path{ste:ta:q:compact}, \path{ste:ta:q:effective}, \path{ste:ta:q:singularities}, \path{ste:ta:q:couplings}, \path{ste:ta:q:cutoff}, \path{ste:ta:q:formalize}, \path{ste:ta:app:reproduction}.
\item \emph{Numbering.} Source~02's Section~$n$ is Section~$n+1$ here (Sections~2--16) and its Appendices~A--B keep their letters; source~10's Section~$n$ is Section~$n+16$ (Sections~17--30) and its Appendices~A--B are Appendices~C--D; source~03's Section~$n$ is Section~$n+30$ (Sections~31--41) and its Appendix~A is Appendix~E. Theorem numbers follow the sections.
\item \emph{Bibliography.} The three bibliographies are merged. Source~03's \texttt{matiyasevich} (the 2010 finite-fold paper) is \texttt{matiyasevichFF}, because source~02 uses that key for the 1993 book. The repository items read by several sources at different pins are merged into one entry each, recording every pin: \texttt{repoH10} (02's \texttt{repoH10}, 03's \texttt{reporeadme}, 10's \texttt{repo-h10}), \texttt{repoReview} (02's \texttt{repoReview}, 03's \texttt{repoaudit}) and \texttt{repoProb} (02's \texttt{repoProb}, 10's \texttt{repo-substrates}). Entries marked \wtag\ were added.
\item \emph{Front matter and typography.} The three delivered title pages are replaced by one; each source's title, subtitle, author line, abstract and status statement open its Part verbatim. Source~02's remarks used the definition style; all remarks now use the remark style. No symbol was renamed (Section~\ref{ste:sec:notation}).
\item \emph{Files.} Delivered programs, data and audit notes are shipped byte-identical under prefixed names (the README maps them). Not shipped: the three PDFs, the delivered READMEs of all three and the manuscripts of sources~10 and~03 (printed as Parts~II and~III), the checksum ledgers of sources~02 and~03 (verified at placement), the two in-archive byte copies \path{verification/run_output.txt} (source~02) and \path{artifacts/test_run.txt} (source~10), and source~10's two regenerable certificates \path{countdown_certificate.json} and \path{fixedpoint_certificate.json}, which the README shows how to rebuild. All survive in the archives of their arrival commits.
\item \emph{Added text.} Every sentence written in the merge is marked \wtag: Section~\ref{ste:sec:front}; the openings of the three Parts; the notes in Part~I on the MathOverflow claim, the repository relationship and review, the $(3,6)$ attribution, the nonnegative-mortality boundary, the radix-three encoding, geometric clocks, the circuit quadratization, the seven-state wording, the no-bound argument and shipped file names; the notes in Part~II on the collection's contractions, the structural-subclasses question, the complementarity device, the quartic corollary, the no-cutoff and threshold arguments, the excluded certificates, the checker defect, the unpinned README read, the loader-scale defect and shipped file names; the notes in Part~III on the fixed-quartic route, the circuit quadratization, the Boolean energy landscape, the finite-fold template with Remark~\ref{ste:ta:rem:universalcompact}, the no-cutoff clause, the compiler defect, the degree floor, the build passes and shipped file names; the notes opening each appendix; and this appendix.
\end{itemize}

\begin{thebibliography}{99}
\bibitem{cassaigne}
J. Cassaigne, V. Halava, T. Harju, and F. Nicolas,
\emph{Tighter Undecidability Bounds for Matrix Mortality, Zero-in-the-Corner
Problems, and More}, arXiv:1404.0644v3, September 5, 2014.
\url{https://arxiv.org/html/1404.0644v3}.
The adopted bounds are $(d,m)=(3,6),(5,4),(9,3),(15,2)$; the paper's general
reducibility convention is oracle Turing reducibility.

\bibitem{forster}
Y. Forster, E. Heiter, and G. Smolka,
\emph{Verification of PCP-Related Computational Reductions in Coq},
arXiv:1711.07023v2, July 18, 2018.
\url{https://arxiv.org/abs/1711.07023v2}.
Related publication DOI: \url{https://doi.org/10.1007/978-3-319-94821-8_15}.

\bibitem{matiyasevich}
Y. V. Matiyasevich, \emph{Hilbert's Tenth Problem}, MIT Press, 1993.
Author-hosted preface, including the statement of the
computably-enumerable/Diophantine correspondence:
\url{https://logic.pdmi.ras.ru/~yumat/H10Pbook/preface.htm}.

\bibitem{miller}
V. Miller, \emph{Decidability of a matrix product being rank 1},
MathOverflow question 511960, June 2, 2026; subsequent discussion with
B. Steinberg, June 2--4, 2026. The common-invariant-line restriction appears in
the June 3 discussion. Accessed October 2, 2026.
\url{https://mathoverflow.net/questions/511960/decidability-of-a-matrix-product-being-rank-1}.

\bibitem{repoH10}
V. Reshetnikov and contributors, \emph{ProveIt: Hilbert's tenth problem} (titled \emph{Hilbert's tenth problem: Jones's articles on Diophantine representation} by sources~10 and~03), repository README \path{Computability/HilbertTenthProblem/README.md}, accessed October 2, 2026.
Source~02 read it at commit \repo:
\url{https://github.com/VladimirReshetnikov/ProveIt/blob/928ea97017a25ebe56d240c84f27d2275d818c75/Computability/HilbertTenthProblem/README.md}.
Source~03 read it at the snapshot of~\cite{repoReview}, commit \texttt{6914ccca6685baf53b7a35f25efc89366c76ba74}:
\url{https://github.com/VladimirReshetnikov/ProveIt/blob/6914ccca6685baf53b7a35f25efc89366c76ba74/Computability/HilbertTenthProblem/README.md}.
Source~10 read it through the unpinned link
\url{https://github.com/VladimirReshetnikov/ProveIt/blob/main/Computability/HilbertTenthProblem/README.md}; the returned blob is recorded in Appendix~\ref{ste:ec:app:audit}.
\wtag\ The three sources' keys \texttt{repoH10} (02), \texttt{reporeadme} (03) and \texttt{repo-h10} (10) are merged into this entry.

\bibitem{repoReview}
\emph{Review of the imported computational-substrate reports, 2 October 2026},
ProveIt research review,
\path{Computability/HilbertTenthProblem/Papers/research-wip/native-stream-queue/imported_substrate_review_20261002.md}, blob \texttt{8bd07142c8b4298fb636e47b40d2d84e993a40ce}.
Source~02 read it at commit \repo:
\url{https://github.com/VladimirReshetnikov/ProveIt/blob/928ea97017a25ebe56d240c84f27d2275d818c75/Computability/HilbertTenthProblem/Papers/research-wip/native-stream-queue/imported_substrate_review_20261002.md}.
Source~03 read it at snapshot commit \texttt{6914ccca6685baf53b7a35f25efc89366c76ba74}:
\url{https://github.com/VladimirReshetnikov/ProveIt/blob/6914ccca6685baf53b7a35f25efc89366c76ba74/Computability/HilbertTenthProblem/Papers/research-wip/native-stream-queue/imported_substrate_review_20261002.md}.
\wtag\ The keys \texttt{repoReview} (02) and \texttt{repoaudit} (03) are merged; both pins carry the same blob.

\bibitem{repoProb}
\emph{Diophantine Laws of Probabilistic, Quantum and Continuous Computation},
ProveIt combined report README,
\path{SetTheory/Cardinals/docs/reports/hilbert-tenth-problem/probabilistic-quantum-and-continuous-computation/README.md}, an AI-assisted, unrefereed research collection; its source inventory includes \emph{Strict Contractions with Diophantine-Complete Transients} and \emph{Precision Is Memory}.
Source~02 read it at commit \repo:
\url{https://github.com/VladimirReshetnikov/ProveIt/blob/928ea97017a25ebe56d240c84f27d2275d818c75/SetTheory/Cardinals/docs/reports/hilbert-tenth-problem/probabilistic-quantum-and-continuous-computation/README.md}.
Source~10 read it at commit \texttt{297eb58e43157b7723cc28982083fecf8d2f29bf}, October 2, 2026:
\url{https://github.com/VladimirReshetnikov/ProveIt/blob/297eb58e43157b7723cc28982083fecf8d2f29bf/SetTheory/Cardinals/docs/reports/hilbert-tenth-problem/probabilistic-quantum-and-continuous-computation/README.md}.
\wtag\ The keys \texttt{repoProb} (02) and \texttt{repo-substrates} (10) are merged. The report itself is~\cite{pqc}.

\bibitem{shapley}
L.~S. Shapley.
Stochastic games.
\emph{Proceedings of the National Academy of Sciences} 39(10) (1953),
1095--1100.
\href{https://doi.org/10.1073/pnas.39.10.1095}{doi:10.1073/pnas.39.10.1095}.

\bibitem{dudenhefner}
A. Dudenhefner.
Certified decision procedures for two-counter machines.
\emph{FSCD 2022}, LIPIcs 228, 16:1--16:18.
\href{https://doi.org/10.4230/LIPIcs.FSCD.2022.16}{doi:10.4230/LIPIcs.FSCD.2022.16}.

\bibitem{varonka}
A. Varonka and K. Watanabe.
On piecewise affine reachability with Bellman operators.
\href{https://arxiv.org/abs/2502.19923v3}{arXiv:2502.19923v3},
January 26, 2026. Extended version of the MFCS 2025 paper,
LIPIcs 345, 92:1--92:18.

\bibitem{vahanwala}
M. Vahanwala.
Skolem and positivity completeness of ergodic Markov chains.
\emph{Information Processing Letters} 186 (2024), 106481.
\href{https://doi.org/10.1016/j.ipl.2024.106481}{doi:10.1016/j.ipl.2024.106481}.
Author preprint: \href{https://arxiv.org/abs/2305.04881}{arXiv:2305.04881}.

\bibitem{skolem2026}
F. Luca, J. Ouaknine, and J. Worrell.
Conjectural decidability of the Skolem problem.
\href{https://arxiv.org/abs/2607.15510}{arXiv:2607.15510}, July 2026.

\bibitem{grobelna}
M. Grobelna, J. K\v{r}et\'{\i}nsk\'{y}, and M. Weininger.
Stopping criteria for value iteration on concurrent stochastic reachability
and safety games.
\href{https://arxiv.org/abs/2505.21087v2}{arXiv:2505.21087v2},
September 10, 2025. Full version of the LICS 2025 paper.

\bibitem{h10coq}
D. Larchey-Wendling and Y. Forster.
Hilbert's tenth problem in Coq.
\emph{FSCD 2019}, LIPIcs 131, 27:1--27:20.
\href{https://doi.org/10.4230/LIPIcs.FSCD.2019.27}{doi:10.4230/LIPIcs.FSCD.2019.27}.

\bibitem{repo-mrdp}
V. Reshetnikov, ProveIt repository.
\emph{Diophantine/MRDP.lean}.
\href{https://github.com/VladimirReshetnikov/ProveIt/blob/297eb58e43157b7723cc28982083fecf8d2f29bf/Computability/HilbertTenthProblem/Lean/Diophantine/MRDP.lean}{Pinned repository source},
commit \texttt{297eb58e43157b7723cc28982083fecf8d2f29bf},
read October 2, 2026.

\bibitem{jones}
James P. Jones,
\emph{Universal Diophantine equation},
Journal of Symbolic Logic \textbf{47} (1982), no.~3, 549--571.
DOI: \href{https://doi.org/10.2307/2273588}{10.2307/2273588}.
Primary publisher page consulted October 2, 2026.

\bibitem{kieu}
Tien D. Kieu,
\emph{A reformulation of Hilbert's tenth problem through quantum mechanics},
arXiv:quant-ph/0111063, version~2 (2003; initial submission 2001).
\url{https://arxiv.org/abs/quant-ph/0111063}.
The squared number-operator construction is the only external operator representation used here; proposed physical consequences are not assumed.

\bibitem{matiyasevichFF}
Yuri Matiyasevich,
\emph{Towards finite-fold Diophantine representations},
Journal of Mathematical Sciences \textbf{171} (2010), no.~6, 745--752.
DOI: \href{https://doi.org/10.1007/s10958-010-0179-4}{10.1007/s10958-010-0179-4}.
\wtag\ Source~03's key \texttt{matiyasevich}, renamed because source~02 uses that key for a different item.

\bibitem{cantone}
D. Cantone, A. Casagrande, F. Fabris, and E. G. Omodeo,
\emph{The quest for Diophantine finite-fold-ness},
Le Matematiche \textbf{76} (2021), no.~1.
\url{https://lematematiche.dmi.unict.it/index.php/lematematiche/article/view/2044}.

\bibitem{neary}
\wtag\ T. Neary,
\emph{Undecidability in Binary Tag Systems and the Post Correspondence Problem for Five Pairs of Words},
in \emph{32nd International Symposium on Theoretical Aspects of Computer Science (STACS 2015)},
LIPIcs \textbf{30} (2015), 649--661. A copy is kept at \path{Computability/HilbertTenthProblem/Lean/LIPIcs.STACS.2015.649.pdf}.

\bibitem{pqc}
\wtag\ ProveIt research report \emph{Diophantine Laws of Probabilistic, Quantum and Continuous Computation}, collection \path{SetTheory/Cardinals/docs/reports/hilbert-tenth-problem/probabilistic-quantum-and-continuous-computation/}; Parts~II (quantum substrates), IV (questions), V (exact probabilities), VII (exact observation in continuous-state dynamics), VIII (continuous time) and IX (exact stopping laws). Cited by label names, as at the write of this report.

\bibitem{cdc}
\wtag\ ProveIt research report \emph{Canonical Diophantine Certificates}, collection \path{SetTheory/Cardinals/docs/reports/hilbert-tenth-problem/canonical-diophantine-certificates/}; its witness-faithful foundations, Part~III (accelerators), Part~IX (the boundary of fixed-arity compression) and Part~XI (positivity-restricted degree threshold). Cited by label names, as at the write of this report.

\bibitem{review1977}
\wtag\ ProveIt research programme, \emph{Six new substrate reports: review, constructor repairs and local reductions}, \path{Computability/HilbertTenthProblem/Papers/research-wip/native-stream-queue/incoming_substrate_review_1977e6ea6.md} (commit \texttt{44b28c395}), sections \emph{Exact erasure} and \emph{Thermal arithmetic and its repair}, with its receipt \path{incoming_substrate_review_1977e6ea6.json} (keys \texttt{erasure}, \texttt{thermal}, \texttt{thermal\_repair}) and the patch \path{thermal_exact_dag_guards.patch} in the same directory.

\bibitem{review808}
\wtag\ ProveIt research programme, \emph{Seven incoming computational-substrate reports} and \emph{Exact Convergence and Total Quadratic Semantics: review of 808b53ed8}, \path{Computability/HilbertTenthProblem/Papers/research-wip/native-stream-queue/incoming_substrate_review_808b53ed8.md} and \path{incoming_convergence_review_808b.md} (commit \texttt{126028588}), with the receipt \path{incoming_substrate_review_808b53ed8.json} (keys \texttt{independent.convergence}, \texttt{independent.convergence\_repairs}, \texttt{independent.convergence\_scale}) and the patches \path{exact_convergence_input_binding.patch} and \path{exact_convergence_integer_scale.patch} in the same directory.

\bibitem{matrixpcp}
\wtag\ ProveIt research programme, research note \path{Computability/HilbertTenthProblem/Papers/research-wip/native-stream-queue/matrix_pcp_trace.md}, Section~1, as at the write of this report.
\end{thebibliography}
\end{document}
